% ======================================================================
%  La Lune et l'eau -- Disque secondaire circumterrestre
%  et passage non collisionnel d'un corps planétaire massif
%  Une théorie de chemin rétrodictive du système proto-terrestre au système Terre--Lune
%  Michel Debailleul — © 2026 — Licence CC BY-NC 4.0
%  Compilation : pdflatex (trois passes pour la table des matières et les renvois)
% ======================================================================
\documentclass[11pt,a4paper,openany]{report}

\usepackage[utf8]{inputenc}
\usepackage[T1]{fontenc}
\usepackage{lmodern}
\usepackage[french]{babel}
\usepackage{amsmath,amssymb}
\usepackage[a4paper,margin=2.4cm,headheight=14pt]{geometry}
\usepackage{microtype}
\usepackage{booktabs,array,tabularx,longtable}
\usepackage{xcolor}
\usepackage{graphicx}
\graphicspath{{./}{../figures/}}
\usepackage[most]{tcolorbox}
\usepackage{tikz}
\usetikzlibrary{arrows.meta,calc,positioning,decorations.pathmorphing,babel}
\usepackage{pgfplots}
\usepgfplotslibrary{fillbetween}
\pgfplotsset{compat=1.18}
\usepackage{titlesec}
\usepackage{fancyhdr}
\usepackage{hyperref}
\hypersetup{
  colorlinks=true,
  linkcolor=blue!45!black,
  citecolor=blue!45!black,
  urlcolor=blue!45!black,
  pdftitle={La Lune et l'eau -- Disque secondaire circumterrestre et passage non collisionnel d'un corps planétaire massif},
  pdfauthor={Michel Debailleul},
  pdfsubject={Une théorie de chemin rétrodictive du système proto-terrestre au système Terre--Lune},
  pdfkeywords={Lune, eau lunaire, disque circumterrestre, accrétion secondaire, rencontre non collisionnelle,
  moment cinétique, tri orbital, géochimie lunaire, volatils, rapport D/H, instabilité orbitale primitive, planète éliminée}
}

% ---------------------------------------------------------------------
%  Titres de parties et de chapitres
% ---------------------------------------------------------------------
\titleclass{\part}{top}
\titleformat{\part}[display]
  {\centering\normalfont\bfseries\color{blue!40!black}}
  {\Large\partname~\thepart}{10pt}{\Huge}
\titlespacing*{\part}{0pt}{30pt}{28pt}
\titleformat{\chapter}[display]
  {\normalfont\bfseries}{\large\chaptertitlename~\thechapter}{6pt}{\LARGE}
\titlespacing*{\chapter}{0pt}{-8pt}{18pt}

% ---------------------------------------------------------------------
%  Encadrés
% ---------------------------------------------------------------------
\tcbset{boite/.style={enhanced,breakable,boxrule=0.6pt,arc=2pt,
  left=6pt,right=6pt,top=4pt,bottom=4pt,fonttitle=\bfseries\small,
  before skip=8pt,after skip=10pt}}
\newtcolorbox{ideephys}{boite,colback=blue!3,colframe=blue!45!black,title=Idée physique}
\newtcolorbox{lecture}{boite,colback=teal!3,colframe=teal!55!black,title=Lecture physique}
\newtcolorbox{falsif}{boite,colback=red!3,colframe=red!55!black,title=Critère de falsification}
\newtcolorbox{controle}{boite,colback=orange!4,colframe=orange!70!black,title=Contrôle quantitatif}
\newtcolorbox{apport}{boite,colback=violet!4,colframe=violet!65!black,title=Apport original de la théorie}
\newtcolorbox{consequence}{boite,colback=green!4,colframe=green!40!black,title=Conséquence observable}
\newtcolorbox{definition}{boite,colback=brown!5,colframe=brown!65!black,title=Définition}
\newtcolorbox{rappel}{boite,colback=gray!5,colframe=gray!60!black,title=Rappel didactique}
\newtcolorbox{principe}[1]{boite,colback=gray!4,colframe=gray!55!black,title=#1}
\newtcolorbox{prediction}[1]{boite,colback=cyan!3,colframe=cyan!40!black,title=Prédiction --- #1}
\newcommand{\sila}{\par\noindent\textbf{Si la théorie est correcte.}\ }
\newcommand{\verif}{\par\smallskip\noindent\textbf{Comment la vérifier.}\ }

% ---------------------------------------------------------------------
%  Notations
% ---------------------------------------------------------------------
\newcommand{\Mt}{M_{\oplus}}
\newcommand{\Rt}{R_{\oplus}}
\newcommand{\Req}{R_{\mathrm{eq}}}
\newcommand{\aR}{a_{\mathrm{Roche}}}
\newcommand{\dd}{\mathrm{d}}
\newcommand{\Mtot}{M_{\mathrm{tot}}}
\newcommand{\RT}{\mathcal{R}_{\oplus}}
\newcommand{\RD}{\mathcal{R}_{D}}
\newcommand{\RB}{\mathcal{R}_{B}}
\newcommand{\Rinf}{\mathcal{R}_{\infty}}

% ---------------------------------------------------------------------
%  En-têtes
% ---------------------------------------------------------------------
\pagestyle{fancy}
\fancyhf{}
\fancyhead[L]{\small\nouppercase{\leftmark}}
\fancyhead[R]{\small Michel Debailleul}
\fancyfoot[C]{\small\thepage}
\renewcommand{\headrulewidth}{0.4pt}
\fancypagestyle{plain}{\fancyhf{}\fancyfoot[C]{\small\thepage}\renewcommand{\headrulewidth}{0pt}}

\begin{document}

% =====================================================================
%  PAGE DE TITRE
% =====================================================================
\pagenumbering{Alph}
\begin{titlepage}
\thispagestyle{empty}
\centering
\vspace*{0.8cm}
{\Huge\bfseries La Lune et l'eau\par}
\vspace{0.5cm}
{\LARGE Disque secondaire circumterrestre\\[3pt]
et passage non collisionnel d'un corps planétaire massif\par}
\vspace{0.45cm}
{\Large Un chemin physique du système proto-terrestre\\ au système Terre--Lune\par}
\vspace{0.35cm}
{\normalsize Réservoirs, flux, passage planétaire non collisionnel,\\
géochimie, eau et domaine de viabilité\par}
\vspace{1.1cm}

\begin{tikzpicture}[scale=1.05]
  \fill[blue!30,opacity=0.55,even odd rule] (0,0) ellipse (3.5 and 0.62) (0,0) ellipse (1.25 and 0.22);
  \shade[ball color=orange!75!brown] (0,0) ellipse (1.02 and 0.9);
  \begin{scope}
    \clip (-4,-1.2) rectangle (4,0);
    \fill[blue!30,opacity=0.55,even odd rule] (0,0) ellipse (3.5 and 0.62) (0,0) ellipse (1.25 and 0.22);
  \end{scope}
  \draw[dashed,thick,gray!60!black,-{Stealth[length=3mm]}]
      (6.2,3.5) .. controls (2.0,2.3) and (2.0,1.3) .. (6.2,-0.1);
  \shade[ball color=gray!55] (3.0,1.72) circle (0.38);
  \node[font=\small] at (2.15,2.2) {planète $B$};
  \node[font=\small] at (0,-1.35) {proto-Terre};
  \node[font=\small,blue!45!black] at (-2.9,0.95) {disque secondaire};
  \draw[->,thick,red!60!black] (2.55,1.45) -- (1.5,0.75);
  \node[font=\scriptsize,red!60!black,align=center] at (1.35,1.55) {couple\\gravitationnel};
\end{tikzpicture}

\vspace{1.2cm}
{\large\bfseries Michel Debailleul\par}
\vspace{0.15cm}
Géophysicien --- Master, Université libre de Bruxelles (ULB)\\
Chercheur indépendant\\
ORCID : 0009-0003-1222-1433 \quad | \quad \href{mailto:michel.debailleul@yahoo.fr}{michel.debailleul@yahoo.fr}\\[4pt]
\copyright\ 2026 Michel Debailleul --- Licence : CC BY-NC 4.0
\vfill
\end{titlepage}
\pagenumbering{roman}

% =====================================================================
%  RÉSUMÉ
% =====================================================================
\thispagestyle{plain}
\begin{center}{\Large\bfseries Résumé}\end{center}
\medskip

Le système Terre--Lune présente aujourd'hui une architecture dynamique, une distribution de
masse, un moment cinétique, des signatures isotopiques, des compositions géochimiques, une
teneur en eau et une chronologie qui constituent autant de contraintes observationnelles
indépendantes. La question traitée ici n'est pas de décrire ces propriétés, mais d'identifier un
\emph{chemin physique} continu reliant l'état initial du système proto-terrestre à l'état observé.

L'ancrage astrophysique de cette démarche est que des satellites naturels sont déjà interprétés
comme les produits de réservoirs circumplanétaires : les satellites réguliers de Jupiter et de
Saturne dans des disques circumplanétaires, et les satellites internes de Neptune dans un système
remanié où un disque de débris aurait suivi la capture de Triton. Le cas terrestre proposé ne
transpose pas leurs masses ni leurs compositions ; il reprend le principe physique d'un réservoir
lié à une planète, capable de transporter masse et moment cinétique puis de former des satellites.

Cet ouvrage explore un tel chemin. Un disque secondaire circumterrestre, contemporain de
l'accrétion de la proto-Terre, constitue un réservoir orbital de masse et de moment cinétique
avant la formation de la Lune. Un corps planétaire massif $B$, en passage rasant et non
collisionnel, n'apporte pas l'essentiel de la matière lunaire : il agit comme un \emph{agent
gravitationnel transitoire de redistribution des flux de masse et de moment cinétique} entre
quatre réservoirs --- la proto-Terre, le disque, le corps $B$ et le milieu extérieur.

La logique héliocentrique du modèle est que $B$ n'a pas à subsister dans le Système solaire
actuel. Il peut appartenir à la population des corps planétaires primitifs ensuite éliminés par
l'instabilité orbitale globale du Système solaire. Cette extension ne l'identifie pas nécessairement
à une cinquième géante de glace ; elle exige seulement que son passage rasant précède son
élimination dynamique, selon la condition $t_B<t_{\mathrm{ej}}$. Elle relie ainsi le problème
Terre--Lune à celui, plus général, de la migration et de l'instabilité planétaires primitives
\cite{Tsiganis2005,Gomes2005,Nesvorny2011,NesvornyMorbidelli2012,Liu2022,Clement2018}.

Le passage réalise un tri orbital. Une cohorte de matière liée acquiert un moment cinétique qui
la place, après circularisation, au-delà de la limite de Roche, dans la fenêtre où la Lune peut
s'accréter ; une fraction réaccrète sur la proto-Terre, une autre s'échappe. Un passage parallèle
à la bande dynamique de basses latitudes d'une proto-Terre oblate en rotation rapide peut
en outre injecter dans le disque une matière superficielle déjà appauvrie en métal et en éléments
sidérophiles, dont les traceurs W, V, Cr, Mn, FeO, O et Ti permettent de tester la lignée.

L'ouvrage formule successivement le problème scientifique, les fondements astrophysiques,
l'architecture des réservoirs, la chronologie physique, la dynamique du passage, l'évolution du
réservoir circumterrestre, la géochimie par familles de traceurs, le domaine de viabilité et les
prédictions, chacune accompagnée de sa méthode de vérification. Un calcul de Phase~B par
particules test, conduit avec la configuration de référence, place directement $18{,}5\,\%$ du
réservoir dans la fenêtre lunaire à $i_B=25^\circ$, avec $84\,\%$ de la masse conservée dans le
domaine circumterrestre et un plateau entre $25$ et $30^\circ$ ; le contrôle rétrograde ramène
cette fraction à $4{,}0\,\%$ et porte la fraction échappée à $42{,}5\,\%$.

Le même passage fournit un mécanisme physique d'apport d'eau. L'enveloppe de volatils d'un corps
venu des régions froides est dilatée par le rayonnement de l'océan magmatique jusqu'à déborder son
propre lobe de Roche, produisant un flux vers le système terrestre et un autre vers une traînée
dont la réaccrétion s'étale sur des dizaines de millions d'années. L'eau lunaire, dont la
signature isotopique est indiscernable de l'eau terrestre, appartient à la même lignée matérielle
que l'oxygène, le titane, le chrome et le tungstène. La théorie ne prétend pas démontrer une
histoire unique : elle définit un chemin physique quantitatif et falsifiable.

\bigskip
\noindent\textbf{Mots-clés :} formation de la Lune ; eau lunaire ; volatils ; rapport D/H ;
disque circumterrestre ; accrétion secondaire ; rencontre non collisionnelle ; tri orbital ;
moment cinétique ; limite de Roche ; géochimie lunaire ; instabilité orbitale primitive ;
planète éliminée.

\bigskip
\begin{otherlanguage}{english}
\begin{center}{\Large\bfseries Abstract}\end{center}
\medskip

The Earth--Moon system currently exhibits a dynamical architecture, a mass distribution, an
angular momentum, isotopic signatures, geochemical compositions, a water content and a chronology
that together form independent observational constraints. The question addressed here is not to
describe these properties, but to identify a continuous \emph{physical path} linking the initial
state of the proto-terrestrial system to the observed state.

The astrophysical anchor of this approach is that natural satellites are already interpreted as
the products of circumplanetary reservoirs: the regular satellites of Jupiter and Saturn in
circumplanetary disks, and the inner satellites of Neptune in a reworked system where a debris
disk probably followed the capture of Triton. The proposed terrestrial case does not import their
masses or compositions; it retains the physical principle of a reservoir bound to a planet, able
to transport mass and angular momentum and then to form satellites.

This work explores such a path. A secondary circumterrestrial disk, contemporaneous with the
accretion of the proto-Earth, forms an orbital reservoir of mass and angular momentum before the
Moon forms. A massive planetary body $B$, on a grazing and non-collisional passage, does not
supply the bulk of the lunar material: it acts as a \emph{transient gravitational agent
redistributing mass and angular-momentum fluxes} between four reservoirs --- the proto-Earth, the
disk, body $B$ and the external medium.

The heliocentric logic of the model is that $B$ need not survive in the present Solar System. It
may belong to the population of primordial planetary bodies later eliminated by the global orbital
instability of the Solar System. This extension does not necessarily identify it with a fifth ice
giant; it only requires its grazing passage to precede its dynamical elimination, according to
$t_B<t_{\mathrm{ej}}$. It thus links the Earth--Moon problem to the broader problem of primordial
planetary migration and instability
\cite{Tsiganis2005,Gomes2005,Nesvorny2011,NesvornyMorbidelli2012,Liu2022,Clement2018}.

The passage performs an orbital sorting. A bound cohort of matter acquires an angular momentum
that places it, after circularization, beyond the Roche limit, in the window where the Moon can
accrete; one fraction re-accretes onto the proto-Earth, another escapes. A passage parallel to the
dynamical low-latitude band of an oblate rapidly rotating proto-Earth can also inject into the disk
surface material already depleted in metal and siderophile elements; W, V, Cr, Mn, FeO, O and Ti
then provide tracers with which to test the lineage.

The work successively formulates the scientific problem, the astrophysical foundations, the
reservoir architecture, the physical chronology, the passage dynamics, the evolution of the
circumterrestrial reservoir, geochemistry by tracer families, the viability domain and the
predictions, each with a verification method. A Phase~B test-particle calculation carried out with
the reference configuration places $18.5\,\%$ of the reservoir directly in the lunar window at
$i_B=25^\circ$, with $84\,\%$ of the mass retained in the circumterrestrial domain and a plateau
between $25$ and $30^\circ$; the retrograde control reduces this fraction to $4.0\,\%$ and raises
the escaping fraction to $42.5\,\%$.

The same passage provides a physical mechanism for the supply of water. The volatile envelope of a
body arriving from the cold regions is dilated by the radiation of the magma ocean until it
overflows its own Roche lobe, producing one flow toward the terrestrial system and one into a
trail whose re-accretion is spread over tens of millions of years. Lunar water, whose isotopic
signature is indistinguishable from terrestrial water, belongs to the same material lineage as
oxygen, titanium, chromium and tungsten. The theory does not claim to prove a unique history: it
defines a quantitative and falsifiable physical path.

\bigskip
\noindent\textbf{Keywords:} Moon formation; lunar water; volatiles; D/H ratio; circumterrestrial
disk; secondary accretion; non-collisional encounter; orbital sorting; angular momentum; Roche
limit; lunar geochemistry; primordial orbital instability; eliminated planet.
\end{otherlanguage}

% =====================================================================
%  PRÉAMBULE MÉTHODOLOGIQUE
% =====================================================================
\clearpage
\thispagestyle{plain}
\phantomsection
\addcontentsline{toc}{chapter}{Préambule méthodologique}
\begin{center}
{\Large\bfseries Préambule méthodologique}\\[4pt]
{\large\bfseries Une théorie de chemin rétrodictive}
\end{center}
\medskip

La présente monographie ne prétend pas reconstruire de manière unique ce qui
s'est réellement produit dans l'histoire du Système solaire. Une telle
reconstruction serait probablement inaccessible : les conditions initiales
exactes du système proto-terrestre, les trajectoires individuelles des corps
planétaires, l'état thermodynamique détaillé des réservoirs et la chronologie
fine des rencontres ont disparu.

L'objectif est différent. Il s'agit de construire un \emph{chemin physique
possible} reliant un système proto-terrestre admissible au système Terre--Lune
observé. La démarche est donc rétrodictive : elle part de l'état final connu,
puis remonte, contrainte après contrainte, vers une famille d'états initiaux et
d'étapes dynamiques capables de produire cet état final.

Le problème traité n'est donc pas :

\begin{center}
\emph{Quel événement unique s'est réellement produit ?}
\end{center}

mais plutôt :

\begin{center}
\emph{Existe-t-il un chemin physique continu, quantitativement contraint et
falsifiable, reliant un état proto-terrestre possible à l'état Terre--Lune
observé ?}
\end{center}

Ce statut peut s'écrire sous la forme
\begin{equation}
\boxed{
\exists\,\Gamma(t)\in\mathcal{D}_{\mathrm{viable}}
\quad\text{tel que}\quad
\Gamma(t_0)=\mathcal{S}_0,\qquad
\Gamma(t_f)=\mathcal{S}_f .
}
\label{eq:statut-retrodictif}
\end{equation}

La trajectoire $\Gamma(t)$ ne désigne pas nécessairement l'histoire réelle du
Système solaire. Elle désigne un chemin admissible dans l'espace des états
physiques : masse, moment cinétique, énergie, composition, phases de la matière
et chronologie. La valeur scientifique de la théorie ne repose donc pas sur une
revendication d'unicité historique, mais sur l'existence éventuelle d'un domaine
continu de paramètres satisfaisant simultanément les contraintes dynamiques,
géochimiques, isotopiques, thermodynamiques et chronologiques.

Dans ce cadre, le disque secondaire circumterrestre n'est pas introduit comme
un artifice narratif, mais comme un réservoir physique à tester. De même, le
corps planétaire massif $B$ n'apporte pas l'essentiel de la matière
lunaire : il est un agent gravitationnel transitoire de redistribution des
flux de masse et de moment cinétique. Son destin ultérieur appartient au même
cadre rétrodictif. Il peut appartenir à la population des corps planétaires du
Système solaire primitif ensuite éliminés par collision, accrétion, diffusion
gravitationnelle ou éjection.

La condition minimale associée à cette extension héliocentrique est
\begin{equation}
\boxed{
B\in\mathcal{P}_{\mathrm{elim}},
\qquad
 t_B<t_{\mathrm{ej}} .
}
\label{eq:B-elim-preambule}
\end{equation}

Ici, $\mathcal{P}_{\mathrm{elim}}$ désigne la population des corps planétaires
présents dans le Système solaire primitif mais absents du système actuel ;
$t_B$ est l'instant du passage rasant près de la proto-Terre ; et
$t_{\mathrm{ej}}$ l'instant de l'élimination dynamique de $B$. Cette formulation
laisse ouvertes les deux possibilités : une instabilité orbitale précoce ou une
instabilité plus tardive.

\begin{principe}{Statut de la théorie}
La théorie du disque secondaire circumterrestre est une théorie de chemin, non
une revendication d'unicité historique. Elle devient crédible si un domaine
continu de paramètres permet de franchir ensemble les portes dynamiques,
géochimiques, isotopiques, thermodynamiques et chronologiques. Elle devient
falsifiée si aucun tel domaine n'existe.
\end{principe}

\paragraph{Une filiation avec un mécanisme publié.}
Les rencontres non collisionnelles dans le système Terre--Lune primitif ont déjà
été invoquées pour expliquer l'inclinaison orbitale lunaire
\cite{PahlevanMorbidelli2015}. Le mécanisme proposé ici agit à un moment
différent de la même histoire : il ne s'applique pas à une Lune déjà formée, mais
au réservoir circumterrestre qui la précède. Les deux mécanismes ne se
concurrencent pas ; ils occupent des phases distinctes de la même évolution. Le
premier a passé le comité de lecture le plus exigeant ; le second en est
l'extension antérieure, et c'est cette filiation qui autorise à présenter le
passage rasant non collisionnel comme un processus de la physique connue, et non
comme une construction ad hoc.

% =====================================================================
\clearpage
\pagenumbering{arabic}
\part{Le problème scientifique}
% =====================================================================

\noindent
Cette première partie pose la question à laquelle l'ouvrage répond. Elle ne commence ni par le
disque ni par le corps $B$, mais par le système Terre--Lune tel qu'il est observé, et par
l'exigence que ces observations imposent à toute histoire de sa formation. Elle établit aussi un
point d'ancrage essentiel : la nature connaît déjà des satellites issus de réservoirs
circumplanétaires. Le disque secondaire, le corps $B$, les flux entre réservoirs, la géochimie,
l'eau et la dynamique apparaissent ensuite comme les éléments d'une même réponse.

% ---------------------------------------------------------------------
\chapter{Du système proto-terrestre au système Terre--Lune}

La théorie proposée résout un problème inverse. L'état final est connu et non négociable : la Lune
existe, avec sa masse, son déficit de fer, sa composition isotopique, sa teneur en eau et le moment
cinétique du couple qu'elle forme avec la Terre. Le travail consiste donc à remonter aux conditions
capables d'y conduire, et à les déterminer à trois instants : l'état de la proto-Terre et du disque
secondaire avant la rencontre, la configuration du passage du corps planétaire massif $B$ pendant
celle-ci, et l'état du système immédiatement après, qui sert de point de départ à l'évolution
maréale ultérieure. Ces trois états sont les inconnues du problème, et le domaine viable est
l'ensemble de leurs valeurs mutuellement compatibles.
\label{ch:probleme}
% ---------------------------------------------------------------------

\section{Le problème scientifique}

La Lune existe. Le système Terre--Lune présente aujourd'hui une architecture dynamique, une
distribution de masse, un moment cinétique, des signatures isotopiques, des compositions
géochimiques, une teneur en eau et une chronologie qui constituent autant de contraintes
observationnelles indépendantes.

Ces observations sont désormais suffisamment nombreuses et suffisamment précises pour
imposer un cadre physique exigeant à toute théorie de la formation lunaire. Aucune de ces
contraintes ne peut être considérée isolément : la dynamique orbitale, la géochimie, les
systèmes isotopiques, la thermodynamique, la teneur en eau et la chronologie décrivent un même
système physique et doivent pouvoir être expliquées de manière cohérente.

La question scientifique n'est donc pas de savoir si la Lune existe ni quelles sont aujourd'hui
ses propriétés, mais de déterminer :

\begin{center}
\emph{Quel chemin physique a conduit naturellement du système proto-terrestre\\ au système
Terre--Lune actuellement observé ?}
\end{center}

Autrement dit, il s'agit d'identifier une évolution physiquement admissible reliant un état
initial de la proto-Terre à l'état final représenté par le système Terre--Lune, tout en respectant
simultanément les contraintes imposées par les observations. Cette démarche se représente de
manière générale par
\begin{equation}
\boxed{\;\mathcal{S}_0\longrightarrow\mathcal{S}_f\;}
\label{eq:S0Sf}
\end{equation}
où $\mathcal{S}_0$ désigne l'état initial du système proto-terrestre et $\mathcal{S}_f$ l'état
actuellement observé. Le problème scientifique consiste à déterminer s'il existe un chemin
physique continu reliant ces deux états.

\begin{figure}[htbp]
\centering
\begin{tikzpicture}[node distance=0.45cm,
  st/.style={draw,rounded corners=3pt,align=center,font=\small,minimum height=1.15cm,
             text width=2.25cm,fill=#1},
  st/.default=white,
  ts/.style={draw,dashed,rounded corners=2pt,align=center,font=\scriptsize,text width=2.25cm,
             fill=green!4}]
  \node[st=blue!8] (s0) {$\mathcal{S}_0$\\proto-Terre\\+ réservoir $\RD$};
  \node[st=violet!10,right=of s0] (s1) {passage\\non collisionnel\\de $B$};
  \node[st=blue!8,right=of s1] (s2) {réservoir\\réorganisé\\(tri orbital)};
  \node[st=teal!8,right=of s2] (s3) {accrétion\\lunaire};
  \node[st=orange!10,right=of s3] (sf) {$\mathcal{S}_f$\\système\\Terre--Lune};
  \foreach \a/\b in {s0/s1,s1/s2,s2/s3,s3/sf} \draw[-{Stealth},thick] (\a) -- (\b);
  \node[ts,below=0.9cm of s1] (t1) {moment cinétique\\et masse};
  \node[ts,below=0.9cm of s2] (t2) {dynamique\\et architecture};
  \node[ts,below=0.9cm of s3] (t3) {géochimie\\et isotopes};
  \node[ts,below=0.9cm of sf] (t4) {chronologie};
  \foreach \t/\s in {t1/s1,t2/s2,t3/s3,t4/sf} \draw[-{Stealth},dashed,green!40!black] (\t) -- (\s);
  \node[font=\scriptsize,green!40!black,anchor=east] at ($(t1.west)+(-0.1,0)$) {tests indépendants};
\end{tikzpicture}
\caption{Le problème posé comme la recherche d'un chemin physique continu entre l'état initial
$\mathcal{S}_0$ et l'état observé $\mathcal{S}_f$. Chaque famille d'observations teste le chemin de
façon indépendante.}
\label{fig:chemin}
\end{figure}

\section{Les contraintes observationnelles}

Le tableau~\ref{tab:contraintes} rassemble les familles de contraintes que le chemin doit
satisfaire simultanément. Elles serviront, dans la partie~\ref{part:predictions}, de tests
indépendants.

\begin{table}[htbp]
\centering
\small
\caption{Familles de contraintes observationnelles imposées au chemin physique.}
\label{tab:contraintes}
\begin{tabularx}{\textwidth}{>{\raggedright\arraybackslash}p{2.6cm}
>{\raggedright\arraybackslash}X >{\raggedright\arraybackslash}p{5.4cm}}
\toprule
\textbf{Famille} & \textbf{Contrainte} & \textbf{Nature ou ordre de grandeur} \\
\midrule
Masse & rapport $M_L/\Mt$ & $1{,}23\times10^{-2}$ \\
Architecture & un satellite dominant & $f_{\mathrm{dom}}\simeq1$ \\
Moment cinétique & moment total du système & $L_{\mathrm{EM}}\approx3{,}5\times10^{34}$ kg\,m$^2$\,s$^{-1}$ \\
Orbite & inclinaison orbitale lunaire & environ $5^\circ$ sur l'écliptique \\
Structure interne & petit noyau lunaire & quelques pour cent de la masse au plus \\
Isotopes & O, Ti, Cr, W & Terre et Lune très proches, souvent indiscernables \\
Sidérophiles & V, Cr, Mn, W & appauvrissements semblables à ceux du manteau terrestre \\
Fer & FeO du manteau lunaire & plus élevé que dans le manteau terrestre \\
Volatils & K/U, Zn, isotopes & Lune appauvrie en éléments volatils \\
Eau & H$_2$O pré-éruptive, D/H & quelques centaines de ppm ; D/H terrestre \\
Chronologie & âge de formation & plusieurs dizaines de millions d'années après les premiers solides \\
\bottomrule
\end{tabularx}
\end{table}


\section{Ancrage astrophysique : des satellites naturels issus de disques}

Avant d'introduire le cas terrestre, un fait astrophysique doit être posé clairement : la formation
de satellites à partir d'un réservoir circumplanétaire n'est pas une construction ad hoc. Dans le
Système solaire, plusieurs familles de satellites naturels sont interprétées comme les produits
d'une matière liée à une planète, organisée dans un disque, un anneau massif ou un disque de
débris, puis transformée en satellites par accrétion, migration, dissipation et franchissement de
la limite de Roche.

\begin{principe}{Ancrage observationnel}
La théorie ne part pas de l'idée abstraite qu'un disque pourrait former une Lune. Elle part d'un
principe observé à plusieurs reprises :
\begin{equation}
\boxed{\;\text{planète en croissance ou remaniée}+\text{réservoir circumplanétaire}
\longrightarrow \text{satellites naturels}.\;}
\label{eq:ancrage-satellites}
\end{equation}
Le cas terrestre doit ensuite être traité avec ses propres contraintes : une planète tellurique,
un réservoir dominé par les solides, une masse lunaire élevée et un passage non collisionnel de
$B$ jouant le rôle d'agent de redistribution.
\end{principe}

Trois exemples jouent un rôle particulier dans le raisonnement. Le système galiléen de Jupiter
constitue l'exemple le plus classique d'un système régulier dont la formation est liée à un disque
circumjovien alimenté pendant la croissance de la planète. Les modèles discutent l'influx de gaz
et de solides, le transport de moment cinétique, la migration et les conditions d'accrétion des
satellites \cite{CanupWard2002,MosqueiraEstrada2003a,CanupWard2006,Nimmo2026}. Saturne apporte
un second appui : ses satellites réguliers et son système d'anneaux montrent qu'un réservoir
circumplanétaire peut transporter de la matière, s'étaler visqueusement, franchir la limite de
Roche et former des satellites \cite{Charnoz2010,CridaCharnoz2012,Blanc2025}. Neptune enfin est
un cas plus subtil : Triton est très probablement capturé, mais son arrivée a pu détruire ou
remanier un système régulier antérieur et produire un disque de débris à partir duquel les
satellites internes actuels se sont reformés ou réaccrétés
\cite{AgnorHamilton2006,BanfieldMurray1992,Showalter2019,Szulagyi2018,Davis2026}.

\begin{table}[htbp]
\centering
\small
\caption{Exemples naturels montrant qu'un réservoir circumplanétaire peut former ou reformer des
satellites. Ces exemples établissent le principe physique ; ils ne fixent pas les paramètres du
cas terrestre.}
\label{tab:ancrage-satellites}
\begin{tabularx}{\textwidth}{>{\raggedright\arraybackslash}p{2.4cm}
>{\raggedright\arraybackslash}X >{\raggedright\arraybackslash}X}
\toprule
\textbf{Système} & \textbf{Réservoir ou processus invoqué} & \textbf{Leçon physique pour le chemin proposé} \\
\midrule
Jupiter & disque circumjovien alimenté en gaz et solides pendant la croissance de Jupiter &
un disque circumplanétaire peut produire plusieurs satellites réguliers organisés \\
Saturne & disque circumplanétaire, anneaux massifs, étalement visqueux et franchissement de Roche &
un réservoir particulaire peut migrer, franchir Roche et alimenter une accrétion satellitaire \\
Neptune & capture de Triton, destruction ou remaniement d'un système antérieur, disque de débris et
réaccrétion des satellites internes &
un événement gravitationnel extérieur peut réorganiser un système circumplanétaire et laisser un
réservoir à partir duquel des satellites se reforment \\
\bottomrule
\end{tabularx}
\end{table}

Cette distinction est essentielle. Les satellites de Jupiter, Saturne et Neptune ne démontrent pas
l'histoire de la Lune. Ils démontrent que la nature dispose déjà du mécanisme de base : stocker de
la matière dans un réservoir circumplanétaire, lui faire transporter du moment cinétique, puis
convertir une partie de ce réservoir en satellites. La théorie proposée applique ce principe au
cas proto-terrestre et ajoute une étape propre : le passage rasant non collisionnel de $B$, qui ne
crée pas le réservoir mais le trie et le réorganise.

\section{Hypothèse étudiée}

La présente théorie explore l'un des chemins possibles. Elle transpose au cas proto-terrestre le
principe physique établi par les systèmes satellitaires naturels : une planète peut posséder un
réservoir circumplanétaire capable de stocker de la matière, de transporter du moment cinétique
et de former des satellites. Elle repose sur l'hypothèse qu'un disque secondaire
circumterrestre, contemporain de l'accrétion de la proto-Terre, constituait un tel réservoir
orbital de matière et de moment cinétique avant la formation de la Lune.

Dans ce cadre, un corps planétaire massif $B$ n'apporte pas l'essentiel de la matière lunaire. Son
passage demeure non collisionnel et agit comme un agent gravitationnel transitoire de
redistribution des flux de masse et de moment cinétique entre les différents réservoirs du système
proto-terrestre.

\begin{definition}
\textbf{Rôle du corps planétaire $B$.} Le corps $B$ est un \emph{agent gravitationnel transitoire
de redistribution des flux de masse et de moment cinétique entre les réservoirs du système
proto-terrestre}. Il n'apporte pas l'essentiel de la matière lunaire. Son action commence et
s'achève avec son passage. Cette définition est la référence de tout l'ouvrage.
\end{definition}

La Lune résulte alors de l'évolution de ce système réorganisé, principalement à partir de matière
déjà présente dans le disque circumterrestre, éventuellement enrichie par une contribution
limitée provenant des couches superficielles de la proto-Terre et par les volatils cédés par $B$
lui-même.

\section{Destin héliocentrique de \texorpdfstring{$B$}{B} et planète éliminée}
\label{sec:destinB}

La définition de $B$ comme agent transitoire pose immédiatement une question légitime : si un
corps planétaire massif a frôlé la proto-Terre sans collision, où se trouve-t-il aujourd'hui ? La
théorie répond en raccordant le problème Terre--Lune à l'évolution globale du Système solaire
primitif. $B$ n'est pas nécessairement un corps encore présent ; il peut appartenir à la
population des corps planétaires éliminés lors des instabilités orbitales primitives.

On note cette population
\begin{equation}
\mathcal{P}_{\mathrm{elim}}=
\left\{\begin{array}{l}
\text{corps planétaires primitifs éliminés par collision, accrétion,}\\
\text{diffusion gravitationnelle ou éjection}
\end{array}\right\}.
\label{eq:Pelim}
\end{equation}
La proposition héliocentrique minimale est alors
\begin{equation}
\boxed{\;B\in\mathcal{P}_{\mathrm{elim}},\qquad t_B<t_{\mathrm{ej}}\;}
\label{eq:tB-tej}
\end{equation}
où $t_B$ est l'instant du passage rasant près de la proto-Terre, et $t_{\mathrm{ej}}$ l'instant de
l'éjection ou de l'élimination dynamique de $B$.

Cette formulation est volontairement plus robuste que l'énoncé fort « $B$ est la planète
manquante ». Les modèles d'instabilité du Système solaire externe montrent qu'une migration des
planètes géantes a pu remodeler l'architecture planétaire \cite{Tsiganis2005,Gomes2005}. Des
versions à cinq ou six planètes invoquent l'éjection d'une planète supplémentaire, souvent de type
géante de glace \cite{Nesvorny2011,NesvornyMorbidelli2012}. D'autres travaux placent l'instabilité
beaucoup plus tôt, au moment de la dissipation du disque gazeux, et montrent que les planètes
telluriques en croissance peuvent être sculptées par cette instabilité \cite{Liu2022,Clement2018}.
La théorie ne choisit donc pas, à ce stade, entre une instabilité précoce et une instabilité plus
tardive : elle impose seulement que le passage Terre--$B$ appartienne à l'histoire de $B$ avant son
élimination.

\begin{ideephys}
Le corps $B$ n'est pas un objet artificiellement ajouté au seul problème Terre--Lune. Il peut
être compris comme un membre temporaire du Système solaire primitif : il passe près de la
proto-Terre, réorganise le réservoir circumterrestre, puis rejoint la dynamique héliocentrique
globale qui élimine une partie des corps planétaires excédentaires. Le passage rasant et la
disparition ultérieure de $B$ deviennent ainsi deux étapes possibles d'une même histoire.
\end{ideephys}

Le chemin complet prend alors la forme
\begin{equation}
\boxed{\begin{gathered}
B_{\mathrm{proto}}\longrightarrow\text{passage rasant près de la proto-Terre}
\longrightarrow\mathcal{D}_1\\
\longrightarrow\text{diffusion héliocentrique}
\longrightarrow\text{élimination de }B.
\end{gathered}}
\label{eq:cheminB}
\end{equation}
Cette extension ne remplace pas le mécanisme central de la théorie. Elle lui donne une logique
plus large : le disque secondaire et le tri orbital expliquent la formation lunaire, tandis que
l'instabilité orbitale primitive fournit un destin naturel à l'agent transitoire $B$.

\section{Portée de la théorie}

Cette théorie n'est pas présentée comme l'unique histoire possible de la formation de la Lune.
Son objectif est plus limité et plus précis : déterminer si le chemin physique proposé définit un
domaine de paramètres compatible avec l'ensemble des contraintes dynamiques, géochimiques,
isotopiques, thermodynamiques, chronologiques et de teneur en eau actuellement connues.

Les différentes observations ne sont donc pas considérées comme des validations présupposées,
mais comme des tests indépendants permettant de confirmer, de restreindre ou de falsifier le
domaine de plausibilité du modèle.

Ainsi, même si toutes ces contraintes étaient satisfaites simultanément, le résultat établirait
l'existence d'un chemin physique viable et quantitativement cohérent parmi les scénarios
envisageables de formation de la Lune, sans prétendre démontrer une histoire unique.

\begin{ideephys}
L'objet de la théorie est le chemin physique. Le disque est une condition initiale, $B$ est l'agent
de la transformation, les flux entre réservoirs en sont le langage, la géochimie, la teneur en eau
et la dynamique en sont les tests. Le destin héliocentrique de $B$ n'est pas un supplément
narratif : il ferme la logique du rôle transitoire en reliant $B$ à la population des corps
planétaires éliminés. Chaque partie de l'ouvrage répond à une partie de la même question :
\emph{comment le système Terre--Lune observé a-t-il pu émerger par une évolution
physiquement cohérente ?}
\end{ideephys}

\paragraph{Une filiation avec un mécanisme publié.}
Les rencontres non collisionnelles dans le système Terre--Lune primitif ont déjà été invoquées
pour rendre compte de l'inclinaison orbitale lunaire \cite{PahlevanMorbidelli2015}. Le mécanisme
proposé ici agit à un moment différent de la même histoire : il ne s'applique pas à une Lune déjà
formée, mais au réservoir circumterrestre qui la précède. Les deux mécanismes ne se concurrencent
pas ; ils occupent des phases distinctes de la même évolution. Le premier a passé le comité de
lecture le plus exigeant ; le second en est l'extension antérieure, et c'est cette filiation qui
autorise à présenter le passage rasant non collisionnel comme un processus de la physique connue,
et non comme une construction ad hoc.

\section{Statut des trois éléments du chemin}

Le chemin proposé articule trois éléments dont le statut épistémique diffère, et il est utile de
les distinguer d'emblée : l'un est emprunté à un cadre déjà établi, le deuxième est propre à la
théorie et constitue ce qu'elle demande d'examiner, le troisième est une hypothèse de raccordement
qui ne porte aucune charge démonstrative.

Le \emph{disque secondaire} relève d'un cadre établi. L'idée qu'un réservoir circumterrestre
contemporain de l'accrétion puisse être le site de formation de la Lune est celle des modèles de
co-accrétion \cite{Ruskol1960,Weidenschilling1986}, et le principe général d'accrétion secondaire
est illustré par les systèmes satellitaires naturels (chapitre~\ref{sec:satellites}).

Le \emph{passage rasant non collisionnel d'un corps planétaire massif $B$} est l'élément propre à
la théorie. Le disque existe avant la rencontre ; $B$ n'a donc pas à le créer, et son action se
limite à agir sur lui par son champ gravitationnel. Cette économie de moyens est le trait
distinctif du chemin proposé : la matière est déjà là, le moment cinétique est déjà là, et l'agent
extérieur ne fait que redistribuer l'un et l'autre.

Le \emph{destin héliocentrique de $B$} est une hypothèse de raccordement : il ne sert pas à produire
la Lune, mais à inscrire l'agent transitoire dans l'histoire globale du Système solaire primitif.
L'identification forte de $B$ à une planète manquante précise n'est pas requise ; il suffit que $B$
appartienne à la population des corps planétaires ensuite éliminés, avec la condition
$ t_B<t_{\mathrm{ej}} $.

\begin{ideephys}[Le principe des circonstances précises]
La Lune existe, et elle a les propriétés particulières que l'observation lui connaît : une masse
inhabituelle pour un satellite de planète tellurique, un déficit de fer, une composition isotopique
indiscernable de celle du manteau terrestre, un moment cinétique élevé, une teneur en eau
comparable à celle du manteau supérieur terrestre. Ces propriétés ne se déduisent pas d'une loi
générale ; elles demandent des \emph{circonstances précises}. Le travail proposé consiste donc
moins à raconter ce qui s'est passé qu'à délimiter, dans l'espace des paramètres, la région où ces
circonstances sont simultanément réalisables. Cette région est le \emph{domaine viable}
$\mathcal{D}_{\mathrm{viable}}$ défini au chapitre~\ref{ch:portes}. Si elle est non vide, le chemin
proposé est physiquement praticable ; si elle est vide, il ne l'est pas. C'est en ce sens que la
théorie est une théorie de chemin, et c'est aussi en ce sens qu'elle est falsifiable.
\end{ideephys}

\begin{lecture}
Dans le chemin proposé, une interaction gravitationnelle sans contact suffit à réorganiser un
réservoir déjà constitué. $B$ conserve son identité physique et passe. Son passage introduit un
couple extérieur capable de redistribuer le moment cinétique du réservoir préexistant.
\end{lecture}


\section{Organisation de l'ouvrage}

L'ouvrage suit le chemin physique dans l'ordre où la question l'impose
(figure~\ref{fig:plan}). La partie~\ref{part:fondements} établit les fondements astrophysiques :
l'accrétion secondaire autour des planètes est un phénomène réel et répété. La
partie~\ref{part:architecture} définit les quatre réservoirs et les lois qui gouvernent leurs
échanges. La partie~\ref{part:chronologie} ordonne le chemin par ses échelles de temps. La
partie~\ref{part:dynamique} décrit l'action de $B$ et son résultat central, le tri orbital. La
partie~\ref{part:reservoir} suit le réservoir circumterrestre avant et après le passage. La
partie~\ref{part:geochimie} organise les tests géochimiques par familles de traceurs. La
partie~\ref{part:viabilite} rassemble les paramètres, les portes et le programme numérique. La
partie~\ref{part:predictions} énonce les prédictions sous la forme : \emph{si la théorie est
correcte\dots\ voici comment le vérifier}.
Le destin héliocentrique de $B$ traverse ces parties comme une contrainte de cohérence : l'agent
gravitationnel doit pouvoir passer près de la Terre, puis être diffusé ou éliminé ultérieurement
sans rester dans l'architecture planétaire actuelle.

\begin{figure}[htbp]
\centering
\begin{tikzpicture}[node distance=0.35cm and 0.5cm,
  pt/.style={draw,rounded corners=3pt,align=center,font=\scriptsize,minimum height=1.05cm,
             text width=3.6cm,fill=#1},pt/.default=white]
  \node[pt=gray!8] (p1) {\textbf{I} Le problème scientifique};
  \node[pt=gray!8,right=of p1] (p2) {\textbf{II} Fondements astrophysiques};
  \node[pt=blue!8,right=of p2] (p3) {\textbf{III} Architecture : quatre réservoirs};
  \node[pt=blue!8,below=of p3] (p4) {\textbf{IV} Chronologie physique};
  \node[pt=violet!10,left=of p4] (p5) {\textbf{V} Dynamique : l'agent $B$ et le tri orbital};
  \node[pt=blue!8,left=of p5] (p6) {\textbf{VI} Le réservoir évolutif $\RD$};
  \node[pt=green!8,below=of p6] (p7) {\textbf{VII} Géochimie par familles de traceurs};
  \node[pt=orange!10,right=of p7] (p8) {\textbf{VIII} Domaine de viabilité};
  \node[pt=cyan!8,right=of p8] (p9) {\textbf{IX} Prédictions et vérifications};
  \foreach \a/\b in {p1/p2,p2/p3,p3/p4,p4/p5,p5/p6,p6/p7,p7/p8,p8/p9}
    \draw[-{Stealth},thick] (\a) -- (\b);
\end{tikzpicture}
\caption{Architecture de l'ouvrage : chaque partie répond à une partie de la même question.}
\label{fig:plan}
\end{figure}
% =====================================================================
\part{Les fondements astrophysiques}
\label{part:fondements}
% =====================================================================

\noindent
Avant de proposer un chemin pour la Terre, il faut établir que son premier élément appartient à
la physique connue. Cette partie montre que l'accrétion secondaire autour des planètes est un
phénomène réel et répété. Jupiter et Saturne ancrent le principe des disques circumplanétaires
réguliers ; Neptune ajoute le cas d'un système remanié par la capture de Triton, où un disque de
débris et la réaccrétion de satellites internes jouent un rôle central. Ces analogies sont utiles
parce qu'elles établissent un mécanisme, mais elles ne fixent ni la masse, ni la composition, ni
la chronologie du cas terrestre. La partie se termine donc sur la spécificité de la proto-Terre :
un disque dominé par les solides, dont le mode d'alimentation diffère de celui des planètes
géantes.

% ---------------------------------------------------------------------
\chapter{Le principe d'accrétion secondaire}
\label{ch:accretion-secondaire}\label{sec:satellites}
% ---------------------------------------------------------------------

\section{Le principe circumplanétaire}

Les systèmes de satellites naturels fournissent un ancrage direct pour la notion de système
d'accrétion secondaire. Les grands satellites réguliers de Jupiter, de Saturne et d'Uranus sont
majoritairement progrades, peu excentriques et organisés autour du plan équatorial ou du plan
de Laplace de leur planète. Les satellites internes de Neptune, bien que probablement issus d'une
histoire secondaire après la capture de Triton, constituent un autre exemple de matière liée à une
planète, réorganisée en réservoir circumplanétaire puis réassemblée en satellites. Ces systèmes
sont des archives dynamiques d'une évolution circumplanétaire
\cite{CanupWard2002,CanupWard2006,Nimmo2026,Blanc2025,BanfieldMurray1992,Showalter2019,JPL}.
Les systèmes suivants sont retenus comme exemples de référence.

\begin{table}[htbp]
\centering
\caption{Systèmes satellitaires servant d'appui au principe d'accrétion circumplanétaire. La
comparaison porte sur le principe physique du disque, non sur une identité de masse ou de
thermodynamique avec la proto-Terre.}
\label{tab:systemes}
\begin{tabular}{ll}
\toprule
\textbf{Planète} & \textbf{Satellites retenus} \\
\midrule
Jupiter & Io, Europe, Ganymède, Callisto \\
Saturne & Mimas, Encelade, Téthys, Dioné, Rhéa, Titan, Japet \\
Uranus  & Miranda, Ariel, Umbriel, Titania, Obéron \\
\bottomrule
\end{tabular}
\end{table}

Les modèles de formation des satellites réguliers traitent précisément l'alimentation en
matière, le transport de moment cinétique, le refroidissement, la migration et l'accrétion dans
un disque entourant la planète centrale \cite{CanupWard2002,CanupWard2006,Nimmo2026}.

\begin{ideephys}
La théorie ne transpose pas à la proto-Terre la masse du disque de Jupiter ou de Saturne. Elle
retient un résultat plus fondamental : un corps planétaire en croissance peut être entouré d'un
réservoir en rotation qui transporte de la masse et du moment cinétique et qui peut devenir le
siège d'une accrétion satellitaire.
\end{ideephys}

\section{Sens physique du terme « secondaire »}

Le disque est qualifié de secondaire relativement au disque circumsolaire primaire. La
hiérarchie des deux systèmes d'accrétion s'écrit
\begin{equation}
\text{disque circumsolaire primaire}\longrightarrow\text{proto-Terre}\longrightarrow
\text{disque circumterrestre secondaire.}
\label{eq:hierarchie}
\end{equation}
Le mot « secondaire » décrit une hiérarchie gravitationnelle et accrétionnelle. Il ne signifie
pas que le disque apparaît tardivement, ni qu'il est créé par le corps $B$.


\section{Trois propositions emboîtées}

La théorie distingue trois niveaux logiques afin de séparer ce qui relève de la physique générale
des disques, ce que montrent les systèmes naturels et ce qui constitue la proposition
spécifiquement terrestre.

\begin{principe}{Principe d'accrétion secondaire circumplanétaire}
\textbf{P1 --- Physique générale.} Une planète en croissance peut recevoir de la matière porteuse
de moment cinétique et en organiser une fraction dans un réservoir circumplanétaire en rotation.
Ce réservoir transporte masse et moment cinétique et peut devenir un site d'accrétion
satellitaire \cite{Pringle1981,CanupWard2002,CanupWard2006}.

\smallskip
\textbf{P2 --- Manifestations naturelles.} Les systèmes réguliers de Jupiter, Saturne et Uranus
sont des expressions naturelles de l'accrétion secondaire circumplanétaire. Leur histoire
détaillée diffère d'un système à l'autre, mais leur architecture exige qu'une matière liée à la
planète ait été organisée et accrétée dans son environnement orbital
\cite{Nimmo2026,Blanc2025,Helled2020}.

\smallskip
\textbf{P3 --- Proposition terrestre.} La proto-Terre possédait un disque secondaire
circumterrestre lorsque le corps planétaire massif $B$ l'a frôlée sans collision. Cette rencontre
a réorganisé le réservoir et préparé l'accrétion d'un satellite dominant : la Lune.
\end{principe}

Cette séparation protège la logique de la théorie. Les systèmes jovien, saturnien et uranien
montrent plusieurs réalisations physiques du même principe général ; la branche terrestre doit
ensuite être calculée avec les paramètres propres à la proto-Terre et avec la perturbation
spécifique de $B$.

\section{Une hiérarchie de systèmes d'accrétion}

La hiérarchie générale s'écrit
\begin{equation}
\text{disque circumstellaire}\to\text{planète en croissance}\to
\text{disque circumplanétaire}\to\text{satellites}
\label{eq:hierarchie-gen}
\end{equation}
Le cas terrestre proposé est
\begin{equation}
\text{disque circumsolaire}\rightarrow\text{proto-Terre + disque secondaire}
\xrightarrow{\;\text{passage de }B\;}\text{disque réorganisé}\rightarrow\text{Lune}.
\label{eq:hierarchie-terre}
\end{equation}

\begin{figure}[htbp]
\centering
\begin{tikzpicture}[scale=0.95]
  \fill[yellow!25,even odd rule] (0,0) ellipse (6.2 and 1.2) (0,0) ellipse (1.0 and 0.2);
  \shade[ball color=yellow!80!orange] (0,0) circle (0.45);
  \node[font=\small] at (0,-0.8) {étoile};
  \fill[blue!30,even odd rule] (4.1,0.1) ellipse (1.1 and 0.28) (4.1,0.1) ellipse (0.35 and 0.09);
  \shade[ball color=orange!70!brown] (4.1,0.1) circle (0.22);
  \node[font=\small] at (4.1,-0.55) {planète};
  \node[font=\small,blue!45!black] at (4.3,0.75) {disque secondaire};
  \node[font=\small,align=center] at (-3.6,0.55) {disque circumstellaire\\primaire};
\end{tikzpicture}
\caption{Principe d'accrétion secondaire : le disque circumplanétaire est un système d'accrétion
emboîté dans le système circumstellaire primaire.}
\label{fig:hierarchie}
\end{figure}

\section{Jupiter : le système galiléen comme référence dynamique}

Io, Europe, Ganymède et Callisto forment un système compact de satellites progrades dont les
excentricités, de l'ordre de $10^{-3}$ à $10^{-2}$, et les inclinaisons relatives au plan de
Laplace jovien sont faibles \cite{JPL}. Cette organisation est cohérente avec une accrétion dans
un disque dissipatif capable d'amortir les mouvements aléatoires.

La revue récente du système galiléen \cite{Nimmo2026} identifie comme incertitudes majeures la
distribution du moment cinétique de la matière qui alimente le disque, l'origine des solides et
la structure thermique et visqueuse du disque. Ce sont précisément les grandeurs qui décrivent
le disque circumterrestre :
\begin{equation}
\mathcal{V}_D=\left\{\dot M_{\mathrm{in}},\ \Sigma(r,t),\ T(r,t),\ \nu(r,t),\ L_D(t),\
\mathbf{X}(r,t)\right\}.
\label{eq:VD}
\end{equation}
Le gradient de densité des satellites galiléens, avec une fraction de glace croissante vers
l'extérieur, montre qu'un disque satellitaire imprime une information thermochimique radiale
aux corps qui s'y forment. La fonction $T(r,t)$ doit donc être traitée comme une variable de
formation des matériaux, et pas seulement comme un paramètre dynamique.

\section{Saturne : transport visqueux, limite de Roche et naissance de satellites}

Les grandes lunes telles que Titan et Japet peuvent être reliées à une phase circumplanétaire
ancienne ; Japet, plus incliné que les satellites internes, est retenu comme cas limite
dynamiquement distinct. L'environnement des anneaux illustre un autre mécanisme : un
réservoir particulaire s'étale visqueusement et forme de petits satellites lorsqu'une fraction de
sa matière franchit la limite de Roche \cite{Charnoz2010,Blanc2025} :
\begin{equation}
\text{réservoir circumplanétaire}\to\text{étalement}\to
\text{franchissement de Roche}\to\text{accrétion satellitaire}
\label{eq:saturne}
\end{equation}
Cette séquence est directement utile au cas terrestre : la limite de Roche ne fixe pas le bord
interne du disque ; elle marque une transition dans la capacité de la matière à construire des
agrégats autogravitants durables.

\section{Uranus : un système régulier, une origine du réservoir encore discutée}

Miranda, Ariel, Umbriel, Titania et Obéron forment un système régulier organisé autour de
l'équateur d'Uranus. L'origine exacte de ce réservoir reste discutée : des modèles de disque
circumplanétaire primordial et de disque produit plus tardivement ont été proposés
\cite{Helled2020,Blanc2025}. Uranus est donc utilisé comme exemple du principe d'accrétion
dans un réservoir circumplanétaire, sans imposer d'identité d'origine avec le réservoir terrestre.


\section{Neptune : capture de Triton, disque de débris et réaccrétion}

Neptune n'apporte pas le même type de preuve que Jupiter ou Saturne. Son satellite dominant,
Triton, est rétrograde et fortement incliné ; il est donc généralement interprété comme un corps
capturé, probablement lors d'une rencontre gravitationnelle impliquant un binaire transneptunien
\cite{AgnorHamilton2006}. Cette capture a profondément remanié le système neptunien.

Le point important pour la présente théorie n'est pas Triton lui-même, mais les satellites
internes de Neptune. Dans le scénario de Banfield et Murray, l'orbite très excentrique de Triton
après sa capture excite les anciens satellites internes, provoque des collisions, forme un disque
de débris, puis laisse se reformer un système interne de satellites sur des orbites proches de
l'équateur \cite{BanfieldMurray1992}. La découverte et l'analyse d'Hippocampe renforcent l'idée
d'un système interne jeune, fragmenté et réassemblé, fortement marqué par les collisions
\cite{Showalter2019}. Des simulations de formation de satellites autour des géantes de glace
montrent en outre qu'Uranus et Neptune peuvent former des disques circumplanétaires capables de
produire des systèmes réguliers ; Neptune aurait donc pu posséder un système initial de type
uranien avant le remaniement provoqué par Triton \cite{Szulagyi2018}. Des observations récentes
des anneaux et lunes internes de Neptune par le JWST suggèrent même que leur matière pourrait
provenir d'intérieurs de corps glacés plus anciens détruits ou remaniés \cite{Davis2026}.

\begin{ideephys}
Le cas neptunien apporte une analogie différente de celle de Jupiter et de Saturne. Il montre
qu'un événement gravitationnel extérieur peut détruire, trier ou réorganiser un système de
matière circumplanétaire, puis laisser un réservoir à partir duquel des satellites internes se
reforment. Dans le chemin terrestre proposé, $B$ ne capture pas la Lune et ne fournit pas sa masse
principale ; il joue cependant un rôle analogue d'agent gravitationnel de réorganisation.
\end{ideephys}

\section{Observation directe : le disque circumplanétaire de PDS 70c}

ALMA a résolu autour de la jeune planète PDS~70c une source compacte de poussière
spatialement séparée du disque circumstellaire, interprétée comme un disque circumplanétaire
\cite{Benisty2021}. Le système fournit une observation directe de deux niveaux d'accrétion
emboîtés :
\begin{equation}
\text{disque circumstellaire de PDS~70}\ \supset\ \text{PDS~70c}+\text{disque circumplanétaire}.
\label{eq:pds70}
\end{equation}

\begin{ideephys}
PDS~70c montre qu'une planète encore en formation peut posséder son propre disque à
l'intérieur d'un disque circumstellaire plus vaste. La théorie transpose ce principe de
hiérarchie au cas de la proto-Terre, puis y ajoute la perturbation qui lui est propre : le passage
non collisionnel de $B$.
\end{ideephys}


% ---------------------------------------------------------------------
\chapter{Du principe au cas terrestre}
\label{ch:cas-terrestre}
% ---------------------------------------------------------------------

Les systèmes naturels établissent le principe ; ils ne fixent pas les valeurs terrestres. Ce
chapitre compare quantitativement les systèmes satellitaires, formule la contrainte d'un satellite
dominant et identifie la nature du disque terrestre.

\section{Comparaison quantitative des systèmes naturels}

On définit la fraction de masse satellitaire
\begin{equation}
\mu_{\mathrm{sat}}=\frac{\sum_i M_i}{M_p}.
\label{eq:musat}
\end{equation}
À partir des paramètres gravitationnels du JPL \cite{JPL}, on obtient les valeurs du
tableau~\ref{tab:musat}.

\begin{table}[htbp]
\centering
\caption{Fraction de masse des systèmes satellitaires retenus, calculée à partir des paramètres
$GM$ (indépendante de la valeur adoptée de $G$).}
\label{tab:musat}
\begin{tabular}{llc}
\toprule
\textbf{Système} & \textbf{Satellites inclus} & $\mu_{\mathrm{sat}}$ \\
\midrule
Jupiter & Io, Europe, Ganymède, Callisto & $2{,}07\times10^{-4}$ \\
Saturne & Mimas, Encelade, Téthys, Dioné, Rhéa, Titan, Japet & $2{,}47\times10^{-4}$ \\
Uranus  & Miranda, Ariel, Umbriel, Titania, Obéron & $1{,}04\times10^{-4}$ \\
Terre   & Lune & $1{,}23\times10^{-2}$ \\
\bottomrule
\end{tabular}
\end{table}

\begin{figure}[htbp]
\centering
\begin{tikzpicture}[x=5cm,y=1cm]
  \draw[thick,-{Stealth}] (-4.25,0) -- (-1.7,0) node[right] {$\mu_{\mathrm{sat}}$};
  \foreach \e/\lab in {-4.2/{-4{,}2},-3.8/{-3{,}8},-3.4/{-3{,}4},-3.0/{-3{,}0},-2.6/{-2{,}6},-2.2/{-2{,}2},-1.8/{-1{,}8}}
    \draw (\e,0.08) -- (\e,-0.08) node[below,font=\scriptsize] {$10^{\lab}$};
  \foreach \x/\y/\n/\v in {-3.983/0.8/Uranus/{1{,}04\times10^{-4}},-3.684/1.45/Jupiter/{2{,}07\times10^{-4}},-3.607/2.1/Saturne/{2{,}47\times10^{-4}},-1.910/1.45/{Terre--Lune}/{1{,}23\times10^{-2}}}{
    \draw[gray] (\x,0) -- (\x,\y-0.25);
    \fill[blue!50!black] (\x,0) circle (2pt);
    \node[font=\small,align=center] at (\x,\y) {\n\\[-2pt]{\scriptsize$\v$}};
  }
  \draw[{Stealth}-{Stealth},red!60!black] (-3.55,2.7) -- (-1.96,2.7);
  \node[font=\small,red!60!black] at (-2.75,3.0) {contraste d'environ 50 à 120};
\end{tikzpicture}
\caption{Contraste de masse entre les systèmes réguliers des planètes géantes et le système
Terre--Lune (échelle logarithmique). Ce contraste est une contrainte quantitative du cas
terrestre, non une objection au principe circumplanétaire.}
\label{fig:contraste}
\end{figure}

Les trois systèmes de planètes géantes se situent à l'ordre de grandeur $10^{-4}$, cohérent avec
le régime de régulation décrit pour des disques gazeux alimentés \cite{CanupWard2006}. On
définit le contraste de masse
\begin{equation}
C_M=\frac{M_L/\Mt}{\mu_{\text{sat,géant}}},
\qquad 50\lesssim C_M\lesssim 120.
\label{eq:CM}
\end{equation}

\begin{controle}
Cette différence ne réfute pas le principe d'accrétion secondaire. Elle fixe une contrainte
propre au cas terrestre : le disque circumterrestre, réorganisé par $B$, doit produire un résultat
bien plus concentré en masse satellitaire que les systèmes réguliers des planètes géantes. La
théorie doit expliquer quantitativement pourquoi un satellite dominant de masse lunaire émerge
du système terrestre.
\end{controle}


\begin{table}[htbp]
\centering
\small
\caption{Contrastes entre les disques circumplanétaires des géantes et le réservoir
circumterrestre proposé. Les différences ne sont pas des anomalies à corriger : elles définissent
le problème que la théorie cherche à résoudre.}
\label{tab:contraste-disques}
\begin{tabularx}{\textwidth}{>{\raggedright\arraybackslash}p{4.2cm}
>{\raggedright\arraybackslash}X >{\raggedright\arraybackslash}X}
\toprule
 & \textbf{Jupiter, Saturne} & \textbf{Terre (proposé)} \\
\midrule
Masse du disque & $\sim10^{-4}\,M_p$ & $\sim10^{-2}\,M_p$, à expliquer \\[3pt]
Phase dominante & gaz et solides & condensats, vapeur de silicate \\[3pt]
Support de pression & oui, assuré par le gaz & marginal \\[3pt]
Durée de vie & millions d'années & quelques années, éq.~\eqref{eq:tnu-reservoir} \\
\bottomrule
\end{tabularx}
\end{table}


\paragraph{Pourquoi le réservoir est dominé par les condensats.} Deux raisons indépendantes, qu'il
vaut la peine de distinguer.

La première est chronologique. Les disques protoplanétaires vivent trois à cinq millions d'années,
et la chronologie au tungstène place la formation lunaire plus de cinquante millions d'années après
les CAI. Un facteur dix à vingt sépare les deux : l'hydrogène et l'hélium nébulaires ont disparu
depuis longtemps, et aucun gaz léger n'est disponible. Ce point ne relève pas d'une hypothèse, il
est fixé par la datation.

La seconde est énergétique. La chaleur latente de vaporisation du silicate vaut environ
$1{,}2\times10^{7}$~J\,kg$^{-1}$, tandis que l'énergie thermique d'un magma à $3000$~K n'en
représente qu'un dixième. Une matière isolée ne peut donc vaporiser qu'une fraction d'elle-même,
de l'ordre de $10$ à $20\,\%$, la vaporisation étant endothermique et se refroidissant en cours de
route. Vaporiser l'ensemble demanderait de déposer une énergie comparable à l'énergie de liaison,
soit $3\times10^{7}$~J\,kg$^{-1}$ à $1{,}7\,\Rt$. Or le passage n'en dépose aucune : la matière est
soulevée exactement à la limite où elle cesse d'être liée. \emph{La nature du réservoir découle donc
directement du caractère non collisionnel du mécanisme}, et c'est une conséquence de la théorie
plutôt qu'une hypothèse ajoutée.

\paragraph{Mais le réservoir n'est pas dans le vide.} Il baigne dans le rayonnement d'un océan
magmatique. Un élément situé à la distance $r$ d'une surface portée à la température $T_s$ atteint
l'équilibre radiatif à $T_d\simeq T_s\sqrt{\Req/r}/\sqrt2$. Pour $\Req=1{,}7\,\Rt$, le disque à la
limite de Roche est à $0{,}54\,T_s$. Le silicate condensant entre $1700$ et $2100$~K, il en résulte
un critère simple : le réservoir reste vapeur jusqu'à $\aR$ si la surface dépasse environ
$3700$~K, et il condense en gouttelettes en deçà. Ce seuil est la forme quantitative que prend la
porte $\mathcal{D}_{\mathrm{thermo}}$.

Le réservoir proposé est donc un milieu à deux phases : dominé en masse par les condensats, pour
les raisons énergétiques ci-dessus, mais avec une composante vapeur qui contrôle la pression et la
viscosité et dont l'abondance dépend directement de la température de surface. Le mot
\emph{condensats} est préféré à \emph{solides} à dessein : ce qui quitte l'océan magmatique est
liquide, et ne se solidifie qu'ensuite par rayonnement une fois en orbite.

\begin{ideephys}[Ce que le tri orbital vient remplacer]
Dans un disque circumjovien, c'est le gaz qui fournit le support de pression, la viscosité et la
durée : il laisse au système des millions d'années pour transporter lentement le moment cinétique
vers l'extérieur et former des satellites. Le réservoir circumterrestre proposé n'a rien de tout
cela. Il est dominé par les solides, sans support de pression, et sa durée de vie se compte en
années. Il lui faut donc un mécanisme qui accomplisse en une fois ce que le gaz accomplit
lentement. C'est la fonction que l'on propose d'attribuer au passage de $B$ : le tri orbital est
l'équivalent impulsif du transport visqueux. Les quatre contrastes du
tableau~\ref{tab:contraste-disques} ne sont donc pas des faiblesses de l'analogie ; ils sont la
raison d'être du mécanisme proposé.
\end{ideephys}

\section{De plusieurs satellites à un satellite dominant}

Une mesure simple de concentration de la masse satellitaire est
\begin{equation}
f_{\mathrm{dom}}=\frac{M_{\max}}{\sum_i M_i}.
\label{eq:fdom}
\end{equation}
Dans le système terrestre actuel, $f_{\mathrm{dom}}\simeq1$. Le rôle dynamique spécifique de
$B$ s'exprime sous la forme
\begin{equation}
\mathcal{D}_0\xrightarrow{\;B\;}\mathcal{D}_1
\xrightarrow{\;\text{dissipation + accrétion}\;}\left\{M_L,\ f_{\mathrm{dom}}\simeq1\right\}.
\label{eq:roleB}
\end{equation}

\begin{falsif}
Si, dans le domaine de paramètres qui respecte les autres contraintes, le disque post-rencontre
produit systématiquement plusieurs satellites de masse comparable ou une masse satellitaire
totale très inférieure à la masse lunaire, le mécanisme de réorganisation proposé est
insuffisant.
\end{falsif}

\section{Spécificité du cas terrestre : un disque dominé par les solides}
\label{sec:soustherm}

Jupiter, Saturne et PDS~70c dépassent la \emph{masse thermique} de leur disque,
\begin{equation}
M_{\mathrm{th}}\simeq\left(\frac{H}{r}\right)^{3}M_\odot\approx10\,\Mt
\quad\text{à 1 UA},
\label{eq:Mth}
\end{equation}
au-delà de laquelle le gaz capté forme un disque en rotation. La proto-Terre, dix fois moins
massive, est dans le régime \emph{sous-thermique}. Son rayon de Bondi,
\begin{equation}
R_{\mathrm{B}}=\frac{G\Mt}{c_s^2}\approx4{,}0\times10^{8}\ \mathrm{m}\approx63\,\Rt
\qquad(T\approx280\ \mathrm{K},\ \mu=2{,}34),
\label{eq:Bondi}
\end{equation}
est inférieur à son rayon de Hill ($\approx235\,\Rt$), et c'est lui qui fixe la zone
d'accrétion ; les simulations à haute résolution montrent que les galets qui franchissent la
sphère de Bondi d'un embryon rocheux sont effectivement accrétés \cite{Popovas2018}.

Le rayon de circularisation d'une matière de moment cinétique spécifique
$j\simeq\Omega R_{\mathrm{acc}}^2/4$ vaut $r_c=j^2/(G\Mt)$ \cite{QuillenTrilling1998} :
\begin{equation}
r_c(R_{\mathrm{acc}}=R_H)=\frac{R_H}{48}\approx4{,}9\,\Rt,
\qquad
r_c(R_{\mathrm{acc}}=R_{\mathrm{B}})=\left(\frac{R_{\mathrm{B}}}{R_H}\right)^{4}
\frac{R_H}{48}\approx0{,}03\,\Rt.
\label{eq:rc}
\end{equation}
Le gaz nébulaire capté à travers la sphère de Bondi retombe donc sur la planète au lieu de
former un disque, conformément aux simulations d'enveloppes d'embryons sous-thermiques,
soutenues par la pression plutôt que par la rotation \cite{Ormel2015,Fung2015}. Le disque
secondaire terrestre est en conséquence dominé par les solides, avec trois canaux
d'alimentation :
\begin{equation}
\dot M_{\mathrm{in}}=\dot M_{\mathrm{coacc}}+\dot M_{\mathrm{ej}}+\dot M_{\mathrm{rot}},
\label{eq:canaux}
\end{equation}
où $\dot M_{\mathrm{coacc}}$ représente la capture par collisions de planétésimaux dans la sphère
de Hill \cite{Ruskol1960,Weidenschilling1986}, $\dot M_{\mathrm{ej}}$ les éjectas des impacts
ordinaires d'accrétion et $\dot M_{\mathrm{rot}}$ la perte de masse équatoriale d'une proto-Terre
en rotation proche de la limite critique.

\begin{lecture}
L'analogie avec les planètes géantes porte sur le principe d'un réservoir circumplanétaire, non
sur son mode d'alimentation. Le canal rotationnel $\dot M_{\mathrm{rot}}$ est le plus directement
couplé au reste de la théorie : il fournit une matière de composition silicatée terrestre, orientée
selon le spin, et relie la rotation rapide de la proto-Terre à l'existence même du disque.
\end{lecture}

% =====================================================================

\begin{rappel}
La Partie~\ref{part:fondements} a établi deux résultats. D'une part, un réservoir circumplanétaire
en rotation est un objet astrophysique réel. D'autre part, dans le cas terrestre, ce réservoir est
dominé par les solides et doit produire une masse satellitaire concentrée dans un corps unique.
La partie suivante traduit ces résultats dans un langage physique commun : celui des réservoirs et
des flux.
\end{rappel}
% =====================================================================
\part{Architecture physique : les quatre réservoirs}
\label{part:architecture}
% =====================================================================

\noindent
Le chemin physique est décrit comme une succession d'échanges entre réservoirs. Ce langage a
un avantage décisif : il impose de fermer les bilans de masse, de moment cinétique, d'énergie et
de composition à chaque étape, et il donne au rôle de $B$ une formulation précise. Cette partie
définit les réservoirs, les flux qui les relient et les lois de conservation qui les contraignent,
puis fixe les variables d'état du Niveau~I et les échelles numériques de référence.

% ---------------------------------------------------------------------
\chapter{Réservoirs, flux et lois de conservation}
\label{ch:reservoirs}
% ---------------------------------------------------------------------

\section{Les quatre réservoirs du système}

Le système proto-terrestre est décomposé en quatre réservoirs :
\begin{equation}
\mathcal{R}=\left\{\RT,\ \RD,\ \RB,\ \Rinf\right\}.
\label{eq:reservoirs}
\end{equation}

\begin{definition}
\textbf{$\RT$} --- la proto-Terre : sa masse, son spin, sa structure interne et la composition de ses
couches.\\
\textbf{$\RD$} --- le réservoir circumterrestre secondaire : disque gazeux, multiphasé ou particulaire,
et population d'agrégats ou de lunules qu'il contient.\\
\textbf{$\RB$} --- le corps planétaire massif $B$ : sa masse, son spin et son mouvement orbital
relatif à la proto-Terre.\\
\textbf{$\Rinf$} --- le milieu extérieur non lié au système Terre--disque : matière héliocentrique
qui entre dans le système (co-accrétion) ou qui en sort (échappement).
\end{definition}

Chaque réservoir $i$ est décrit par un état
\begin{equation}
\mathcal{X}_i=\left\{M_i,\ \mathbf{J}_i,\ E_i,\ \mathbf{X}_i\right\},
\label{eq:etat-reservoir}
\end{equation}
qui rassemble sa masse, son moment cinétique, son énergie et sa composition. Les moments
cinétiques sont définis dans le repère barycentrique du système :
\begin{equation}
\mathbf{J}_\oplus=\mathbf{S}_\oplus,\qquad
\mathbf{J}_D=\mathbf{L}_D+\sum_k\mathbf{L}_k,\qquad
\mathbf{J}_B=\mathbf{S}_B+\mathbf{L}_{\mathrm{orb}},\qquad
\mathbf{J}_\infty=\mathbf{L}_{\mathrm{esc}},
\label{eq:J-reservoirs}
\end{equation}
où $\mathbf{L}_k$ est le moment orbital des lunules et $\mathbf{L}_{\mathrm{orb}}$ celui du
mouvement relatif Terre--$B$.

\begin{figure}[htbp]
\centering
\begin{tikzpicture}[
  rs/.style={draw,thick,rounded corners=4pt,align=center,font=\small,minimum width=2.7cm,
             minimum height=1.35cm,fill=#1},rs/.default=white]
  \node[rs=orange!12] (T) at (0,0) {$\RT$\\proto-Terre};
  \node[rs=blue!10] (D) at (6.3,0) {$\RD$\\réservoir\\circumterrestre};
  \node[rs=gray!15] (B) at (11.7,0) {$\RB$\\corps $B$};
  \node[rs=white] (I) at (6.3,3.4) {$\Rinf$\\milieu extérieur};
  \draw[-{Stealth},thick] ([yshift=5pt]T.east) -- node[above,font=\scriptsize,align=center]{$\Phi_{\oplus\to D}$\\perte équatoriale\\et injection} ([yshift=5pt]D.west);
  \draw[-{Stealth},thick] ([yshift=-5pt]D.west) -- node[below,font=\scriptsize]{$\Phi_{D\to\oplus}$ réaccrétion} ([yshift=-5pt]T.east);
  \draw[-{Stealth},thick] ([xshift=-6pt]I.south) -- node[left,font=\scriptsize,align=right]{$\Phi_{\infty\to D}$\\co-accrétion} ([xshift=-6pt]D.north);
  \draw[-{Stealth},thick] ([xshift=6pt]D.north) -- node[right,font=\scriptsize]{$\Phi_{D\to\infty}$ échappement} ([xshift=6pt]I.south);
  \draw[-{Stealth},thick] (I.west) to[bend right=20] node[above left,font=\scriptsize]{$\Phi_{\infty\to\oplus}$ accrétion} (T.north);
  \draw[-{Stealth},gray!60!black] ([yshift=5pt]D.east) -- node[above,font=\scriptsize]{$\Phi_{D\to B}$} ([yshift=5pt]B.west);
  \draw[-{Stealth},gray!60!black] ([yshift=-5pt]B.west) -- node[below,font=\scriptsize]{$\Phi_{B\to D}$} ([yshift=-5pt]D.east);
  \draw[red!60!black,thick,decorate,decoration={snake,amplitude=1.2pt,segment length=6pt},-{Stealth}]
     (B.south) to[bend left=28] node[below,font=\scriptsize]{couples $\mathcal{T}_{B\to D}$, $\mathcal{T}_{B\to\oplus}$} (T.south);
\end{tikzpicture}
\caption{Les quatre réservoirs et leurs échanges. Les flèches pleines représentent des flux de
masse, porteurs de moment cinétique, d'énergie et de composition ; la flèche ondulée représente
les couples gravitationnels exercés par $B$ pendant son passage.}
\label{fig:reservoirs}
\end{figure}

\section{Flux de masse}

On note $\Phi_{i\to j}\ge0$ le flux de masse du réservoir $i$ vers le réservoir $j$. L'évolution de
chaque réservoir s'écrit
\begin{equation}
\frac{\dd M_i}{\dd t}=\sum_{j\neq i}\left(\Phi_{j\to i}-\Phi_{i\to j}\right),
\qquad
\sum_i M_i=\text{constante}.
\label{eq:flux-masse}
\end{equation}
Les flux principaux sont regroupés dans le tableau~\ref{tab:flux}. L'équation~\eqref{eq:bilanM} du
disque est le cas particulier $i=D$ de cette relation.

\begin{table}[htbp]
\centering
\small
\caption{Flux de masse principaux entre réservoirs et processus associés.}
\label{tab:flux}
\begin{tabular}{lll}
\toprule
\textbf{Flux} & \textbf{Processus} & \textbf{Phase dominante} \\
\midrule
$\Phi_{\infty\to\oplus}$ & croissance de la proto-Terre & croissance \\
$\Phi_{\infty\to D}$ & co-accrétion, capture dans la sphère de Hill & croissance \\
$\Phi_{\oplus\to D}^{\mathrm{rot}}$ & perte de masse équatoriale rotationnelle & croissance \\
$\Phi_{\oplus\to D}^{\mathrm{mar}}$ & injection par la marée de $B$ & rencontre \\
$\Phi_{D\to\oplus}$ & réaccrétion & rencontre, relaxation \\
$\Phi_{D\to\infty}$ & échappement & rencontre, relaxation \\
$\Phi_{D\to B},\ \Phi_{B\to D}$ & capture par $B$, cession de matière par $B$ & rencontre \\
\bottomrule
\end{tabular}
\end{table}

\section{Moment cinétique}

Un flux de masse transporte le moment cinétique spécifique $\mathbf{j}_{i\to j}$ de la matière
transférée, $\boldsymbol{\Lambda}_{i\to j}=\Phi_{i\to j}\,\mathbf{j}_{i\to j}$. Les réservoirs échangent
en outre du moment cinétique par les couples gravitationnels $\boldsymbol{\mathcal{T}}_{j\to i}$ :
\begin{equation}
\frac{\dd\mathbf{J}_i}{\dd t}=\sum_{j\neq i}\left(\boldsymbol{\Lambda}_{j\to i}-\boldsymbol{\Lambda}_{i\to j}\right)
+\sum_{j\neq i}\boldsymbol{\mathcal{T}}_{j\to i},
\qquad
\boldsymbol{\mathcal{T}}_{j\to i}=-\boldsymbol{\mathcal{T}}_{i\to j}.
\label{eq:flux-J}
\end{equation}
Les couples internes se compensent deux à deux, de sorte que
\begin{equation}
\boxed{\;\sum_i\mathbf{J}_i=\text{constante}\;}
\label{eq:conservation-J}
\end{equation}
tant que le couple solaire est négligeable, ce qui est le cas pendant la rencontre. Le bilan de
moment cinétique de la rencontre (chapitre~\ref{ch:bilans}) est l'application directe de cette loi.

\begin{ideephys}
Dans ce langage, le rôle de $B$ devient exact. Pendant son passage, les couples
$\boldsymbol{\mathcal{T}}_{B\to D}$ et $\boldsymbol{\mathcal{T}}_{B\to\oplus}$ dominent tous les autres
termes, et ils déclenchent les flux $\Phi_{\oplus\to D}$, $\Phi_{D\to\oplus}$ et $\Phi_{D\to\infty}$. Lorsque
$B$ s'éloigne, ces termes s'annulent : c'est le sens du mot \emph{transitoire} dans la définition
de $B$.
\end{ideephys}

\section{Énergie}

L'énergie mécanique n'est pas conservée au cours du chemin : la relaxation du réservoir
convertit de l'énergie orbitale en chaleur, qui est ensuite rayonnée. Le bilan global s'écrit
\begin{equation}
\frac{\dd}{\dd t}\left(\sum_iE_i+E_{\mathrm{int}}+E_{\mathrm{th}}\right)=-\mathcal{L}_{\mathrm{rad}},
\label{eq:energie}
\end{equation}
où $E_{\mathrm{int}}$ est l'énergie d'interaction gravitationnelle entre réservoirs, $E_{\mathrm{th}}$
l'énergie thermique et $\mathcal{L}_{\mathrm{rad}}$ la puissance rayonnée. La dissipation par chocs,
collisions et viscosité, détaillée au chapitre~\ref{ch:reorganisation}, est le terme source de
$E_{\mathrm{th}}$.

\section{Composition}

Chaque flux transporte une composition $\mathbf{X}_{i\to j}$ qui n'est pas nécessairement celle du
réservoir moyen : elle dépend de la profondeur, de la latitude et de la phase de la matière
transférée. La composition de chaque réservoir évolue selon
\begin{equation}
\frac{\dd(M_i\mathbf{X}_i)}{\dd t}=\sum_{j\neq i}\left(\Phi_{j\to i}\mathbf{X}_{j\to i}
-\Phi_{i\to j}\mathbf{X}_{i\to j}\right)+\boldsymbol{\Pi}_i,
\label{eq:flux-X}
\end{equation}
où $\boldsymbol{\Pi}_i$ représente les transformations internes au réservoir : fusion, condensation,
évaporation, ségrégation du métal.

\begin{lecture}
L'équation~\eqref{eq:flux-X} est la porte d'entrée de la géochimie. Le flux
$\Phi_{\oplus\to D}^{\mathrm{mar}}$ prélève la matière d'une couche précise de la proto-Terre ; sa
composition $\mathbf{X}_{\oplus\to D}$ porte donc la mémoire de la latitude et de la profondeur
arrachées. La partie~\ref{part:geochimie} exploite cette mémoire.
\end{lecture}

\section{Variables d'état du Niveau I}

Le Niveau I caractérise l'état admissible du système Terre--disque juste avant le passage de $B$,
puis mesure sa transformation pendant la rencontre. Les systèmes satellitaires naturels fixent
des classes de comportement à reproduire ou à distinguer : transport de moment cinétique,
refroidissement, franchissement de la limite de Roche, migration, accrétion et concentration de
la masse satellitaire.

L'état proto-terrestre est
\begin{equation}
\Psi_\oplus=\left\{\Mt,\ \Req,\ \Omega_\oplus,\ S_\oplus,\ T_\oplus(r),\ \rho_\oplus(r),\
\mathbf{X}_\oplus(r)\right\},
\label{eq:Psiterre}
\end{equation}
et l'état du disque
\begin{equation}
\begin{aligned}
\mathcal{D}_0=\Big\{&M_{D,0},\ \Sigma_0(r),\ R_{\mathrm{in}},\ R_{\mathrm{out}},\ T_0(r),\
P_0(r),\ H_0(r),\\
&L_{D,0},\ v_{r,0}(r),\ \kappa_0(r),\ \mathbf{X}_0(r),\
\mathbf{f}_{\mathrm{phase},0}(r),\ \mathcal{G}_0\Big\},
\end{aligned}
\label{eq:D0}
\end{equation}
où $\mathcal{G}_0=\{m_k,a_k,e_k,i_k,\mathbf{X}_k\}_{k=1}^{N_0}$ décrit la population
d'agrégats ou de lunules déjà formés à $t_B$ (section~\ref{sec:regimes}). Le vecteur des fractions
de phase est
\begin{equation}
\mathbf{f}_{\mathrm{phase}}=(f_{\mathrm{vap}},f_{\mathrm{liq}},f_{\mathrm{sol}}),\qquad
f_{\mathrm{vap}}+f_{\mathrm{liq}}+f_{\mathrm{sol}}=1.
\label{eq:fphase}
\end{equation}

\section{Paramètres de la rencontre}

L'état de $B$ est décrit par
\begin{equation}
\Psi_B=\left\{M_B,\ R_B,\ \rho_B,\ q_B,\ v_\infty,\ i_B,\ \omega_B,\ \varpi_B,\ \mathbf{S}_B\right\},
\label{eq:PsiB}
\end{equation}
où $q_B$ est la distance géocentrique minimale de la trajectoire nominale, $v_\infty$ la vitesse
asymptotique relative, $i_B$ l'inclinaison de la trajectoire sur le plan du disque, $\omega_B$
l'argument du périastre compté depuis la ligne des nœuds, $\varpi_B$ la longitude du périastre et
$\mathbf{S}_B$ le spin propre de $B$. Le vecteur minimal de paramètres est
\begin{equation}
\Theta_I=\left\{\Psi_\oplus,\ \mathcal{D}_0,\ \Psi_B,\ t_B\right\}.
\label{eq:Theta}
\end{equation}
Pour relier la rencontre au destin ultérieur de $B$, on définit aussi l'extension héliocentrique
\begin{equation}
\Theta_I^{+}=\left\{\Theta_I,\ \mathcal{H}_{B,1},\ t_{\mathrm{ej}}\right\},
\qquad
\mathcal{H}_{B,1}=\left\{a_{B,1},e_{B,1},i_{B,1},\Omega_{B,1},\omega_{B,1}\right\},
\label{eq:Theta-plus}
\end{equation}
où $\mathcal{H}_{B,1}$ désigne les éléments héliocentriques de $B$ après son passage près de la
proto-Terre. Cette extension n'est pas nécessaire pour calculer le tri orbital local, mais elle devient
nécessaire pour tester le raccord avec l'instabilité globale du Système solaire.

\begin{falsif}
Le Niveau I échoue si aucun domaine continu de $\Theta_I$ ne permet à la fois un disque
préexistant physiquement persistant, une rencontre non collisionnelle, un bilan de moment
cinétique fermé, une masse de disque post-rencontre suffisante et une composition compatible
avec les contraintes lunaires.
\end{falsif}


\section{Échelles de référence du système proto-Terre--disque--\texorpdfstring{$B$}{B}}

Le tableau~\ref{tab:echelles} rassemble les échelles numériques qui encadrent le Niveau I. Elles
sont calculées pour une proto-Terre de masse et de rayon terrestres actuels à 1~UA ; les valeurs
associées à $B$ sont des \emph{valeurs d'échelle}, établies pour $M_B=0{,}1\,\Mt$, $q_B=2\,\Rt$ et
une trajectoire parabolique. Elles servent à fixer des ordres de grandeur et ne doivent pas être
confondues avec la configuration de référence de la théorie, définie à la
section~\ref{sec:equation-lune} par $\tilde\omega=0{,}97$, $\Req=1{,}70\,\Rt$, $\mu_B=0{,}95$ et
$q_B=3{,}0\,\Rt$.

\begin{table}[htbp]
\centering
\small
\caption{Échelles de référence. Ces valeurs ne sont pas des solutions : elles fixent les ordres de
grandeur par rapport auxquels les sorties du calcul doivent être situées.}
\label{tab:echelles}
\begin{tabular}{lll}
\toprule
\textbf{Grandeur} & \textbf{Expression} & \textbf{Valeur} \\
\midrule
Rayon de Hill & $a_\oplus(\Mt/3M_\odot)^{1/3}$ & $1{,}50\times10^{9}$ m $\approx235\,\Rt$ \\
Bord externe tronqué & $\xi_H R_H$, $\xi_H\approx0{,}3$--$0{,}4$ & $\approx70$--$95\,\Rt$ \\
Masse thermique & $(H/r)^3M_\odot$ & $\approx10\,\Mt$ \\
Rayon de Bondi & $G\Mt/c_s^2$ & $\approx63\,\Rt$ \\
Circularisation (Hill) & $R_H/48$ & $\approx4{,}9\,\Rt$ \\
Circularisation (Bondi) & $(R_{\mathrm{B}}/R_H)^4R_H/48$ & $\approx0{,}03\,\Rt$ \\
Limite de Roche ($\rho_s=3{,}34$ g\,cm$^{-3}$) & éq.~\eqref{eq:Roche} & $1{,}85\times10^{7}$ m $\approx2{,}9\,\Rt$ \\
Période orbitale à $\aR$ & $2\pi(\aR^3/G\Mt)^{1/2}$ & $\approx7{,}0$ h \\
Orbite synchrone ($P_\oplus=4$--$5$ h) & $(G\Mt/\Omega_\oplus^2)^{1/3}$ & $\approx2{,}0$--$2{,}3\,\Rt$ \\
Période critique de rotation & $2\pi(\Rt^3/G\Mt)^{1/2}$ & $\approx1{,}4$ h \\
Spin critique ($k\approx0{,}33$) & $k\Mt\Rt^2\Omega_{\mathrm{c}}$ & $\approx1{,}0\times10^{35}$ kg\,m$^2$\,s$^{-1}$ \\
Moment cinétique Terre--Lune actuel & $L_{\mathrm{EM}}$ & $\approx3{,}5\times10^{34}$ kg\,m$^2$\,s$^{-1}$ \\
Moment orbital d'une Lune à $\aR$ & $M_L(G\Mt\aR)^{1/2}$ & $\approx6{,}3\times10^{33}$ kg\,m$^2$\,s$^{-1}$ \\
Moment orbital relatif de $B$ & éq.~\eqref{eq:LBrel} & $\approx5{,}7\times10^{34}$ kg\,m$^2$\,s$^{-1}$ \\
Vitesse au périastre de $B$ & $(2G\Mtot/q_B)^{1/2}$ & $\approx8{,}3$ km\,s$^{-1}$ \\
Durée caractéristique de la rencontre & $q_B/v_p$ & $\approx26$ min \\
Distance de Roche de $B$ ($\rho_B=3{,}93$ g\,cm$^{-3}$) & éq.~\eqref{eq:RocheB} & $\approx2{,}75\,\Rt$ \\
\bottomrule
\end{tabular}
\end{table}

\begin{controle}
Trois comparaisons structurent la suite. Le moment orbital relatif porté par $B$ est du même
ordre que le moment cinétique total du système Terre--Lune : un passage rasant dispose, sans
collision, d'un réservoir de moment cinétique à l'échelle du problème. La durée de la rencontre
est courte devant la période orbitale à la limite de Roche : le disque subit un forçage
impulsif. Enfin, la distance de Roche de $B$ est voisine des distances de passage envisagées :
l'intégrité de $B$ devient une contrainte explicite.
\end{controle}
% =====================================================================
\part{Chronologie physique}
\label{part:chronologie}
% =====================================================================

\noindent
La chronologie n'est pas racontée ici comme une suite d'épisodes. Elle est construite à partir de
la physique : chaque phase du chemin est définie par les termes qui dominent les équations de
flux~\eqref{eq:flux-masse}--\eqref{eq:flux-X}, et elle est caractérisée par son échelle de temps
propre. La séparation de ces échelles justifie le découpage du calcul en phases successives.

% ---------------------------------------------------------------------
\chapter{Chronologie physique du chemin}
\label{ch:chronologie}
% ---------------------------------------------------------------------

\section{Les phases définies par les flux dominants}

Le tableau~\ref{tab:phases} définit les phases du chemin par leurs termes dominants.

\begin{table}[htbp]
\centering
\small
\caption{Phases du chemin physique, définies par les termes dominants des équations de flux.}
\label{tab:phases}
\begin{tabularx}{\textwidth}{>{\raggedright\arraybackslash}p{3.1cm}
>{\raggedright\arraybackslash}X >{\raggedright\arraybackslash}p{3.6cm}}
\toprule
\textbf{Phase} & \textbf{Termes dominants} & \textbf{Échelle de temps} \\
\midrule
0 --- Croissance et alimentation & $\Phi_{\infty\to\oplus}$, $\Phi_{\infty\to D}$, $\Phi^{\mathrm{rot}}_{\oplus\to D}$, couples $\mathcal{T}_{\oplus\to D}$ & $10^{6}$--$10^{8}$ ans \\
I --- État pré-rencontre & transport visqueux, accrétion en agrégats et lunules & $t_\nu$, $t_{\mathrm{sat}}$ \\
II --- Rencontre & $\mathcal{T}_{B\to D}$, $\mathcal{T}_{B\to\oplus}$, $\Phi^{\mathrm{mar}}_{\oplus\to D}$ & $t_{p}\sim$ 0,5--2 h \\
III --- Tri et relaxation & $\Phi_{D\to\oplus}$, $\Phi_{D\to\infty}$, dissipation & $10$--$10^{3}$ périodes orbitales \\
IV --- Accrétion lunaire & accrétion au-delà de Roche, transport depuis le disque intérieur & mois à $10^{2}$ ans \\
V --- Océan magmatique lunaire & refroidissement et différenciation de la Lune & $10^{7}$--$10^{8}$ ans \\
VI --- Destin de $B$ & diffusion héliocentrique, interactions planétaires, élimination ou éjection & instabilité précoce ou tardive \\
\bottomrule
\end{tabularx}
\end{table}

\section{Hiérarchie des échelles de temps}

Les échelles de temps caractéristiques se classent selon
\begin{equation}
\boxed{\;t_{p}<P_{\mathrm{orb}}(\aR)\ll t_{\mathrm{relax}}\lesssim t_{\mathrm{acc},L}
\ll t_{\mathrm{LMO}}\;}
\label{eq:hierarchie-temps}
\end{equation}
où $t_{p}$ désigne l'échelle périastrique $q_B/v_p$, qui vaut environ $26$~min pour les
valeurs d'échelle du tableau~\ref{tab:echelles} et environ $0{,}95$~h pour la configuration de
référence. Cette échelle mesure la traversée du périastre et non la durée du forçage : la durée
pendant laquelle le champ différentiel reste comparable à sa valeur maximale atteint $5$ à $12$~h
selon la distance retenue, comme le détaille le paragraphe~\ref{par:duree}. Puisque
$P_{\mathrm{orb}}(\aR)\approx7$~h, la hiérarchie ci-dessus vaut pour l'échelle périastrique, mais
l'hypothèse d'un réservoir figé pendant tout le forçage cesse d'être uniformément valable. C'est
la raison pour laquelle le calcul de Phase~B intègre la trajectoire hyperbolique complète au lieu
d'appliquer une impulsion. La figure~\ref{fig:temps} les place sur une échelle logarithmique.

\begin{figure}[htbp]
\centering
\begin{tikzpicture}[x=1cm,y=1cm]
  \draw[thick,-{Stealth}] (0,0) -- (14.4,0) node[right,font=\scriptsize] {$t$ (s)};
  \foreach \k in {2,...,16}{
    \pgfmathsetmacro{\xx}{\k-2}
    \draw (\xx,0.08) -- (\xx,-0.08) node[below,font=\scriptsize] {$10^{\k}$};}
  \fill[violet!30] (1.176,0.25) rectangle (1.845,0.55);
  \node[font=\scriptsize,violet!60!black,align=center] at (1.477,1.05) {rencontre\\$t_{p}$};
  \fill[black] (2.398,0.4) circle (2pt);
  \node[font=\scriptsize,align=center] at (2.398,1.75) {$P_{\mathrm{orb}}(\aR)$};
  \draw[thin] (2.398,0.45) -- (2.398,1.55);
  \fill[teal!30] (3.398,0.25) rectangle (5.398,0.55);
  \node[font=\scriptsize,teal!50!black,align=center] at (4.398,1.05) {tri et\\relaxation};
  \fill[orange!35] (4.415,-0.95) rectangle (7.477,-0.65);
  \node[font=\scriptsize,orange!60!black] at (5.954,-1.25) {accrétion lunaire};
  \fill[gray!40] (12.477,0.25) rectangle (13.477,0.55);
  \node[font=\scriptsize,align=center] at (13.000,1.05) {océan magmatique\\lunaire};
  \fill[blue!30] (11.477,-0.95) rectangle (12.477,-0.65);
  \node[font=\scriptsize,blue!50!black,align=center] at (12.000,-1.4) {dissipation du\\gaz nébulaire};
\end{tikzpicture}
\caption{Échelles de temps du chemin physique (valeurs d'ordre de grandeur). Onze ordres de
grandeur séparent la durée de la rencontre de celle de l'océan magmatique lunaire.}
\label{fig:temps}
\end{figure}

\begin{ideephys}
La rencontre dure moins d'une période orbitale à la limite de Roche : pendant le passage, le
réservoir n'a pas le temps d'évoluer visqueusement, et il peut être traité comme figé. Après le
passage, la relaxation s'étend sur des centaines de périodes, sans que $B$ intervienne encore.
Cette séparation des échelles justifie le découpage du programme numérique en Phases~A, B et
C (chapitre~\ref{ch:programme}).
\end{ideephys}

\section{Ordre des étapes}

La chronologie minimale est
\begin{equation}
t^{\mathrm{start}}_{\mathrm{accr},\oplus}<t^{\mathrm{form}}_D<t_B<t^{\mathrm{reorg}}_D
<t^{\mathrm{accr}}_L<t_{\mathrm{LMO}},
\label{eq:ordre}
\end{equation}
Cette chaîne fixe l'ordre interne du problème Terre--Lune. Le raccord héliocentrique ajoute une
condition indépendante :
\begin{equation}
\boxed{\;t_B<t_{\mathrm{ej}}\;},
\label{eq:ordreB}
\end{equation}
sans imposer que $t_{\mathrm{ej}}$ appartienne à une instabilité précoce ou tardive. Le calcul doit
simplement montrer qu'une orbite post-passage de $B$ peut rejoindre une trajectoire éliminée par
la dynamique globale.

et un vecteur chronologique plus complet est
\begin{equation}
\mathbf{t}=\left\{t^{\mathrm{form}}_D,\ t^{\mathrm{feed}}_D,\ t^{\mathrm{cool}}_D,\
t^{\mathrm{cond}}_D,\ t_B,\ t_{\mathrm{relax}},\ t_{\mathrm{Roche}},\
t^{\mathrm{start}}_{\mathrm{acc},L},\ t^{\mathrm{end}}_{\mathrm{acc},L},\ t_{\mathrm{LMO}},\ t_{\mathrm{ej}}\right\}.
\label{eq:vecteur-t}
\end{equation}

\section{Disque gazeux et réservoir résiduel}

La théorie distingue les états physiques du réservoir circumterrestre :
\begin{equation}
\mathcal{D}(t)=\mathcal{D}_{\mathrm{gaz}}\cup\mathcal{D}_{\mathrm{multiphase}}\cup
\mathcal{D}_{\mathrm{particulaire}},
\label{eq:etats}
\end{equation}
la pondération entre ces régimes étant calculée par $\mathbf{f}_{\mathrm{phase}}(r,t)$.

Le gaz nébulaire se dissipe en quelques millions d'années, alors que les chronomètres lunaires
placent l'accrétion de la Lune plusieurs dizaines de millions d'années après la formation des
premiers solides ; la croissance rapide de la proto-Terre par accrétion de galets
\cite{Johansen2021} et son achèvement ultérieur encadrent l'intervalle
$[t^{\mathrm{start}}_{\mathrm{accr},\oplus},t_B]$. À l'instant du passage, le réservoir est donc
décrit par ses composantes multiphasée et particulaire, alimentées par les canaux de
l'équation~\eqref{eq:canaux}.

\begin{lecture}
Cette distinction évite de faire dépendre la théorie de la persistance d'un disque riche en gaz.
Le réservoir évolue vers un état multiphasé ou particulaire tout en conservant masse et moment
cinétique circumterrestres, et c'est cet état que le passage de $B$ réorganise.
\end{lecture}

% =====================================================================

\section{Trois régimes chronologiques}
\label{sec:regimes}

Soit $t_{\mathrm{sat}}$ le temps caractéristique d'accrétion satellitaire et
$\Delta t_{\mathrm{preB}}=t_B-t_{D,\mathrm{form}}$ le temps disponible avant le passage. On
distingue
\begin{align}
&\text{régime d'accrétion tardive :} && t_{\mathrm{sat}}\gg\Delta t_{\mathrm{preB}},
\label{eq:tardif}\\
&\text{régime intermédiaire :} && t_{\mathrm{sat}}\sim\Delta t_{\mathrm{preB}},
\label{eq:interm}\\
&\text{régime d'accrétion précoce :} && t_{\mathrm{sat}}\ll\Delta t_{\mathrm{preB}}.
\label{eq:precoce}
\end{align}

Les échelles de temps disponibles orientent ce choix. Une composante gazeuse non alimentée
s'étale en
\begin{equation}
t_\nu=\left[\alpha\,h^2\,\Omega_K\right]^{-1}\approx80\ \text{ans}
\qquad(\alpha=10^{-3},\ h=H/r=0{,}1,\ r=10\,\Rt),
\label{eq:tnu-num}
\end{equation}
Cette valeur est évaluée à $10\,\Rt$. Il est utile de la reprendre au rayon où se trouve
effectivement le réservoir, puisque $\Omega_K\propto r^{-3/2}$ :
\begin{equation}
t_\nu\simeq7\ \text{ans à }2\,\Rt,\qquad t_\nu\simeq13\ \text{ans à }2{,}9\,\Rt
\qquad(\alpha=10^{-3},\ h=0{,}1),
\label{eq:tnu-reservoir}
\end{equation}
et de noter que pour un réservoir massif, dont le transport est dominé par les couples
gravitationnels, le $\alpha$ effectif monte à $10^{-2}$--$10^{-1}$ et ramène ces durées à quelques
mois. La fenêtre temporelle offerte au passage est donc une contrainte à part entière, et c'est
l'équation~\eqref{eq:tnu-reservoir} qui doit servir de référence à la
porte~\eqref{eq:tpers}.
Une composante particulaire située au-delà de la limite de Roche s'agrège quant à elle en quelques
mois à environ un an \cite{Ida1997,Kokubo2000}. Ces temps de référence sont à comparer à
$\Delta t_{\mathrm{preB}}$ et au taux d'alimentation du disque ; le calcul détermine lequel des trois
régimes appartient au domaine viable. Lorsque le réservoir comprend des lunules
$\mathcal{G}_0$ à l'instant du passage, celles qui sont situées au-delà de $a_{\mathrm{sync}}$ migrent
vers l'extérieur.

\begin{lecture}
Lorsque des lunules sont présentes, le rôle de $B$ se précise : en excitant excentricités et inclinaisons, il provoque
des croisements d'orbites entre lunules et entre lunules et disque résiduel, donc des fusions.
Le passage rasant devient un mécanisme de concentration de la masse satellitaire, qui donne un
contenu physique à la condition $f_{\mathrm{dom}}\to1$.
\end{lecture}


\begin{rappel}
La chronologie a fixé l'ordre et la durée des phases. La partie suivante entre dans la phase la plus
brève et la plus décisive : la rencontre, pendant laquelle l'agent $B$ redistribue les flux de masse
et de moment cinétique.
\end{rappel}

% =====================================================================
\part{Dynamique : l'agent gravitationnel \texorpdfstring{$B$}{B}}
\label{part:dynamique}
% =====================================================================

\noindent
Cette partie constitue le c\oe ur original de la théorie. Elle décrit comment un corps planétaire
massif $B$, sans jamais toucher la proto-Terre, redistribue les flux de masse et de moment
cinétique entre les réservoirs. Quatre chapitres la composent : le passage lui-même et sa
géométrie dynamique de basse latitude, l'injection de matière superficielle proto-terrestre, le
résultat central du passage --- le tri orbital qui définit la fenêtre lunaire --- et enfin les
bilans qui ferment la rencontre.

% ---------------------------------------------------------------------
\chapter{Le passage non collisionnel}
\label{ch:passage}\label{sec:passage}
% ---------------------------------------------------------------------

Conformément à sa définition (chapitre~\ref{ch:probleme}), $B$ agit sur le système
Terre--réservoir par son seul champ gravitationnel. Ce chapitre établit la fréquence de tels
passages, leur géométrie, leur régime dynamique et le champ de forces exact qu'ils exercent.

\section{Une interaction gravitationnelle sans contact}

Pendant la phase de croissance des planètes telluriques, les embryons planétaires se croisent
et subissent de nombreuses rencontres rapprochées. La section efficace des rencontres dont le
périastre est inférieur à $q$ s'écrit, focalisation gravitationnelle comprise,
\begin{equation}
\sigma(q)=\pi q^2\left[1+\frac{2G\Mtot}{q\,v_\infty^2}\right],
\qquad \Mtot=\Mt+M_B.
\label{eq:sigma}
\end{equation}
Soit $R_c=\Req+R_{B,\mathrm{eff}}$ la distance de contact et $\Theta=2G\Mtot/(R_c v_\infty^2)$ le
paramètre de focalisation. Le rapport entre le nombre de passages rasants
($R_c<q<2R_c$) et le nombre de collisions ($q<R_c$) vaut
\begin{equation}
\frac{\mathcal{N}_{\mathrm{ras}}}{\mathcal{N}_{\mathrm{coll}}}
=\frac{\sigma(2R_c)-\sigma(R_c)}{\sigma(R_c)}
=\frac{3+\Theta}{1+\Theta}\ \in\ [1,\,3].
\label{eq:frequence}
\end{equation}

\begin{apport}
Les passages rasants sont au moins aussi fréquents que les collisions, et jusqu'à trois fois
plus fréquents lorsque la focalisation gravitationnelle est faible. Le mécanisme de la théorie
ne repose donc pas sur un événement exceptionnel : il retient, dans la population naturelle des
rencontres planétaires, celles qui ne se terminent pas par un contact. Les rencontres non
collisionnelles sont par ailleurs reconnues comme un processus efficace de la dynamique
Terre--Lune primitive \cite{PahlevanMorbidelli2015} ; la théorie leur attribue ici un rôle
antérieur et plus fondamental : réorganiser le disque d'où naît la Lune.
\end{apport}

\section{Définition géométrique du caractère non collisionnel}

Le scénario est gravitationnel. La trajectoire de $B$ ne touche pas la proto-Terre :
\begin{equation}
q_B>\Req+R_{B,\mathrm{eff}},
\label{eq:nc}
\end{equation}
où $\Req$ est le rayon équatorial de la proto-Terre en rotation, supérieur au rayon polaire, et
$R_{B,\mathrm{eff}}$ le rayon de $B$ déformé par la marée terrestre. L'excentricité de la
trajectoire hyperbolique relative est
\begin{equation}
e_B=1+\frac{q_B\,v_\infty^2}{G\Mtot}.
\label{eq:eB}
\end{equation}
La trajectoire traverse le plan du disque en deux nœuds situés à
\begin{equation}
r_{\text{nœud}}^{\pm}=\frac{q_B\,(1+e_B)}{1\pm e_B\cos\omega_B}.
\label{eq:noeuds}
\end{equation}
Pour les configurations où la théorie exige que $B$ ne traverse pas la partie dense du disque,
on impose
\begin{equation}
R_{\mathrm{out,dense}}<\min\left(r_{\text{nœud}}^{+},r_{\text{nœud}}^{-}\right)
=\frac{q_B\,(1+e_B)}{1+e_B\,|\cos\omega_B|}.
\label{eq:nontraversee}
\end{equation}
Pour une trajectoire quasi parabolique dont le périastre est au plus haut au-dessus du plan du
disque ($\omega_B=\pi/2$), cette borne vaut $2q_B$.

\begin{lecture}
L'action de $B$ est portée par son champ gravitationnel. La géométrie inclinée permet une forte
perturbation du disque tout en maintenant la séparation physique entre $B$, la proto-Terre et
la partie dense du disque : un disque dense de quelques rayons terrestres, cohérent avec le
rayon de circularisation de l'équation~\eqref{eq:rc}, est entièrement compatible avec un passage
à $q_B$ de quelques $\Rt$.
\end{lecture}

\section{Régime dynamique de la rencontre}

La vitesse relative au périastre et la durée caractéristique de la rencontre sont
\begin{equation}
v_p=\sqrt{v_\infty^2+\frac{2G\Mtot}{q_B}},
\qquad
t_{p}\simeq\frac{q_B}{v_p}.
\label{eq:tenc}
\end{equation}
Le régime de réponse d'une région du disque de rayon $r$ est mesuré par
\begin{equation}
\eta(r)=\Omega_K(r)\,t_{p}.
\label{eq:eta}
\end{equation}
La notation distingue désormais deux durées qu'il ne faut pas confondre : $t_p=q_B/v_p$ est
l'échelle périastrique, tandis que $t_{\mathrm{eff}}$ désigne la durée pendant laquelle le champ
différentiel reste comparable à son maximum. La seconde vaut cinq à dix fois la première.
Pour $\eta\ll1$ la réponse est impulsive, pour $\eta\sim1$ elle est résonante, pour $\eta\gg1$
elle est adiabatique et le forçage s'atténue. Dans le cas de référence du
tableau~\ref{tab:echelles}, l'échelle périastrique vaut $t_{p}\approx26$~min et
$\eta(\aR)\approx0{,}4$ : au voisinage du périastre, la région où s'assemble la Lune reçoit un
forçage impulsif. Rapporté à la durée effective du forçage, $\eta(\aR)$ dépasse l'unité et le
régime devient résonant, ce qui justifie l'intégration complète plutôt qu'une impulsion. Pour la proto-Terre elle-même, le paramètre
analogue $\eta_\oplus=\Omega_\oplus t_{p}$ est de l'ordre de l'unité lorsque la rotation
est proche de la limite critique : la marée de $B$ agit alors sur une échelle de temps voisine de
la période de rotation, ce qui renforce le couplage entre le passage et la matière équatoriale.

\section{Champ différentiel exact}

Dans un repère centré sur la proto-Terre, l'accélération différentielle exacte due à $B$ sur une
parcelle située en $\mathbf{r}$ est
\begin{equation}
\mathbf{a}_{B,\mathrm{diff}}=-GM_B\left[
\frac{\mathbf{r}-\mathbf{r}_B(t)}{|\mathbf{r}-\mathbf{r}_B(t)|^3}
+\frac{\mathbf{r}_B(t)}{|\mathbf{r}_B(t)|^3}\right].
\label{eq:adiff}
\end{equation}
Le premier terme est l'attraction directe de $B$, le second l'accélération de la proto-Terre
elle-même vers $B$ (terme indirect). Cette expression est utilisée sans développement
multipolaire dans le régime rasant, où $|\mathbf{r}|$ et $|\mathbf{r}_B|$ sont comparables.

\section{Couple exercé sur le disque}

Le couple total de $B$ sur le disque et la variation de moment cinétique qui en résulte sont
\begin{equation}
\boldsymbol{\mathcal{T}}_{B\to D}=\int_D\mathbf{r}\times\mathbf{a}_{B,\mathrm{diff}}\,\dd m,
\qquad
\Delta\mathbf{L}_D=\int_{t_i}^{t_f}\boldsymbol{\mathcal{T}}_{B\to D}\,\dd t+\Delta\mathbf{L}_{\mathrm{flux}},
\label{eq:couple}
\end{equation}
où $\Delta\mathbf{L}_{\mathrm{flux}}$ représente le moment cinétique transporté par la matière
qui entre dans le disque ou le quitte pendant la rencontre. Les sorties dynamiques principales
sont
\begin{equation}
\mathcal{O}_B=\left\{\Delta M_D,\ \Delta\mathbf{L}_D,\ \Delta e_D,\ \Delta i_D,\
\Delta R_{\mathrm{out}},\ \Delta T_D,\ \Delta\mathcal{G}\right\}.
\label{eq:OB}
\end{equation}


\section{Géométrie dynamique de basse latitude du passage}
\label{sec:basse-latitude}

La géométrie retenue par la théorie est celle d'un passage \emph{parallèle à la bande
dynamique de basses latitudes} d'une proto-Terre oblate en rotation rapide, dans le sens de sa rotation. Elle
réunit quatre effets favorables.

\paragraph{La latitude visée.}
Au périastre, le point de la proto-Terre situé à la verticale de $B$ se trouve à la latitude
\begin{equation}
\lambda_B=\arcsin\left(\sin i_B\,\sin\omega_B\right),
\label{eq:lambdaB}
\end{equation}
où $i_B$ est l'inclinaison de la trajectoire sur l'équateur rotationnel et $\omega_B$ l'argument
du périastre. La condition de passage de basse latitude correspond à
\begin{equation}
|\lambda_B|\le\lambda_c,
\label{eq:basse-latitude}
\end{equation}
soit $i_B\lesssim\lambda_c$ pour $\omega_B=90^\circ$. La grandeur $\lambda_c$ définit une bande
dynamique de basses latitudes autour de l'équateur rotationnel instantané de la proto-Terre. Sa
valeur n'est pas imposée par l'obliquité actuelle ; elle est déterminée par le couplage entre
rotation, aplatissement, marée de $B$ et réponse du réservoir.

\paragraph{Le forçage équatorial est préservé.}
L'accélération de marée radiale varie comme $(3\cos^2\psi-1)$, où $\psi$ est l'angle au point
sous $B$. À l'équateur, à la longitude du point visé, $\psi=\lambda_B$ et le forçage relatif vaut
\begin{equation}
F(\lambda_B)=\frac{3\cos^2\lambda_B-1}{2},
\qquad F(20^\circ)\approx0{,}82.
\label{eq:Flambda}
\end{equation}
Une légère inclinaison ne coûte donc que peu de forçage sur l'équateur.

\paragraph{L'aplatissement renforce l'action de \texorpdfstring{$B$}{B}.}
Pour une proto-Terre oblate, la gravité effective équatoriale s'écrit, au premier ordre en $J_2$,
\begin{equation}
g_{\mathrm{eff,eq}}=\frac{G\Mt}{\Req^2}\left(1+\frac32J_2\right)-\Omega_\oplus^2\Req.
\label{eq:geffJ2}
\end{equation}
Le seuil d'injection (chapitre~\ref{ch:injection}) contient le facteur $(\Req/q_B)^3$ : un rayon
équatorial plus grand de 10 à 20~\% augmente le terme de marée de 33 à 73~\%, précisément là
où la gravité effective est la plus faible.

\paragraph{Le sens prograde prolonge le forçage.}
Si $B$ passe dans le sens de rotation de la proto-Terre, une même longitude reste plus longtemps
sous le bourrelet de marée. Avec $\Omega_B=v_p/q_B$ la vitesse angulaire de $B$ au périastre, le
rapport des durées de forçage prograde et rétrograde vaut
\begin{equation}
\frac{\tau_{\mathrm{pro}}}{\tau_{\text{rétro}}}=\frac{\Omega_\oplus+\Omega_B}{\Omega_\oplus-\Omega_B}.
\label{eq:prograde}
\end{equation}
Pour $q_B=3\,\Rt$, $\Omega_B\approx3{,}5\times10^{-4}$~rad\,s$^{-1}$ ; pour une proto-Terre proche de
sa rotation critique, $\Omega_\oplus\approx1{,}2\times10^{-3}$~rad\,s$^{-1}$, et le rapport vaut
environ $1{,}8$. Ce mécanisme est l'analogue planétaire de l'efficacité des rencontres progrades
mise en évidence pour les galaxies \cite{ToomreToomre1972}.

\paragraph{La géométrie reste compatible avec la non-traversée du disque dense.}
Avec $\omega_B=90^\circ$, $B$ passe au périastre à la hauteur $q_B\sin i_B$ au-dessus du plan
équatorial, et ne recoupe ce plan qu'aux nœuds situés à environ $2q_B$
(équation~\eqref{eq:noeuds}). Un réservoir dense s'étendant jusqu'à $5\,\Rt$ n'est pas traversé dès
que $q_B\gtrsim2{,}5\,\Rt$.

\begin{figure}[htbp]
\centering
\begin{tikzpicture}[scale=1.0]
  % disque vu par la tranche
  \fill[blue!25] (-4.6,-0.07) rectangle (-1.35,0.07);
  \fill[blue!25] (1.35,-0.07) rectangle (4.6,0.07);
  % proto-Terre oblate
  \shade[ball color=orange!75!brown] (0,0) ellipse (1.3 and 1.0);
  \draw[dashed,gray!70!black] (-1.3,0) -- (1.3,0);
  \draw[gray!60!black] ({-1.3*cos(asin(0.39))},0.39) -- ({1.3*cos(asin(0.39))},0.39);
  \draw[gray!60!black] ({-1.3*cos(asin(0.39))},-0.39) -- ({1.3*cos(asin(0.39))},-0.39);
  \fill[red!55,opacity=0.45] ({-1.3*cos(asin(0.39))},-0.39) -- ({1.3*cos(asin(0.39))},-0.39)
        -- (1.3,0) -- ({1.3*cos(asin(0.39))},0.39) -- ({-1.3*cos(asin(0.39))},0.39) -- (-1.3,0) -- cycle;
  \node[font=\scriptsize,gray!30!black] at (0,0.62) {$+\lambda_c$};
  \node[font=\scriptsize,gray!30!black] at (0,-0.62) {$-\lambda_c$};
  % trajectoire inclinée de B
  \draw[dashed,thick,gray!60!black,-{Stealth[length=3mm]}] (-2.2,3.3) -- (6.0,0.4);
  \shade[ball color=gray!55] (2.9,1.50) circle (0.33);
  \node[font=\small] at (2.9,2.08) {$B$};
  % angle d'inclinaison
  \draw[gray] (4.49,0.93) -- (6.0,0.93);
  \draw[gray] (5.69,0.93) arc (0:-19.5:1.2);
  \node[font=\scriptsize] at (6.25,0.72) {$i_B$};
  % marée
  \draw[-{Stealth},thick,red!60!black] (1.3,0.05) -- (2.05,0.35);
  \draw[-{Stealth},thick,red!60!black] (-1.3,0.05) -- (-2.05,-0.25);
  \node[font=\scriptsize,red!60!black,align=center] at (-2.2,0.75) {bourrelet\\de marée};
  \node[font=\scriptsize,blue!45!black] at (3.6,-0.35) {réservoir $\RD$};
  \node[font=\small] at (0,-1.35) {proto-Terre oblate};
  \draw[-{Stealth},thin] (0,1.05) arc (95:40:0.55);
  \node[font=\scriptsize] at (0.75,1.35) {$\Omega_\oplus$};
\end{tikzpicture}
\caption{Géométrie de basses latitudes : $B$ passe sur une trajectoire faiblement inclinée sur
l'équateur rotationnel d'une proto-Terre oblate, dans le sens de sa rotation. Le bourrelet de
marée se forme dans la bande dynamique de basses latitudes, où la gravité effective est minimale ; le réservoir
dense reste en deçà des nœuds de la trajectoire (échelle non respectée).}
\label{fig:basse-latitude}
\end{figure}

\paragraph{Effet sur le réservoir.}
Les lois d'échelle établies pour des rencontres entre un disque de particules et un perturbateur
sur trajectoire parabolique coplanaire et prograde indiquent un rayon de troncature
\begin{equation}
r_{\mathrm{tr}}\approx0{,}28\,q_B\left(\frac{M_B}{\Mt}\right)^{-0{,}32}
\label{eq:troncature}
\end{equation}
\cite{Breslau2014}, soit $r_{\mathrm{tr}}\approx0{,}6\,q_B$ pour $M_B=0{,}1\,\Mt$. La bande dynamique de basses latitudes prograde est donc la plus active du domaine des paramètres : elle maximise à la
fois l'injection et le tri orbital. La fraction échappée $\eta_{\mathrm{esc}}$
(chapitre~\ref{ch:fenetre}) est la variable de contrôle qui délimite, dans cette bande, les
passages favorables.

\begin{ideephys}
La couche arrachée est la plus externe, donc la moins dense : du silicate déjà appauvri en
métal Fe--Ni et en éléments sidérophiles par la formation du noyau. Cela n'implique pas un
appauvrissement en FeO silicaté, qui reste une contrainte propre de la composition lunaire.
\end{ideephys}

\begin{consequence}
La matière injectée porte la signature de la couche silicatée externe des basses latitudes :
appauvrissement en métal, en Ni, Co, W et V hérité de la formation du noyau terrestre, et
enrichissement relatif en éléments incompatibles si la couche est différenciée. Les traceurs de la
partie~\ref{part:geochimie} testent directement cette signature.
\end{consequence}

\section{Réservoir de moment cinétique porté par \texorpdfstring{$B$}{B}}

Dans le repère barycentrique du couple Terre--$B$, le moment cinétique orbital relatif de la
rencontre est
\begin{equation}
L_{B,\mathrm{rel}}=\mu_r\sqrt{G\Mtot\,q_B\,(1+e_B)},
\qquad \mu_r=\frac{\Mt M_B}{\Mtot}.
\label{eq:LBrel}
\end{equation}
Pour $M_B=0{,}1\,\Mt$, $q_B=2\,\Rt$ et $e_B=1$, on obtient
$L_{B,\mathrm{rel}}\approx5{,}7\times10^{34}$~kg\,m$^2$\,s$^{-1}$, soit environ $1{,}6\,L_{\mathrm{EM}}$.
La fraction échangée avec le système Terre--disque,
\begin{equation}
\chi_J=\frac{|\Delta J_{\oplus D}|}{L_{B,\mathrm{rel}}},
\label{eq:chiJ}
\end{equation}
est une sortie de la Phase~B du programme numérique.

\begin{apport}
Sans aucune collision, la trajectoire de $B$ porte un moment cinétique de l'ordre de celui du
système Terre--Lune. Le passage rasant peut donc jouer, par le seul couple gravitationnel, le
rôle que le scénario d'impact attribue à la collision : fixer le budget de moment cinétique du
système Terre--disque. L'échange peut se faire dans les deux sens ; si la proto-Terre en rotation
rapide porte un spin supérieur à $L_{\mathrm{EM}}$ (tableau~\ref{tab:echelles}), la trajectoire de $B$
peut en recevoir une partie.
\end{apport}

\section{Intégrité de \texorpdfstring{$B$}{B} pendant la rencontre}

La marée terrestre à la surface de $B$ est caractérisée par
\begin{equation}
\mathcal{T}_B=2\,\frac{\Mt}{M_B}\left(\frac{R_B}{d}\right)^3,
\label{eq:TB}
\end{equation}
rapport entre l'accélération de marée et la gravité propre de $B$. La distance de Roche de $B$
autour de la proto-Terre est
\begin{equation}
d_{R,B}\simeq2{,}456\left(\frac{3\Mt}{4\pi\rho_B}\right)^{1/3}\approx2{,}75\,\Rt
\qquad(\rho_B=3{,}93\ \mathrm{g\,cm^{-3}}).
\label{eq:RocheB}
\end{equation}
La valeur $d_{R,B}\simeq2{,}75\,\Rt$ sert de repère de marée statique pour $B$. Elle ne constitue
pas une frontière absolue : la réponse réelle dépend de la durée du passage, de la structure
interne, du spin et de la cohésion effective de $B$.

En première approximation, pour $q_B\gtrsim d_{R,B}$, $B$ traverse la rencontre sans perte de masse significative ; pour
$q_B<d_{R,B}$, une fraction de sa matière peut être cédée au disque ($M_{B\to D}>0$) et doit être
suivie par la contrainte isotopique de la section~\ref{sec:bornes}. Le calcul du Niveau~I s'arrête
lorsque
\begin{equation}
|\mathbf{a}_{B,\mathrm{diff}}|<a_{\mathrm{tol}}
\label{eq:atol}
\end{equation}
sur l'ensemble du système Terre--disque, c'est-à-dire lorsque le forçage différentiel de $B$
devient dynamiquement négligeable.

\paragraph{La rotation de \texorpdfstring{$B$}{B} au premier ordre.}
La rotation propre de $B$ n'affecte pas son champ de marée au premier ordre, qui ne dépend que de
sa masse et de la distance. Elle joue sur deux quantités seulement : sa propre limite de
dislocation, légèrement abaissée s'il tourne vite, et la vitesse initiale de la matière qu'il
cède. Ces effets sont secondaires dans le domaine viable et ne modifient pas les conclusions du
chapitre. Le spin $\mathbf{S}_B$ est donc conservé comme paramètre descriptif de l'état de $B$
(équation~\eqref{eq:PsiB}), mais il n'entre pas dans le calcul du tri orbital.

\section{Après le passage : retour au problème héliocentrique}
\label{sec:retour-heliocentrique}

À la fin de la rencontre locale, $B$ n'appartient plus au système Terre--disque. Son état pertinent
n'est plus son périastre géocentrique, mais son orbite héliocentrique post-passage
$\mathcal{H}_{B,1}$ définie par l'équation~\eqref{eq:Theta-plus}. Le problème se décompose donc en
deux niveaux :
\begin{equation}
\boxed{\;\text{niveau local : } \Theta_I\to\mathcal{D}_1,
\qquad
\text{niveau héliocentrique : } \mathcal{H}_{B,1}\to B\in\mathcal{P}_{\mathrm{elim}}.\;}
\label{eq:deux-niveaux-B}
\end{equation}
Le premier niveau calcule la réorganisation du réservoir circumterrestre ; le second vérifie que
l'orbite laissée à $B$ peut être diffusée, déstabilisée puis éliminée dans une instabilité globale.

\begin{controle}
Le passage près de la Terre ne suffit pas à lui seul à faire disparaître $B$. La disparition de $B$
doit être testée par des intégrations héliocentriques incluant au minimum le Soleil, les planètes en
croissance et les géantes migrantes. La théorie demande donc deux fermetures : une fermeture
locale du bilan Terre--disque--$B$, puis une fermeture héliocentrique montrant que $B$ n'a pas à
survivre dans le Système solaire final.
\end{controle}

% =====================================================================

% ---------------------------------------------------------------------
\chapter{Injection de matière superficielle proto-terrestre}
\label{ch:injection}
% ---------------------------------------------------------------------

\section{L'océan magmatique et la bande intertropicale}
\label{sec:bande}

Avant toute mise en équation, il vaut la peine de décrire la scène, car la géométrie du phénomène
se comprend sans calcul et elle commande tout le reste.

\paragraph{L'état de la surface.} La proto-Terre considérée ici n'a pas de croûte solide au sens
actuel. Sa différenciation est déjà faite : le fer métallique, le nickel, le cobalt et les éléments
sidérophiles sont descendus vers le centre et forment le noyau. Ce qui reste au-dessus est un
manteau silicaté, chaud, largement fondu, dont la partie supérieure constitue un océan magmatique.
Sa surface n'est pas lisse. Elle est parcourue d'une croûte mince, brisée, constamment reformée et
constamment engloutie, dont l'aspect le plus proche que nous connaissions est celui des coulées
basaltiques de type \emph{aa} : un pavage de blocs scoriacés porté par une masse fluide, qui se
fracture dès qu'il est sollicité. Cette texture n'est pas un détail pittoresque. Elle signifie que
la cohésion mécanique de la surface est négligeable devant la gravité, et donc que la seule chose
qui retient la matière en place est le poids effectif qui s'exerce sur elle.

\paragraph{Pourquoi une bande, et pourquoi celle-là.} Deux effets se combinent, et tous deux sont
maximaux au même endroit.

Le premier est la rotation. Sur un corps qui tourne vite, l'accélération centrifuge s'oppose à la
gravité proportionnellement au carré du cosinus de la latitude. Elle est donc nulle aux pôles et
maximale à l'équateur. Le poids effectif est déjà le plus faible dans la ceinture équatoriale avant
même que $B$ n'arrive.

Le second est la marée de $B$. Le champ différentiel tire la matière vers l'extérieur au voisinage
du point situé directement sous $B$, et vers l'intérieur au-delà d'une distance angulaire de
$54{,}7^\circ$ de ce point. Si $B$ passe avec une faible inclinaison, son point sous-jacent balaie
précisément les basses latitudes.

Les deux effets se superposent donc dans la même ceinture, et nulle part ailleurs. C'est de là, et
de là seulement, que la matière peut partir.

\paragraph{La condition de soulèvement.} À l'équateur, la matière quitte la surface lorsque la
somme de la composante centrifuge et de la composante radiale de la marée dépasse la gravité. En
divisant par $G\Mt/\Req^2$, cette condition prend une forme sans dimension remarquablement simple :
\begin{equation}
\boxed{\;2\,\mu_B\left(\frac{\Req}{q_B}\right)^{3}+\tilde\omega^{2}\;\ge\;1.\;}
\label{eq:soulevement}
\end{equation}
Elle se lit directement. Le second terme est la part fournie par la rotation propre, le premier la
part fournie par $B$. Ni l'une ni l'autre ne suffit seule : une proto-Terre en rotation critique
atteint tout juste l'égalité sans $B$, et $B$ seul ne soulève rien sur un corps lent. Le
soulèvement est un effet conjoint, et c'est pourquoi le chemin étudié exige les deux conditions à
la fois.

\paragraph{L'étendue de la bande.} En reprenant la même comparaison à la latitude $\lambda$, où la
centrifuge varie en $\cos^2\lambda$ et la marée en $3\cos^2\lambda-1$, la zone soulevée s'étend
jusqu'à la latitude $\lambda_{\max}$ solution de
\begin{equation}
\mu_B\left(\frac{\Req}{q_B}\right)^{3}\left(3\cos^{2}\lambda-1\right)
\;\ge\;1-\tilde\omega^{2}\cos^{2}\lambda .
\label{eq:bande}
\end{equation}
Le résultat est le même sur tout le domaine viable, ce qui n'allait pas de soi :

\begin{center}
\begin{tabular}{cccc}
\toprule
$\tilde\omega$ & $\Req/\Rt$ & $\mu_B$ & $\lambda_{\max}$ \\
\midrule
$0{,}95$ & $2{,}0$ & $0{,}5$ & $22{,}6^\circ$ \\
$0{,}97$ & $2{,}1$ & $0{,}4$ & $23{,}5^\circ$ \\
$0{,}97$ & $2{,}0$ & $0{,}5$ & $24{,}4^\circ$ \\
$0{,}97$ & $1{,}9$ & $0{,}6$ & $24{,}8^\circ$ \\
$0{,}99$ & $2{,}0$ & $0{,}5$ & $26{,}1^\circ$ \\
\bottomrule
\end{tabular}
\end{center}

\begin{ideephys}[Une bande intertropicale]
La zone d'où part la matière s'étend, dans tout le domaine viable, entre $23^\circ$ et $26^\circ$
de latitude de part et d'autre de l'équateur. C'est, à un degré près, la bande intertropicale
terrestre actuelle. La coïncidence n'a évidemment aucune signification causale, mais elle donne une
image juste et immédiatement parlante : la matière dont la Lune est faite ne vient pas de toute la
surface de la proto-Terre, elle vient d'une ceinture équatoriale large d'environ un quart de la
distance du pôle. Tout le reste de la surface reste en place.
\end{ideephys}

\paragraph{Ce que cela implique pour la composition.} La conséquence géochimique suit
immédiatement, et elle n'est pas une hypothèse supplémentaire. La matière soulevée est celle du
sommet de l'océan magmatique, donc du silicate dont le métal est déjà parti. Sa pauvreté en fer, en
nickel et en éléments sidérophiles n'est pas un ingrédient qu'il faut ajouter au modèle : c'est une
propriété de l'endroit d'où elle part. Le mécanisme qui fixe la masse et celui qui fixe la
composition sont donc le même mécanisme, et c'est l'argument le plus économique de toute la
théorie.


Le réservoir préexistant fournit l'essentiel de la matière. Ce chapitre décrit le second canal, plus
limité : l'arrachement, par la marée de $B$, de matière superficielle des basses latitudes.

\section{Rôle de la rotation terrestre}

Le disque préexistant permet à la théorie de fonctionner même si la rencontre n'injecte qu'une
fraction de matière nouvelle. La rotation proto-terrestre garde un rôle central pour la sélection
de matière aux basses latitudes. La gravité effective monopolaire équatoriale et la latitude
dynamique sont
\begin{equation}
g_{\mathrm{eff,eq}}\simeq\frac{G\Mt}{\Req^2}\left(1-\tilde\omega^2\right),
\qquad
\lambda_\Omega=\arcsin\!\left(\hat{\mathbf{r}}\cdot\hat{\boldsymbol{\Omega}}_\oplus\right).
\label{eq:geff}
\end{equation}
La source candidate est une bande de basses latitudes autour de l'équateur rotationnel
instantané de la proto-Terre.

\section{Seuil d'injection}

La matière équatoriale est entraînée vers le disque lorsque la projection radiale de
l'accélération de marée de $B$, d'ordre $2F(\lambda_B)GM_B\Req/q_B^3$, dépasse la gravité
effective~\eqref{eq:geff}. Le seuil s'écrit
\begin{equation}
\boxed{\;1-\tilde\omega^2\lesssim2\,F(\lambda_B)\,\frac{M_B}{\Mt}\left(\frac{\Req}{q_B}\right)^{3}\;}
\label{eq:seuil}
\end{equation}
À latitude nulle, $F(0)=1$ : pour $M_B=0{,}1\,\Mt$ et $q_B=2\,\Req$, il donne
$\tilde\omega>0{,}987$ ; pour $M_B=0{,}5\,\Mt$ et $q_B=2{,}5\,\Req$,
$\tilde\omega>0{,}967$. Cette estimation reste statique ; pour une proto-Terre fortement oblate ou
proche de la rupture, le calcul doit utiliser le potentiel complet $\Phi_\oplus$. La Phase~B la
remplace par le traitement dynamique de la marée, gouverné par $\eta_\oplus$.

\begin{lecture}
L'injection de matière proto-terrestre par $B$ est d'autant plus efficace que la proto-Terre est
proche de sa rotation critique : rotation rapide et passage rasant agissent ensemble. Le disque
préexistant fournit l'essentiel du réservoir ; la matière injectée l'enrichit.
\end{lecture}

\section{Potentiel transitoire}

Dans le repère centré sur la proto-Terre et tournant avec elle, un potentiel instantané de
diagnostic est
\begin{equation}
\Phi_\star(\mathbf{r},t)=\Phi_\oplus(\mathbf{r})-\frac{GM_B}{|\mathbf{r}-\mathbf{r}_B(t)|}
+GM_B\,\frac{\mathbf{r}\cdot\mathbf{r}_B(t)}{|\mathbf{r}_B(t)|^3}
-\frac12\,|\boldsymbol{\Omega}_\oplus\times\mathbf{r}|^2,
\label{eq:Phistar}
\end{equation}
dont le gradient redonne l'accélération~\eqref{eq:adiff} augmentée des termes d'entraînement.
Les points-selles instantanés vérifient
\begin{equation}
\nabla\Phi_\star(\mathbf{r}_{\mathrm{col}},t)=\mathbf{0}
\label{eq:col}
\end{equation}
et localisent les passages gravitationnels transitoires du côté proche et du côté opposé à $B$.

\section{Flux de matière injectée}

La sortie matériellement pertinente et la masse totale injectée sont
\begin{equation}
\frac{\dd M_{\oplus\to D}}{\dd\lambda_\Omega\,\dd z\,\dd t},
\qquad
M_{\oplus\to D}=\iiint\frac{\dd M_{\oplus\to D}}{\dd\lambda_\Omega\,\dd z\,\dd t}\,
\dd\lambda_\Omega\,\dd z\,\dd t,
\label{eq:flux}
\end{equation}
où $z$ mesure la profondeur d'origine.

\begin{ideephys}
La préexistence du disque transforme le problème : la matière séparée de la proto-Terre devient
un flux d'enrichissement du réservoir circumterrestre, plutôt que l'unique origine de ce
réservoir.
\end{ideephys}

% =====================================================================

% ---------------------------------------------------------------------
\chapter{La fenêtre lunaire : le passage comme tri orbital}
\label{ch:fenetre}
% ---------------------------------------------------------------------

Ce chapitre formule le résultat central du passage. Le réservoir circumterrestre est déjà massif
avant la rencontre ; la fonction de $B$ devient alors très précise :

\begin{principe}{Formulation centrale}
Le disque secondaire circumterrestre est déjà massif avant la rencontre. Le passage non
collisionnel de $B$ agit comme un \emph{mécanisme de tri orbital} : son couple gravitationnel
redistribue le moment cinétique du disque. Une fraction de la matière est déplacée vers des états
liés dont le moment cinétique permet une circularisation au-delà de la limite de Roche ; une
fraction de plus faible moment cinétique réaccrète sur la proto-Terre, tandis que les éléments trop
énergétiques peuvent s'échapper. La Lune s'accrète à partir de la cohorte liée conservée dans la
fenêtre satellitaire.
\end{principe}

\section{Critère d'appartenance à la fenêtre lunaire}

Un élément de matière appartient à la fenêtre lunaire s'il vérifie simultanément
\begin{equation}
\boxed{\;\varepsilon<0,\qquad q_p>\Req,\qquad r_{\mathrm{circ}}=\frac{j^2}{G\Mt}\ge\aR.\;}
\label{eq:fenetre}
\end{equation}
Ces trois conditions disent respectivement : la matière reste liée à la Terre ; elle évite une
réaccrétion immédiate sur la proto-Terre ; son moment cinétique lui permet, après
circularisation, de rester dans la région où l'accrétion satellitaire peut se développer. Ici
$\varepsilon$ est l'énergie orbitale spécifique, $q_p$ la distance au périgée et $j$ le moment
cinétique spécifique.

\section{Une relation géométrique entre les trois conditions}

Pour une orbite képlérienne liée de demi-grand axe $a$ et d'excentricité $e$,
\begin{equation}
r_{\mathrm{circ}}=a(1-e^2)=q_p(1+e),\qquad q_p\le r_{\mathrm{circ}}<2q_p
\label{eq:rcirc-qp}
\end{equation}
(démonstration en annexe~\ref{app:rcirc}). La condition $r_{\mathrm{circ}}\ge\aR$ entraîne donc
$q_p>\aR/2\approx1{,}45\,\Rt$ : la condition de non-réaccrétion est presque automatiquement
satisfaite, sauf pour une proto-Terre très aplatie. \emph{C'est le moment cinétique qui décide.}
Pour les éléments inclinés, on utilise la projection $j_z$ sur le moment cinétique final du
réservoir, car les collisions de relaxation compensent les composantes perpendiculaires :
$r_{\mathrm{circ}}=j_z^2/(G\Mt)$.

\section{Mesure du succès du passage}

La grandeur qui mesure directement le succès du passage est la fraction de masse placée dans la
fenêtre lunaire :
\begin{equation}
\boxed{\;\eta_L=\frac{M\!\left(\varepsilon<0,\ q_p>\Req,\ r_{\mathrm{circ}}\ge\aR\right)}{M_{D,0}}.\;}
\label{eq:etaL}
\end{equation}
Les autres destins sont mesurés par
\begin{equation}
\begin{gathered}
\eta_{\mathrm{reaccr}}=\frac{M_{\mathrm{reaccr}}}{M_{D,0}},\qquad
\eta_{\mathrm{esc}}=\frac{M_{\mathrm{esc}}}{M_{D,0}},\qquad
\eta_{D\to B}=\frac{M_{\to B}}{M_{D,0}},\\[4pt]
\eta_{\mathrm{in}}=\frac{M(\varepsilon<0,\ q_p>\Req,\ r_{\mathrm{circ}}<\aR)}{M_{D,0}}.
\end{gathered}
\label{eq:etas}
\end{equation}
La fraction $\eta_{\mathrm{in}}$ forme le disque intérieur à Roche : elle n'est pas perdue et peut
nourrir l'accrétion lunaire plus tard par étalement visqueux \cite{SalmonCanup2012}. La partition
est complète :
\begin{equation}
\boxed{\;\eta_L+\eta_{\mathrm{in}}+\eta_{\mathrm{reaccr}}+\eta_{\mathrm{esc}}+\eta_{D\to B}=1.\;}
\label{eq:partition}
\end{equation}

Une partie du réservoir situé au-delà de Roche satisfait déjà le critère~\eqref{eq:fenetre} avant la
rencontre. La contribution propre de $B$ est donc
\begin{equation}
\Delta\eta_L=\eta_L-\eta_L^{(0)},\qquad
\eta_L^{(0)}=\frac{M(r>\aR)}{M_{D,0}}\ \text{avant la rencontre},
\label{eq:DeltaetaL}
\end{equation}
mesurée par comparaison avec un calcul témoin sans $B$.

\section{L'objectif numérique}

Le véritable objectif numérique de la théorie s'écrit
\begin{equation}
\boxed{\;\left(M_B,\,q_B,\,v_\infty,\,i_B,\,\omega_B,\,\mathcal{D}_0\right)\longrightarrow
\left(\eta_L,\,\Delta\eta_L,\,\eta_{\mathrm{in}},\,\eta_{\mathrm{reaccr}},\,\eta_{\mathrm{esc}},\,
L_{D,1},\,M_{D,1}\right).\;}
\label{eq:objectif}
\end{equation}
Le passage réussi n'est pas celui qui perturbe le plus fortement le système. C'est celui qui
maximise une fenêtre :
\begin{equation}
\boxed{\;\eta_L\ \text{élevé},\qquad \eta_{\mathrm{esc}}\ \text{modéré},\qquad
M_{\mathrm{reaccr}}\ \text{compatible},\;}
\label{eq:fenetre-succes}
\end{equation}
tout en fermant le bilan de moment cinétique~\eqref{eq:fermeture}.

\section{Anatomie du tri : l'impulsion tangentielle}

Considérons un élément sur une orbite circulaire de rayon $r_0$, qui reçoit une impulsion
tangentielle $\Delta v_t$. On pose $\delta=\Delta v_t/v_K(r_0)$. Son rayon de circularisation
devient
\begin{equation}
r_{\mathrm{circ}}=r_0\,(1+\delta)^2,
\label{eq:rcirc-delta}
\end{equation}
et les trois frontières du tri s'écrivent (annexe~\ref{app:tri}) :
\begin{equation}
\delta_{\mathrm{re}}(r_0)=\sqrt{\frac{2\Req}{r_0+\Req}}-1,\qquad
\delta_L(r_0)=\sqrt{\frac{\aR}{r_0}}-1,\qquad
\delta_{\mathrm{esc}}=\sqrt2-1\approx0{,}414.
\label{eq:frontieres}
\end{equation}
Un élément réaccrète si $\delta<\delta_{\mathrm{re}}$, rejoint le disque intérieur si
$\delta_{\mathrm{re}}<\delta<\delta_L$, entre dans la fenêtre lunaire si
$\delta_L\le\delta<\delta_{\mathrm{esc}}$ et s'échappe au-delà.

\begin{table}[htbp]
\centering
\small
\caption{Frontières du tri orbital pour une impulsion tangentielle ($\Req=\Rt$, $\aR=2{,}9\,\Rt$).}
\label{tab:frontieres}
\begin{tabular}{cccc}
\toprule
$r_0/\Rt$ & $\delta_{\mathrm{re}}$ (réaccrétion) & $\delta_L$ (fenêtre lunaire) & $\delta_{\mathrm{esc}}$ (échappement) \\
\midrule
2{,}0 & $-0{,}18$ & $+0{,}20$ & $+0{,}41$ \\
2{,}5 & $-0{,}24$ & $+0{,}08$ & $+0{,}41$ \\
3{,}0 & $-0{,}29$ & $-0{,}02$ & $+0{,}41$ \\
4{,}0 & $-0{,}37$ & $-0{,}15$ & $+0{,}41$ \\
5{,}0 & $-0{,}42$ & $-0{,}24$ & $+0{,}41$ \\
\bottomrule
\end{tabular}
\end{table}

Une impulsion purement radiale ne modifie pas $j$, donc pas $r_{\mathrm{circ}}$ : elle change
l'excentricité et l'énergie, et peut conduire à la réaccrétion ou à l'échappement, mais elle ne
déplace pas un élément vers la fenêtre lunaire. \emph{Le tri vers la fenêtre est gouverné par le
couple}, c'est-à-dire par la composante tangentielle de l'impulsion.

\section{Impulsion de marée : forme analytique et borne d'ordre de grandeur}

L'impulsion reçue par un élément du réservoir est donnée par la forme close de
l'annexe~\ref{app:impulsion}, valable sous les trois hypothèses énoncées : trajectoire rectiligne,
forçage impulsif, particule fixe. L'approximation au premier ordre
$\Delta\mathbf{v}\propto(0,y,-z)$ n'est fiable que pour $r\ll q_B$ ; elle perd toute précision dès
$r/q_B\gtrsim0{,}3$ et ne doit pas servir de référence dans la région qui alimente la fenêtre
lunaire.

Pour fixer les ordres de grandeur, on utilise la borne
\begin{equation}
\frac{\Delta v}{v_K}\sim\sqrt2\,\frac{\mu_B}{\sqrt{1+\mu_B}}
\left(\frac{r}{q_B}\right)^{3/2}
\label{eq:impulsion}
\end{equation}
(annexe~\ref{app:impulsion}). Pour $\mu_B=0{,}1$, l'amplitude vaut environ $0{,}13\,(r/q_B)^{3/2}$ ;
pour $\mu_B=0{,}3$, environ $0{,}37\,(r/q_B)^{3/2}$. Le signe de la composante tangentielle alterne
avec l'azimut de l'élément par rapport à $B$ : une partie du réservoir gagne du moment cinétique,
une autre en perd. C'est littéralement le tri orbital.

La correction de courbure hyperbolique est traitée en Phase~B par intégration directe sur la vraie
conique, et l'écart avec la forme close~\eqref{eq:impulsion-close} est une sortie du calcul. C'est
l'objet de la porte $\mathcal{D}_{\mathrm{imp}}$ définie ci-dessous.

\begin{figure}[htbp]
\centering
\begin{tikzpicture}
\begin{axis}[width=12.5cm,height=8.2cm,xmin=2,xmax=5,ymin=-0.5,ymax=0.5,
  xlabel={rayon initial $r_0/\Rt$},ylabel={$\delta=\Delta v_t/v_K$},
  /pgf/number format/use comma,tick label style={font=\small},label style={font=\small},
  axis on top]
  \addplot[name path=bas,draw=none,domain=2:5] {-0.5};
  \addplot[name path=haut,draw=none,domain=2:5] {0.5};
  \addplot[name path=re,thick,black,domain=2:5,samples=60] {sqrt(2/(x+1))-1};
  \addplot[name path=L,thick,red!60!black,domain=2:5,samples=60] {sqrt(2.9/x)-1};
  \addplot[name path=esc,thick,blue!60!black,domain=2:5] {0.41421};
  \addplot[gray!25] fill between[of=bas and re];
  \addplot[orange!15] fill between[of=re and L];
  \addplot[teal!22] fill between[of=L and esc];
  \addplot[blue!8] fill between[of=esc and haut];
  \addplot[dashed,thick,violet!70!black,domain=2:5,samples=40] {0.3721*(x/5)^1.5};
  \addplot[dashed,thick,violet!70!black,domain=2:5,samples=40] {-0.3721*(x/5)^1.5};
  \addplot[dotted,thick,violet!70!black,domain=2:5,samples=40] {0.1348*(x/5)^1.5};
  \addplot[dotted,thick,violet!70!black,domain=2:5,samples=40] {-0.1348*(x/5)^1.5};
  \node[font=\small] at (axis cs:4.3,0.46) {échappement};
  \node[font=\small,teal!40!black] at (axis cs:4.35,0.0) {fenêtre lunaire};
  \node[font=\small,orange!60!black] at (axis cs:2.35,0.0) {disque intérieur};
  \node[font=\small] at (axis cs:3.2,-0.44) {réaccrétion};
  \node[font=\scriptsize,violet!70!black,align=left,anchor=west] at (axis cs:2.05,0.33) {tirets : $\pm\Delta v/v_K$, $\mu_B=0{,}3$\\pointillés : $\mu_B=0{,}1$ ($q_B=5\,\Rt$)};
\end{axis}
\end{tikzpicture}
\caption{Carte analytique du tri orbital pour une impulsion tangentielle : réaccrétion,
disque intérieur à Roche, fenêtre lunaire et échappement, en fonction du rayon initial $r_0$ et de
l'impulsion relative $\delta$. Les courbes violettes donnent l'amplitude d'ordre de grandeur de
l'impulsion de marée~\eqref{eq:impulsion} pour un passage à $q_B=5\,\Rt$. Carte d'ordre de
grandeur, que la Phase~B remplace par le calcul exact.}
\label{fig:tri}
\end{figure}

\begin{controle}
La figure~\ref{fig:tri} montre que la fenêtre lunaire est accessible : pour $\mu_B\sim0{,}3$ et un
passage à $q_B=5\,\Rt$, la matière initialement vers $2{,}5$--$3\,\Rt$ reçoit des impulsions qui la
font entrer dans la fenêtre, sans atteindre le seuil d'échappement ni celui de réaccrétion. Les
limites exactes de cette région sont une sortie de la Phase~B.
\end{controle}

\section{Du tri orbital au satellite dominant}

Le tri orbital donne aussi une explication physique à la formation d'un satellite dominant. Si $B$
excite un réservoir contenant déjà des agrégats ou des lunules, il peut provoquer
\begin{equation}
\begin{gathered}
\text{excitation des orbites}\rightarrow\text{croisements orbitaux}\rightarrow
\text{collisions entre corps circumterrestres}\\
\rightarrow\text{coalescence}\rightarrow f_{\mathrm{dom}}\rightarrow1.
\end{gathered}
\label{eq:chaine-dominance}
\end{equation}
Les collisions concernent ici la matière du réservoir et ses agrégats ; la rencontre Terre--$B$
demeure non collisionnelle. Les simulations à $N$ corps d'un disque circumterrestre montrent
d'ailleurs qu'un tel disque tend à former un satellite unique \cite{Kokubo2000}.

\begin{consequence}
Dans les simulations de Phase~C, le nombre de grands corps survivants doit décroître après le
passage, et la masse satellitaire doit se concentrer dans un corps dominant : $N_{\mathrm{sat}}\to1$,
$f_{\mathrm{dom}}\to1$.
\end{consequence}

% ---------------------------------------------------------------------
\chapter{Bilans de la rencontre}
\label{ch:bilans}
% ---------------------------------------------------------------------

La rencontre s'achève lorsque le forçage de $B$ devient négligeable. Ce chapitre ferme ses bilans
de masse et de moment cinétique, dans le langage des réservoirs : la masse du réservoir se
répartit selon la partition~\eqref{eq:partition}, et le moment cinétique total~\eqref{eq:conservation-J}
est conservé.

\section{Bilan de masse}

Le disque post-rencontre vérifie
\begin{equation}
M_{D,1}=M_{D,0}+M_{\oplus\to D}+M_{B\to D}-M_{\mathrm{reaccr}}-M_{\mathrm{esc}}-M_{\to B}.
\label{eq:bilanM1}
\end{equation}
Il se décompose par provenance, chaque terme étant net des pertes subies par la matière
correspondante :
\begin{equation}
M_{D,1}=M_{D,1}^{(0)}+M_{D,1}^{(\oplus)}+M_{D,1}^{(B)},
\label{eq:provenance}
\end{equation}
où $M_{D,1}^{(0)}$ est la masse du disque préexistant conservée dans le disque réorganisé,
$M_{D,1}^{(\oplus)}$ la masse proto-terrestre injectée et retenue, et $M_{D,1}^{(B)}$ la masse
éventuellement cédée par $B$ et retenue. Cette décomposition alimente directement les fractions
géochimiques du chapitre~\ref{sec:geochimie}.

\section{Bilan de moment cinétique}

Le bilan est écrit dans le repère barycentrique du couple Terre--$B$, inertiel à l'échelle de la
rencontre (le couple solaire est négligeable sur la durée $t_{p}$). Le repère
géocentrique de l'équation~\eqref{eq:adiff} n'étant pas inertiel, il sert au calcul des
trajectoires mais pas à la fermeture du bilan. Avant et après la rencontre,
\begin{align}
J_{\mathrm{avant}}&=S_{\oplus,0}+S_{B,0}+L_{D,0}+L_{\mathrm{orb},0},
\label{eq:Javant}\\
J_{\text{après}}&=S_{\oplus,1}+S_{B,1}+L_{D,1}+L_{\mathrm{orb},1}+L_{\mathrm{esc}},
\label{eq:Japres}
\end{align}
où $L_{\mathrm{orb}}$ est le moment orbital du mouvement relatif Terre--$B$ et $S_B$ le spin de $B$,
modifié par les couples de marée. La fermeture exige
\begin{equation}
\boxed{\;J_{\mathrm{avant}}=J_{\text{après}}\;}
\label{eq:fermeture}
\end{equation}
et le système Terre--disque laissé par le passage est caractérisé par
\begin{equation}
J_{\oplus D,1}=S_{\oplus,1}+L_{D,1}.
\label{eq:J1}
\end{equation}

\begin{controle}
Le transfert de moment cinétique vers la trajectoire de $B$ est une sortie de fermeture du bilan
pendant la rencontre. Le diagnostic central est l'état Terre--disque produit après le passage,
$J_{\oplus D,1}$, à comparer au moment cinétique du système Terre--Lune.
\end{controle}

% =====================================================================

\begin{rappel}
À la fin de la rencontre, $B$ a cessé d'agir. Le réservoir circumterrestre a été trié : une cohorte liée
occupe la fenêtre lunaire, une autre forme le disque intérieur, une partie a réaccrété ou s'est
échappée. La partie suivante suit l'évolution de ce réservoir, depuis son alimentation avant la
rencontre jusqu'à l'accrétion de la Lune.
\end{rappel}
% =====================================================================
\part{Le disque secondaire, réservoir évolutif}
\label{part:reservoir}
% =====================================================================

\noindent
Le disque secondaire n'est pas traité ici comme un objet figé, mais comme un réservoir
évolutif $\RD(t)$ : il est alimenté, il transporte, il se refroidit, il condense, il forme des agrégats,
il est trié par le passage de $B$, puis il se relaxe et accrète la Lune. Son état à chaque instant
combine des composantes de natures différentes,
\begin{equation}
\RD(t)=\mathcal{D}_{\mathrm{gaz}}(t)\cup\mathcal{D}_{\mathrm{multiphase}}(t)\cup
\mathcal{D}_{\mathrm{particulaire}}(t)\cup\mathcal{G}(t),
\label{eq:RD}
\end{equation}
où $\mathcal{G}(t)$ désigne la population d'agrégats et de lunules. Les trois chapitres de cette
partie suivent ce réservoir avant, pendant et après le passage.

% ---------------------------------------------------------------------
\chapter{Alimentation et structure du réservoir avant le passage}
\label{ch:structure}
% ---------------------------------------------------------------------

\section{Origine accrétionnelle du réservoir}
\subsection{Apport de matière dans la sphère de Hill}

Le rayon de Hill de la proto-Terre sur son orbite héliocentrique s'écrit
\begin{equation}
R_H=a_\oplus\left(\frac{\Mt}{3M_\odot}\right)^{1/3}.
\label{eq:RH}
\end{equation}
La matière qui pénètre dans cette région peut posséder un moment cinétique géocentrique non
nul. L'accrétion de galets transmet un moment cinétique prograde \cite{JohansenLacerda2010,Visser2020},
et les simulations hydrodynamiques résolvent les flux de gaz et de galets autour d'embryons
rocheux jusqu'à une masse terrestre \cite{Popovas2018}. La physique générale des disques
d'accrétion relie ensuite dissipation et transport du moment cinétique \cite{Pringle1981}.

Le bilan de masse du disque s'écrit
\begin{equation}
\frac{\dd M_D}{\dd t}=\dot M_{\mathrm{in}}-\dot M_{\oplus}-\dot M_{\mathrm{out}},
\label{eq:bilanM}
\end{equation}
où $\dot M_{\mathrm{in}}$ est donné par l'équation~\eqref{eq:canaux}, $\dot M_{\oplus}$ désigne la
réaccrétion sur la proto-Terre et $\dot M_{\mathrm{out}}$ les pertes vers l'extérieur.

\subsection{Masse et densité surfacique}

La masse instantanée du disque et son rapport à la masse proto-terrestre sont
\begin{equation}
M_D(t)=2\pi\int_{R_{\mathrm{in}}(t)}^{R_{\mathrm{out}}(t)}\Sigma(r,t)\,r\,\dd r,
\qquad
\mu_D(t)=\frac{M_D(t)}{\Mt}.
\label{eq:MD}
\end{equation}

\begin{lecture}
$M_D$ est une sortie du bilan d'alimentation et de pertes. Le rapport d'environ $10^{-4}$ obtenu
pour les planètes géantes dans des modèles alimentés en continu \cite{CanupWard2006} est une
référence astrophysique ; il ne fixe pas la valeur de $\mu_D$ pour la proto-Terre, dont le mode
d'alimentation est différent (section~\ref{sec:soustherm}).
\end{lecture}

\subsection{Moment cinétique apporté par l'accrétion}

Pour une surface de contrôle fermée $\mathcal{S}_H$ entourant la région d'influence terrestre, de
normale sortante $\mathbf{n}$, le flux entrant de moment cinétique est
\begin{equation}
\dot L_{\mathrm{in}}=-\oint_{\mathcal{S}_H}\rho\,(\mathbf{r}\times\mathbf{v})_z\,
(\mathbf{v}\cdot\mathbf{n})\,\dd\mathcal{S},
\label{eq:Lin}
\end{equation}
le signe moins rendant positif un apport entrant. Le bilan de moment cinétique du disque est
\begin{equation}
\frac{\dd L_D}{\dd t}=\dot L_{\mathrm{in}}-\dot L_{\oplus}-\dot L_{\mathrm{out}}
+\mathcal{T}_{\oplus D}+\mathcal{T}_{\odot D},
\label{eq:bilanL}
\end{equation}
avec $\mathcal{T}_{\oplus D}$ et $\mathcal{T}_{\odot D}$ les couples exercés par la proto-Terre et par
le Soleil.

\subsection{Origine du spin proto-terrestre}

L'accrétion de galets peut transmettre un moment cinétique prograde au corps central
\cite{JohansenLacerda2010,Visser2020}. La rotation de la proto-Terre appartient donc au même
problème d'accrétion que le disque. On définit la rotation réduite
\begin{equation}
\tilde\omega=\frac{\Omega_\oplus}{\sqrt{G\Mt/\Req^3}}.
\label{eq:omegatilde}
\end{equation}
Pour une contraction approximativement conservatrice du spin,
\begin{equation}
S_\oplus\simeq k\,\Mt\Req^2\Omega_\oplus\simeq\text{constante}
\quad\Longrightarrow\quad
\tilde\omega\propto\Req^{-1/2}
\label{eq:spin}
\end{equation}
lorsque $k$ et $\Mt$ varient peu : la contraction thermique d'une proto-Terre chaude la
rapproche de la rotation critique.

\begin{ideephys}
Le disque et le spin terrestre ne sont pas deux réservoirs indépendants ajoutés l'un à l'autre. Ils
résultent tous deux de la manière dont matière et moment cinétique sont distribués pendant la
croissance de la proto-Terre ; le canal $\dot M_{\mathrm{rot}}$ les relie directement.
\end{ideephys}

% =====================================================================

\section{Structure radiale}
\subsection{Rayon interne}

Le rayon interne correspond à la transition entre la proto-Terre ou son enveloppe dense et un
régime réellement orbital :
\begin{equation}
R_{\mathrm{in}}\simeq R_{\mathrm{tr}}.
\label{eq:Rin}
\end{equation}
Cette transition dépend de la pression, de la rotation, de la température et, si un champ
magnétique significatif est présent, de contraintes magnétohydrodynamiques. Pour une
proto-Terre proche de la rotation critique, $R_{\mathrm{tr}}$ tend vers $\Req$.

\subsection{Rayon externe}

Le rayon externe est paramétré par
\begin{equation}
R_{\mathrm{out}}=\xi_H R_H,\qquad 0<\xi_H<1.
\label{eq:Rout}
\end{equation}
Pour les disques circumplanétaires de planètes géantes, la troncature par les couples solaires
conduit à $\xi_H\approx0{,}3$--$0{,}4$ \cite{QuillenTrilling1998,MartinLubow2011} ; dans la présente
théorie, $\xi_H$ est une sortie de la dynamique héliocentrique et des conditions d'alimentation.
L'essentiel de la masse d'un disque alimenté en solides se concentre toutefois bien à l'intérieur
de ce bord, au voisinage du rayon de circularisation de la matière capturée.

\paragraph{La marée solaire ne déforme pas le disque.}
Le rayon de Laplace, où le Soleil prend le pas sur le bourrelet équatorial, se situe entre 26 et
34 rayons terrestres pour une proto-Terre aplatie. Le réservoir s'arrêtant à $2{,}9\,\Rt$, il est
entièrement dans le domaine où le bourrelet domine, d'un facteur dix. La précession du nœud y est
de vingt-quatre heures au bord interne et de quatre-vingt-dix heures au bord externe : le disque se
met à plat dans le plan équatorial en quelques jours, bien avant l'arrivée de $B$. Cette mise à
plat est cohérente avec ce que dit la littérature du domaine, y compris
\cite{PahlevanMorbidelli2015}. La question de l'alignement du disque sur l'équateur rotationnel
instantané de la proto-Terre est donc close : elle est acquise avant la rencontre.

\subsection{Limite de Roche et orbite synchrone}

Pour une densité satellitaire de référence $\rho_s$, la limite de Roche fluide est
\begin{equation}
\aR\simeq2{,}456\left(\frac{3\Mt}{4\pi\rho_s}\right)^{1/3}
\approx2{,}9\,\Rt\quad(\rho_s=3{,}34\ \mathrm{g\,cm^{-3}}).
\label{eq:Roche}
\end{equation}
Son rôle est distinct de celui de $R_{\mathrm{in}}$ :
\begin{equation}
R_{\mathrm{in}}\neq\aR.
\label{eq:RinRoche}
\end{equation}
Un second rayon joue un rôle dynamique essentiel, l'orbite synchrone
\begin{equation}
a_{\mathrm{sync}}=\left(\frac{G\Mt}{\Omega_\oplus^2}\right)^{1/3}\approx2{,}0\text{--}2{,}3\,\Rt
\quad(P_\oplus=4\text{--}5\ \mathrm{h}).
\label{eq:async}
\end{equation}
Pour une proto-Terre en rotation rapide, $a_{\mathrm{sync}}<\aR$ : tout corps assemblé au-delà de
la limite de Roche migre vers l'extérieur sous l'effet des marées terrestres.

\begin{lecture}
Un disque peut occuper la région intérieure à la limite de Roche. Cette limite contrôle la survie
et l'assemblage des grands agrégats autogravitants ; elle est centrale pour l'accrétion lunaire,
non pour l'existence du disque. L'ordre $a_{\mathrm{sync}}<\aR$ garantit que les corps formés
au-delà de Roche s'éloignent de la Terre plutôt que d'y retomber.
\end{lecture}

\begin{figure}[htbp]
\centering
\begin{tikzpicture}
  \draw[-{Stealth}] (0,0) -- (14.2,0) node[below left,font=\small] {distance géocentrique};
  \fill[orange!70!brown] (0,0) circle (0.35);
  \node[font=\small] at (0,-0.65) {Terre};
  \foreach \x/\lab/\c in {1.2/{$R_{\mathrm{in}}$}/black,2.3/{$a_{\mathrm{sync}}$}/teal!60!black,3.4/{$\aR$}/red!60!black,9.5/{$R_{\mathrm{out}}$}/black,13.2/{échelle de Hill}/gray}{
    \draw[\c,thick] (\x,-0.25) -- (\x,0.25);
    \node[\c,font=\small] at (\x,0.6) {\lab};
  }
  \fill[blue!25,opacity=0.7] (1.2,-0.12) rectangle (9.5,0.12);
  \node[font=\scriptsize,blue!45!black] at (6.3,-0.45) {disque secondaire};
  \node[font=\scriptsize,red!60!black,align=center] at (5.2,1.25) {assemblage de corps\\autogravitants};
  \draw[red!60!black,-{Stealth}] (3.6,1.05) -- (7.0,1.05);
\end{tikzpicture}
\caption{Organisation radiale conceptuelle du disque circumterrestre. Les rayons interne,
synchrone, de Roche, externe et de Hill jouent des rôles physiques différents (échelle non
respectée).}
\label{fig:radial}
\end{figure}

% =====================================================================

\section{Le réservoir comme disque d'accrétion}
\label{sec:disque-accretion}

Il est utile, avant d'entrer dans la structure particulière de $\RD$, de rappeler qu'un disque
d'accrétion est une structure \emph{générique} : de la matière en orbite autour d'un corps central,
une dissipation visqueuse qui transforme l'énergie orbitale en chaleur, et un transport de moment
cinétique vers l'extérieur qui autorise l'accrétion vers l'intérieur. Les trois ingrédients sont
solidaires. C'est exactement la structure que l'on propose pour $\RD$, et les outils développés pour
les disques d'accrétion en général \cite{ShakuraSunyaev1973,Pringle1981,FrankKingRaine2002}
s'appliquent donc à lui sans transposition particulière.

\paragraph{Un analogue pour l'alimentation.} La formation d'un disque par transfert de masse à
travers le point de Lagrange intérieur d'un système binaire est un mécanisme connu et robuste
\cite{FrankKingRaine2002}. Le corps qui remplit son lobe de Roche cède de la matière à son
compagnon, et cette matière, trop riche en moment cinétique pour tomber directement, se circularise
en disque. On propose de voir dans le flux $\Phi_{\oplus\to D}$ l'analogue de ce mécanisme : la
proto-Terre, portée près de sa limite de rotation et déformée par le champ différentiel de $B$,
cède de la matière superficielle qui se circularise à son tour. L'analogie n'est pas une preuve, et
les géométries diffèrent ; elle indique seulement que le type de flux invoqué n'est pas exotique.

\paragraph{Le problème central est le transport du moment cinétique.} Dans tout disque
d'accrétion, la perte de moment cinétique de la matière qui tombe vers le centre doit être
compensée par un gain de moment cinétique de la matière située plus loin. Sans ce transfert, rien
n'accrète. C'est précisément la fonction que l'on attribue au tri orbital du chapitre
\ref{ch:fenetre} : le couple exercé par $B$ déplace du moment cinétique vers l'extérieur du
réservoir, ce qui permet à une fraction de la matière de franchir $\aR$ pendant qu'une autre
fraction retombe vers l'intérieur. Le passage joue donc, en une fois et de manière impulsive, le
rôle que la viscosité joue lentement et continûment.

\paragraph{Une voie de transport supplémentaire à envisager.} La viscosité effective
$\nu\sim G^2\Sigma^2/\Omega_K^3$ retenue plus haut suppose un transport par couples gravitationnels.
Si le réservoir conserve du gaz résiduel ou une fraction ionisée, l'instabilité
magnéto-rotationnelle \cite{BalbusHawley1991} devient un mécanisme de redistribution du moment
cinétique très efficace dans les disques faiblement magnétisés, et pourrait compléter ou remplacer
cette viscosité effective. Cette voie n'est pas retenue ici comme hypothèse de base, parce que le
degré d'ionisation de $\RD$ n'est pas contraint ; on propose de la traiter comme une branche
alternative du domaine viable, à activer si l'état thermique calculé en Phase~A indique une
ionisation suffisante.

\section{Structure orbitale et transport du moment cinétique}
\subsection{Rotation quasi képlérienne}

Lorsque la masse du disque reste petite devant celle de la proto-Terre et que la pression fournit
une correction modérée, la fréquence orbitale dominante est
\begin{equation}
\Omega_K(r)=\sqrt{\frac{G\Mt}{r^3}}.
\label{eq:OmegaK}
\end{equation}
La vitesse azimutale réelle s'écrit
\begin{equation}
v_\phi^2=\frac{G\Mt}{r}+\frac{r}{\rho}\frac{\partial P}{\partial r}+r\,a_{\mathrm{sg}}(r)
+r\,a_{J_2}(r),
\label{eq:vphi}
\end{equation}
où $a_{\mathrm{sg}}$ représente l'autogravitation du disque et $a_{J_2}$ la contribution de
l'aplatissement d'une proto-Terre en rotation rapide. Le moment cinétique spécifique est
$j(r)=r\,v_\phi(r)$.

\subsection{Moment cinétique global du disque}

Dans la limite képlérienne,
\begin{equation}
L_D=2\pi\int_{R_{\mathrm{in}}}^{R_{\mathrm{out}}}\Sigma(r)\sqrt{G\Mt r}\;r\,\dd r,
\label{eq:LD}
\end{equation}
et le système initial Terre--disque possède
\begin{equation}
J_{\oplus D,0}=S_{\oplus,0}+L_{D,0}.
\label{eq:J0}
\end{equation}

\begin{ideephys}
Le moment cinétique destiné à la future Lune peut être en partie orbital avant le passage de
$B$. La rencontre n'a donc pas à créer tout le budget orbital : elle redistribue un budget déjà
présent et y ajoute ou en retire la part échangée avec la trajectoire de $B$.
\end{ideephys}

\subsection{Évolution visqueuse}

Dans l'approximation d'un disque mince, axisymétrique et quasi képlérien, la densité surfacique
suit l'équation classique \cite{Pringle1981}
\begin{equation}
\frac{\partial\Sigma}{\partial t}=\frac{3}{r}\frac{\partial}{\partial r}
\left[r^{1/2}\frac{\partial}{\partial r}\left(\nu\Sigma r^{1/2}\right)\right].
\label{eq:visc}
\end{equation}
Pour une composante gazeuse, $\nu=\alpha c_sH$ avec $H\simeq c_s/\Omega_K$. Pour une
composante particulaire dense, la viscosité effective est dominée par les sillages
gravitationnels et croît comme $G^2\Sigma^2/\Omega_K^3$ \cite{Daisaka2001}. Le temps
visqueux caractéristique est
\begin{equation}
t_\nu(r)\sim\frac{r^2}{\nu}.
\label{eq:tnu}
\end{equation}

% =====================================================================

% ---------------------------------------------------------------------
\chapter{Thermodynamique, stabilité et persistance}
\label{ch:thermo}
% ---------------------------------------------------------------------

\section{État thermique et multiphasé}
\subsection{Bilan énergétique}

Le bilan énergétique local du disque est
\begin{equation}
Q^{+}_{\mathrm{visc}}+Q^{+}_{\mathrm{acc}}+Q^{+}_{\oplus}+Q^{+}_{\odot}
=Q^{-}_{\mathrm{rad}}+Q^{-}_{\mathrm{adv}}+Q^{-}_{\mathrm{phase}}.
\label{eq:bilanE}
\end{equation}
Dans un disque quasi képlérien, la dissipation visqueuse par unité de surface et le
refroidissement radiatif des deux faces sont
\begin{equation}
Q^{+}_{\mathrm{visc}}\simeq\frac94\,\nu\Sigma\Omega_K^2,
\qquad
Q^{-}_{\mathrm{rad}}\simeq2\sigma_{\mathrm{SB}}T_{\mathrm{eff}}^4.
\label{eq:Qvisc}
\end{equation}
Avec l'épaisseur optique verticale $\tau\simeq\kappa\Sigma/2$, l'approximation de diffusion
donne pour un milieu optiquement épais
\begin{equation}
T_{\mathrm{mid}}^4\simeq\frac34\,\tau\,T_{\mathrm{eff}}^4+T_{\mathrm{irr}}^4.
\label{eq:Tmid}
\end{equation}

\begin{lecture}
La température intervient simultanément dans l'épaisseur, la viscosité, l'opacité, la vitesse du
son et l'état de phase. La thermodynamique participe directement à la structure orbitale et à la
durée de vie du disque ; elle n'est pas un module ajouté après la dynamique.
\end{lecture}

\subsection{Condensation et transitions de phase}

Pour chaque espèce $i$, l'état de phase est une fonction
\begin{equation}
\mathcal{P}_i=\mathcal{P}_i\!\left[T(r,t),\,P(r,t),\,f_{\mathrm{O_2}}(r,t),\,\mathbf{X}(r,t)\right],
\label{eq:phase}
\end{equation}
et la chaîne d'évolution solide s'écrit
\begin{equation}
\text{vapeur}\rightarrow\text{condensats}\rightarrow\text{grains}\rightarrow
\text{agrégats}\rightarrow\text{satellitésimaux}.
\label{eq:chaine}
\end{equation}
Pour une population de particules de densité numérique $n_p$, de section efficace $\sigma_c$ et
de vitesse relative $\Delta v$, le temps de collision a pour échelle
\begin{equation}
t_{\mathrm{coll}}\sim\frac{1}{n_p\,\sigma_c\,\Delta v}.
\label{eq:tcoll}
\end{equation}

\subsection{Sédimentation verticale}

Une particule couplée au gaz est caractérisée par son nombre de Stokes
$\mathrm{St}=t_{\mathrm{stop}}\Omega_K$. La compétition entre sédimentation et turbulence fixe
son épaisseur verticale :
\begin{equation}
\frac{H_p}{H}\sim\left(\frac{\alpha}{\alpha+\mathrm{St}}\right)^{1/2}.
\label{eq:Hp}
\end{equation}
Cette relation suit l'apparition d'un plan médian enrichi en solides.

% =====================================================================

\section{Stabilité gravitationnelle et persistance}
\subsection{Critère de Toomre}

La stabilité locale d'un disque mince autogravitant est mesurée par le paramètre de Toomre,
dont la constante dépend de la nature du milieu :
\begin{equation}
Q_T^{\mathrm{gaz}}=\frac{c_s\,\kappa_{\mathrm{epi}}}{\pi G\Sigma}
\quad\text{\cite{Safronov1960,GoldreichLyndenBell1965}},
\qquad
Q_T^{\mathrm{part}}=\frac{\sigma_r\,\kappa_{\mathrm{epi}}}{3{,}36\,G\Sigma}
\quad\text{\cite{Toomre1964}},
\label{eq:Toomre}
\end{equation}
où $\kappa_{\mathrm{epi}}$ est la fréquence épicyclique et $\sigma_r$ la dispersion radiale des
vitesses des particules. Un disque froid et dense s'approche du régime $Q_T\sim1$.

Il est instructif d'évaluer ce critère sur un réservoir de quelques masses lunaires. Pour
$M_{D,0}=5{,}4\,M_L$ réparti selon $\Sigma\propto r^{-1}$ entre $\Req$ et $\aR$, la densité
surfacique vaut $\Sigma\simeq7{,}6\times10^{8}$~kg\,m$^{-2}$ au bord interne et
$4{,}5\times10^{8}$~kg\,m$^{-2}$ au bord externe, et la stabilité $Q_T^{\mathrm{part}}\ge1$ demande
\begin{equation}
\sigma_r\gtrsim0{,}06\text{--}0{,}09\;v_K\qquad(330\text{--}400\ \text{m\,s}^{-1}).
\label{eq:sigmar-reservoir}
\end{equation}
Un réservoir de cette densité n'est donc pas un disque de poussière quiescent : soit il est chaud,
fluide ou partiellement fondu, soit il est déjà organisé en lunules. Le calcul de Phase~A tranche
entre ces deux états, et ce choix détermine la nature des équations de Phase~B, puisqu'un réservoir
fluide relève de l'hydrodynamique et un réservoir de lunules d'un problème à $N$ corps. Plus le
réservoir est léger, plus $\sigma_r$ requis diminue et plus le régime de particules indépendantes
redevient légitime : c'est un argument de plus pour balayer $M_{D,0}$.

\begin{ideephys}
La préexistence du disque pose une question décisive : lorsque $B$ arrive, le disque est-il
encore un réservoir diffus, ou a-t-il déjà formé des agrégats et des lunules ? La réponse est
calculable par la thermodynamique, la stabilité gravitationnelle et les temps de croissance.
\end{ideephys}

\subsection{Persistance du disque}

La persistance est exprimée par l'état du réservoir à $t_B$ plutôt que par une durée arbitraire :
\begin{equation}
M_D(t_B)+\sum_k m_k(t_B)\ \ge\ M_{D,\min},\qquad
L_D(t_B)+\sum_k L_k(t_B)\ \ge\ L_{D,\min},\qquad
R_{\mathrm{out}}(t_B)>R_{\mathrm{in}}(t_B).
\label{eq:persistance}
\end{equation}
À cette condition globale s'ajoute une condition locale, qui compare la durée de vie du réservoir au
délai réellement disponible avant la rencontre :
\begin{equation}
\boxed{\;
t_{\mathrm{pers}}(r)=
\min\!\left(t_\nu(r),\ t_{\mathrm{frag}}(r),\ t_{\mathrm{reaccr}}(r),\ t_{\mathrm{diff}}(r)\right)
>t_B-t_{D,0}.
\;}
\label{eq:tpers}
\end{equation}
Les quatre temps sont respectivement l'étalement visqueux~\eqref{eq:tnu-reservoir}, la
fragmentation gravitationnelle associée à $Q_T$, la réaccrétion sur la proto-Terre et la diffusion
du réservoir. La condition doit être vérifiée sur tout l'intervalle radial dense, et non seulement
en moyenne.

\begin{falsif}
Si les lois d'alimentation, de refroidissement et de transport dissipent ou réaccrètent
systématiquement le réservoir avant toute fenêtre compatible avec le passage de $B$, la branche
correspondante de la théorie est exclue.
\end{falsif}

% =====================================================================

Les trois régimes chronologiques qui décrivent l'état du réservoir à l'instant du passage ont été
présentés au chapitre~\ref{ch:chronologie} (section~\ref{sec:regimes}).

% ---------------------------------------------------------------------
\chapter{Réorganisation du réservoir et accrétion lunaire}
\label{ch:reorganisation}
% ---------------------------------------------------------------------

Après le tri orbital, le réservoir entre dans une phase de relaxation dissipative, puis d'accrétion.
Ce chapitre décrit le nouvel état du réservoir, sa dissipation, la masse disponible au-delà de Roche
et le contrôle de la masse satellitaire accessible.

\section{Nouvel état du disque}

Le passage fournit
\begin{equation}
\mathcal{D}_1=\left\{M_{D,1},\ \Sigma_1,\ T_1,\ \mathbf{X}_1,\ L_{D,1},\ e(r),\ i(r),\
\mathbf{f}_{\mathrm{phase},1},\ \mathcal{G}_1\right\},
\label{eq:D1}
\end{equation}
et l'évolution collective suit
\begin{equation}
\mathcal{D}_1\longrightarrow\mathcal{D}_{\mathrm{relax}}\longrightarrow\mathcal{D}_{\mathrm{accr}}
\longrightarrow\text{proto-Lune}.
\label{eq:evolution}
\end{equation}

\section{Dissipation}

Le bilan énergétique de relaxation s'écrit
\begin{equation}
E_{\mathrm{diss}}=E_{\mathrm{choc}}+E_{\mathrm{circ}}+E_{\mathrm{vis}}+E_{\mathrm{coll}}
-E_{\mathrm{rad}}-E_{\mathrm{esc}}.
\label{eq:Ediss}
\end{equation}
Une parcelle sur une orbite de demi-grand axe $a$ et d'excentricité $e$, circularisée à moment
cinétique conservé, rejoint le rayon $a_{\mathrm{circ}}=a(1-e^2)$. L'énergie spécifique libérée,
positive, est
\begin{equation}
\Delta\varepsilon_{\mathrm{circ}}=\frac{G\Mt}{2a_{\mathrm{circ}}}-\frac{G\Mt}{2a}
=\frac{G\Mt}{2a}\,\frac{e^2}{1-e^2}\ \ge0.
\label{eq:circ}
\end{equation}

\section{Masse disponible au-delà de Roche}

La masse et le moment cinétique situés au-delà de la limite de Roche sont
\begin{equation}
M_{>R}=2\pi\int_{\aR}^{R_{\mathrm{out}}}\Sigma(r)\,r\,\dd r,
\qquad
L_{>R}=2\pi\int_{\aR}^{R_{\mathrm{out}}}\Sigma(r)\sqrt{G\Mt r}\;r\,\dd r.
\label{eq:MR}
\end{equation}

\section{Contrôle d'accrétion disque--satellite}

Une fois le disque réorganisé établi, on définit le moment cinétique spécifique normalisé
\begin{equation}
\ell_D=\frac{L_D}{M_D\sqrt{G\Mt\aR}}.
\label{eq:lD}
\end{equation}
La relation de référence d'accrétion disque--satellite \cite{Ida1997} s'écrit
\begin{equation}
\frac{M_{L,\mathrm{Ida}}}{M_D}\simeq1{,}9\,\ell_D-1{,}1-1{,}9\,\frac{M_{\mathrm{esc}}}{M_D}.
\label{eq:Ida}
\end{equation}
Cette relation, calibrée sur des disques particulaires dont la Lune se forme vers $1{,}3\,\aR$,
sert d'estimateur de premier ordre. Lorsque le disque comporte une composante fluide
intérieure à Roche, la Lune se forme plus loin, vers $2{,}15\,\aR$, et l'estimation révisée de
Salmon et Canup \cite{SalmonCanup2012} doit être employée. La valeur calculée est bornée à
l'intervalle physique $0\le M_{L,\mathrm{Ida}}\le M_D$.

\begin{lecture}
La relation d'Ida intervient après la définition du disque post-rencontre. Elle ne remplace pas
l'hydrodynamique ; elle donne un contrôle rapide de l'échelle de masse accessible à partir du
couple $(M_D,L_D)$ produit par le passage de $B$.
\end{lecture}

\begin{controle}
Ce contrôle est surtout utile comme borne supérieure sur $M_{D,0}$, et il oriente le balayage.
Pour un réservoir de $5{,}4\,M_L$ réparti selon $\Sigma\propto r^{-1}$ entre $\Req$ et $\aR$ en
rotation képlérienne, on obtient $L_{D,0}\simeq3{,}0\times10^{34}$~kg\,m$^2$\,s$^{-1}$, donc
$\ell_D\simeq0{,}89$ et
\begin{equation}
M_{L,\mathrm{Ida}}\simeq0{,}59\,M_D\simeq3{,}2\,M_L.
\label{eq:controle-Ida}
\end{equation}
L'estimateur renvoie donc une masse satellitaire supérieure à la masse observée, ce qui indique que
la branche massive du réservoir n'est pas la branche à explorer en premier. Lu dans l'autre sens, le
même estimateur désigne l'intervalle utile : $M_{D,0}\simeq2$--$3\,M_L$ ramène
$M_{L,\mathrm{Ida}}$ au voisinage de $M_{L,\mathrm{obs}}$. C'est cet intervalle qui devrait ancrer
la grille du tableau~\ref{tab:grille}, la masse manquante étant alors fournie par le canal
d'injection $M_{\oplus\to D}$ (section~\ref{sec:bilan-injection}) plutôt que par le réservoir seul.
\end{controle}

\begin{figure}[htbp]
\centering
\begin{tikzpicture}[node distance=0.55cm,
  bx/.style={draw,rounded corners=3pt,align=center,font=\small,minimum height=1.1cm,text width=2.6cm,fill=#1},
  bx/.default=white]
  \node[bx=blue!8] (a) {disque avant $B$\\$\mathcal{D}_0$};
  \node[bx=violet!10,right=of a] (b) {passage rasant\\couple + redistribution};
  \node[bx=blue!8,right=of b] (c) {disque après $B$\\$\mathcal{D}_1$};
  \node[bx=teal!8,right=of c] (d) {relaxation\\accrétion lunaire};
  \draw[-{Stealth},thick] (a) -- (b);
  \draw[-{Stealth},thick] (b) -- (c);
  \draw[-{Stealth},thick] (c) -- (d);
\end{tikzpicture}
\caption{Le passage de $B$ transforme un disque existant ; l'accrétion lunaire agit ensuite sur le
disque réorganisé.}
\label{fig:flux}
\end{figure}

% =====================================================================
\part{Géochimie : les familles de traceurs}
\label{part:geochimie}
% =====================================================================

\noindent
La géochimie appartient à la même histoire que la dynamique. Chaque élément de matière
conserve, tout au long du chemin, ses étiquettes dynamiques, thermiques et chimiques. Cette
partie pose d'abord le cadre des trois lignées matérielles, puis organise les tests par familles de
traceurs, chacune répondant à une question physique précise. L'eau, qui appartient à la même
lignée que les autres traceurs, y trouve sa place.

% ---------------------------------------------------------------------
\chapter{Trois lignées matérielles}
\label{sec:geochimie}
% ---------------------------------------------------------------------

\section{Fractions de provenance}

La matière lunaire peut contenir trois contributions :
\begin{equation}
f_0+f_\oplus+f_B=1,
\label{eq:fractions}
\end{equation}
où $f_0$ est la fraction héritée du disque préexistant, $f_\oplus$ la fraction injectée depuis la
proto-Terre pendant la rencontre et $f_B$ la contribution éventuelle de $B$. Les masses
correspondantes sont
\begin{equation}
M_L^{(k)}=f_k\,M_L,\qquad k\in\{0,\oplus,B\},
\label{eq:MLk}
\end{equation}
et les $f_k$ résultent de la décomposition~\eqref{eq:provenance} pondérée par l'efficacité
d'accrétion de chaque lignée.

\section{Mélange isotopique pondéré}

Pour un système isotopique ou un rapport $R_X$, le mélange est pondéré par les concentrations :
\begin{equation}
R_X^L=\frac{f_0C_{X,0}R_{X,0}+f_\oplus C_{X,\oplus}R_{X,\oplus}+f_BC_{X,B}R_{X,B}}
{f_0C_{X,0}+f_\oplus C_{X,\oplus}+f_BC_{X,B}}.
\label{eq:melange}
\end{equation}
Cette relation s'applique séparément aux systèmes O, Ti, Cr, Si, W et H, avec leurs concentrations
propres.

\begin{ideephys}
Le disque préexistant possède sa propre histoire chimique ; sa composition ne peut pas être
rendue invisible par définition. La théorie en tire une prédiction : la chimie lunaire résulte d'un
mélange dont les fractions sont calculées par la dynamique.
\end{ideephys}


\section{Une question physique par famille de traceurs}

Les traceurs géochimiques ne sont pas interchangeables : chacun enregistre un processus
différent. La théorie les regroupe en quatre familles, auxquelles s'ajoute l'eau (figure~\ref{fig:familles}).

\begin{figure}[htbp]
\centering
\begin{tikzpicture}[
  fam/.style={draw,rounded corners=3pt,align=center,font=\scriptsize,text width=3.35cm,
              minimum height=2.35cm,fill=#1},fam/.default=white]
  \node[fam=orange!10] (a) at (0,0) {\textbf{Ségrégation du noyau}\\[3pt]V, Cr, Mn\\W/U, Ni, Co\\[3pt]\emph{la matière vient-elle d'un manteau déjà différencié ?}};
  \node[fam=violet!8] (b) at (3.75,0) {\textbf{Identité de source}\\[3pt]$\Delta^{17}$O, $\varepsilon^{50}$Ti, $\varepsilon^{54}$Cr\\$\varepsilon^{182}$W\\[3pt]\emph{quelle est la part de $B$ ?}};
  \node[fam=teal!8] (c) at (7.5,0) {\textbf{Profondeur et fer}\\[3pt]W/V\\FeO, $\phi_{\mathrm{met}}$\\[3pt]\emph{quelle couche a été prélevée ?}};
  \node[fam=blue!6] (d) at (11.25,0) {\textbf{Volatilité et eau}\\[3pt]K/U, Zn, H$_2$O\\D/H\\[3pt]\emph{quelle histoire thermique dans $\RD$ ?}};
  \node[draw,thick,rounded corners=3pt,font=\small,fill=green!5] (L) at (5.6,-2.4) {composition lunaire $\mathbf{X}_L$};
  \foreach \n in {a,b,c,d} \draw[-{Stealth},green!40!black] (\n.south) -- (L);
\end{tikzpicture}
\caption{Les quatre familles de traceurs et la question physique à laquelle chacune répond. L'eau,
dont la signature isotopique est indiscernable de l'eau terrestre, appartient à la famille
« Volatilité ».}
\label{fig:familles}
\end{figure}

% ---------------------------------------------------------------------
\chapter{Les familles de traceurs}
\label{ch:traceurs}
% ---------------------------------------------------------------------

\section{Réfractaires sidérophiles modérés : V, Cr, Mn}

Le chrome, le vanadium et le manganèse ont des abondances semblables dans le manteau terrestre
et dans la Lune, et ils sont nettement appauvris par rapport à leurs abondances primordiales
normalisées au magnésium \cite{Ringwood1990}. Cette similitude a été utilisée pour proposer que la
Lune dérive substantiellement du manteau terrestre après la formation du noyau
\cite{Drake1989}.

Le vanadium occupe une place particulière dans cette famille. Il est moins volatil que le magnésium :
son appauvrissement ne peut pas résulter d'une perte par évaporation, ni dans la nébuleuse ni
dans le réservoir circumterrestre. Il enregistre la ségrégation métal--silicate.

Le débat mérite d'être rapporté fidèlement. Des expériences à 1~bar ont montré que les
appauvrissements de V, Cr et Mn ne suivent pas les coefficients de partage métal/silicate mesurés
dans ces conditions, ce qui a conduit à conclure qu'ils ne s'expliquaient pas d'abord par la
formation du noyau terrestre \cite{Drake1989}. Les expériences à haute pression ont ensuite
montré que, dans les conditions du manteau inférieur, le comportement de ces éléments change
avec la solubilité croissante de l'oxygène dans le fer liquide, de sorte que des motifs
d'appauvrissement semblables dans la Terre et la Lune indiquent qu'une grande part de la Lune
provient du manteau terrestre après la ségrégation du noyau \cite{Ringwood1990,GessmannRubie2000}.

\begin{lecture}
Une Lune petite, à basse pression interne, produit difficilement seule un appauvrissement en
vanadium caractéristique d'une ségrégation à haute pression. Dans la présente théorie, cette
signature s'interprète comme la trace du flux de matière silicatée terrestre vers le réservoir,
$\Phi_{\oplus\to D}$. Le vanadium n'est pas lu isolément : il est interprété conjointement avec le
tungstène et les autres traceurs de la même lignée.
\end{lecture}

\section{Le tungstène : abondance, isotopes et horloge}

Le tungstène apporte deux informations indépendantes.

\paragraph{Son abondance.}
Le rapport W/U, ou W/Th, compare des éléments d'incompatibilité voisine lors de la fusion
silicatée. Il est donc peu sensible à la différenciation magmatique et enregistre la sidérophilie du
tungstène, c'est-à-dire la ségrégation du métal \cite{Newsom1986}. Une valeur lunaire proche de la
valeur terrestre indique une source silicatée ayant subi la même histoire de ségrégation.

\paragraph{Ses isotopes.}
Après correction de l'accrétion tardive, $\varepsilon^{182}$W lunaire est pratiquement identique à
celui du manteau terrestre \cite{Kruijer2015,Touboul2015}. Dans la présente théorie, cette identité
correspond à ce qui est attendu, puisque le tungstène lunaire provient principalement de matière
silicatée terrestre ; cette attente est à vérifier dans le domaine viable.

\begin{apport}
Deux résultats se déduisent ici de la structure même du chemin proposé, sans paramètre ajusté.

D'abord la similitude isotopique Terre--Lune. Si la matière lunaire provient principalement du
manteau terrestre, par le réservoir préexistant et par le flux $\Phi_{\oplus\to D}$, et si la
contribution de $B$ reste faible, alors l'identité en O, Ti, Cr et $^{182}$W n'est pas une
coïncidence à expliquer : c'est la conséquence directe de la provenance. La théorie ne rend pas
cette identité \emph{possible}, elle la rend \emph{attendue}.

Ensuite les appauvrissements en V, Cr et Mn. Ces éléments sont modérément sidérophiles, et leur
déficit dans le manteau terrestre est hérité de la ségrégation du noyau. Si le flux
$\Phi_{\oplus\to D}$ prélève de la matière silicatée terrestre déjà appauvrie, la Lune hérite de
cette signature sans qu'il faille invoquer une seconde ségrégation lunaire. La même lignée de
matière porte donc à la fois la masse, le déficit de fer métallique et les rapports d'éléments
modérément sidérophiles. C'est cette économie d'hypothèses que l'on propose comme argument
principal en faveur du chemin étudié.
\end{apport}

\paragraph{Une horloge.}
Si la couche prélevée avait acquis un rapport Hf/W différent de celui du manteau moyen, la
décroissance du $^{182}$Hf (demi-vie de 8,9~Ma) aurait produit un excès ou un déficit de $^{182}$W,
sauf si le $^{182}$Hf était déjà éteint. Le tungstène contraint donc à la fois la source et l'instant
$t_B$.

\paragraph{Borne sur la composition du réservoir.}
Avec $f_B=0$ et la Terre silicatée comme référence ($\varepsilon^{182}\mathrm{W}_\oplus=0$),
l'équation~\eqref{eq:melange} linéarisée donne (annexe~\ref{app:borneW})
\begin{equation}
\varepsilon^{182}\mathrm{W}_L\simeq f_0\,\frac{C_{\mathrm{W},0}}{C_{\mathrm{W},\oplus}}\,
\varepsilon^{182}\mathrm{W}_0
\quad\Longrightarrow\quad
f_0\lesssim\frac{\delta\varepsilon\;C_{\mathrm{W},\oplus}}{C_{\mathrm{W},0}\,|\varepsilon^{182}\mathrm{W}_0|},
\label{eq:bornef0}
\end{equation}
où $\delta\varepsilon$ est l'écart Terre--Lune toléré ; cette borne s'applique si la matière du réservoir
est chondritique. Avec $C_{\mathrm{W},0}\approx0{,}1$--$0{,}2$~ppm,
$\varepsilon^{182}\mathrm{W}_0\approx-1{,}9$, $C_{\mathrm{W},\oplus}\approx0{,}013$~ppm et
$\delta\varepsilon\approx0{,}1$--$0{,}2$, une matière chondritique ne pourrait représenter que
$0{,}3$ à $1{,}4\,\%$ de la Lune.

\begin{lecture}
Puisque la théorie attribue au réservoir préexistant la part principale de la matière lunaire, la
borne~\eqref{eq:bornef0} impose que ce réservoir soit de nature \emph{silicatée terrestre} :
$\varepsilon^{182}\mathrm{W}_0\simeq0$, et la borne disparaît. Cela désigne la perte de masse
équatoriale $\Phi^{\mathrm{rot}}_{\oplus\to D}$ et les éjectas d'impacts ordinaires sur le manteau comme
alimentation dominante de $\RD$, la co-accrétion de matière chondritique restant minoritaire. C'est
une prédiction de la théorie : la composition du réservoir circumterrestre est terrestre.
\end{lecture}

\section{Le fer et les sidérophiles : FeO, Ni, Co et métal}

La plupart des éléments lithophiles réfractaires et faiblement sidérophiles ont des abondances
semblables dans la Lune et la Terre silicatées, mais FeO paraît nettement plus élevé dans le
manteau lunaire, ce qui constitue une contrainte importante à satisfaire \cite{Canup2023}. Le cadre des trois lignées permet de la traiter quantitativement.

\subsection{FeO et réservoir silicaté}

Pour un mélange simplifié, la teneur en FeO du disque après la rencontre est
\begin{equation}
X^{D,1}_{\mathrm{FeO}}=f_0X^{0}_{\mathrm{FeO}}+f_\oplus X^{\oplus}_{\mathrm{FeO}}+f_BX^{B}_{\mathrm{FeO}},
\label{eq:FeO}
\end{equation}
puis corrigée des fractionnements dus à la fusion, à la condensation et aux pertes de vapeur. Si
une fraction fondue $F$ est extraite d'un réservoir avec un coefficient de partage global $D_X$,
une loi de fusion fractionnée donne par exemple
\begin{equation}
C_X^{\mathrm{liq}}=C_{X,0}\,\frac{1-(1-F)^{1/D_X}}{F}.
\label{eq:fusion}
\end{equation}
Le calcul ne choisit pas séparément une profondeur pour FeO, une autre pour W et une troisième
pour les isotopes : la même lignée matérielle porte toutes les propriétés.

\subsection{Métal et silicate}

Le disque se décompose en
\begin{equation}
M_D=M_D^{\mathrm{sil}}+M_D^{\mathrm{met}}+M_D^{\mathrm{vap}},
\qquad
\phi_{\mathrm{met}}=\frac{M_D^{\mathrm{met}}}{M_D^{\mathrm{sil}}+M_D^{\mathrm{met}}}.
\label{eq:metal}
\end{equation}
Les éléments sidérophiles Ni, Co et W contraignent l'origine, la ségrégation et la réaccrétion
du métal.


\begin{consequence}
L'excès de FeO lunaire doit être reproduit par le calcul du mélange et des transformations
internes au réservoir ($\boldsymbol{\Pi}_D$ dans l'équation~\eqref{eq:flux-X}) : oxydation d'une partie
du métal, fractionnement lors de la condensation, ou composante silicatée plus oxydée. Comme la
borne~\eqref{eq:bornef0} limite la part chondritique, l'excès doit provenir de processus internes au
réservoir.
\end{consequence}

\section{Isotopes O, Ti, Cr : l'empreinte de \texorpdfstring{$B$}{B}}
\label{sec:bornes}

Les systèmes isotopiques de l'oxygène, du titane et du chrome distinguent les corps du Système
solaire selon leur réservoir de formation. Ils mesurent donc la part de $B$ dans la Lune. Pour
l'oxygène, à concentrations comparables,
\begin{equation}
f_B\lesssim\frac{\delta_{17}}{|\Delta^{17}\mathrm{O}_B-\Delta^{17}\mathrm{O}_\oplus|},
\label{eq:bornefB}
\end{equation}
où $\delta_{17}$, de l'ordre de quelques ppm, est l'écart Terre--Lune mesuré \cite{Young2016}. Pour un
corps $B$ de composition isotopique de type martien
($|\Delta^{17}\mathrm{O}_B-\Delta^{17}\mathrm{O}_\oplus|\approx300$~ppm), $f_B\lesssim2\,\%$. Les
contraintes du titane \cite{Zhang2012} et du chrome s'expriment de la même manière.

\begin{apport}
Le mécanisme favorise naturellement une faible contribution matérielle de $B$, car le passage est
non collisionnel. La similitude isotopique Terre--Lune devient alors une contrainte attendue du
domaine viable, à vérifier par les bornes sur $f_B$ (équation~\eqref{eq:bornefB}).
\end{apport}

\section{Volatils et eau}

Pour une espèce volatile $i$, le bilan et le fractionnement isotopique associé à une perte de
vapeur, de type Rayleigh, s'écrivent
\begin{equation}
\frac{\dd M_i}{\dd t}=\dot M_{i,\mathrm{in}}-\dot M_{i,\oplus}-\dot M_{i,\mathrm{esc}}
+\dot M_{i,\mathrm{cond}}-\dot M_{i,\mathrm{vap}},
\qquad
R_i=R_{i,0}\,f_i^{(\alpha_i-1)},
\label{eq:volatils}
\end{equation}
où $f_i$ est la fraction résiduelle et $\alpha_i$ le facteur de fractionnement effectif.

\paragraph{L'eau lunaire appartient à la même lignée.}
Il serait tentant de conclure que la Lune naît sèche, le disque à la limite de Roche étant porté
vers $2000$~K, bien au-dessus de la condensation de l'eau. Les mesures interdisent cette
conclusion. La teneur en eau pré-éruptive des verres volcaniques lunaires est estimée à plusieurs
centaines de parties par million \cite{Saal2008}, valeur comparable à celle du manteau supérieur
terrestre, et la composition isotopique de l'hydrogène lunaire s'est révélée indiscernable de celle
de l'eau terrestre \cite{Saal2013}.

Cette observation renforce la théorie plutôt qu'elle ne la gêne, et pour la raison qui vaut déjà
pour l'oxygène, le titane, le chrome et le tungstène : \emph{c'est la même lignée de matière}. La
Lune contient une eau isotopiquement terrestre parce qu'elle est faite de manteau terrestre. Le
déficit relatif s'explique par la volatilité dans un disque chaud, la parenté isotopique par la
provenance. L'hydrogène rejoint ainsi la liste des traceurs qui pointent vers une origine
terrestre commune.

\begin{controle}
Deux points restent à établir, et ils sont contraignants. D'abord la chronologie : une eau arrivée
pendant les quelques heures de la rencontre aurait peine à se mélanger au manteau avant le départ
de la matière lunaire ; la livraison étalée décrite à la section~\ref{sec:trainee} lève cette
difficulté, mais l'ordre des événements doit être explicité. Ensuite la composition : si $B$
apporte une part de l'eau terrestre, son rapport deutérium sur hydrogène devient une contrainte
forte, puisque l'eau terrestre correspond aux chondrites carbonées et non aux comètes. La région de
formation de $B$ s'en trouverait contrainte, ce qui constitue un test du présent scénario.
\end{controle}

\section{Le rapport W/V et la profondeur d'origine}

Le tungstène et le vanadium ne constituent pas des preuves isolées : ce sont des \emph{traceurs
couplés} des flux entre réservoirs et de la profondeur de la matière prélevée. Leur interprétation
passe par l'équation de composition~\eqref{eq:flux-X} et par la distribution de la matière
injectée~\eqref{eq:flux}.


Le tungstène est très incompatible : lors de la cristallisation d'un océan magmatique, il se
concentre dans les derniers liquides, donc vers la surface. Le vanadium est modérément
compatible, retenu par les pyroxènes et les spinelles. Le rapport W/V, normalisé à des éléments de
référence de même comportement, varie donc avec la profondeur et l'état de différenciation de la
couche prélevée. Il s'écrit, pour la matière injectée,
\begin{equation}
\left(\frac{\mathrm{W}}{\mathrm{V}}\right)_{\oplus\to D}=
\frac{\displaystyle\iiint\frac{\dd M_{\oplus\to D}}{\dd\lambda_\Omega\,\dd z\,\dd t}\,
C_{\mathrm{W}}(\lambda_\Omega,z)\,\dd\lambda_\Omega\,\dd z\,\dd t}
{\displaystyle\iiint\frac{\dd M_{\oplus\to D}}{\dd\lambda_\Omega\,\dd z\,\dd t}\,
C_{\mathrm{V}}(\lambda_\Omega,z)\,\dd\lambda_\Omega\,\dd z\,\dd t}.
\label{eq:WV}
\end{equation}

\begin{ideephys}
C'est l'expression la plus directe de la lignée matérielle unique : la même distribution de
profondeur $z$ et de latitude $\lambda_\Omega$, produite par la dynamique du passage
(équation~\eqref{eq:flux}), fixe simultanément la masse injectée et ses rapports élémentaires. On ne
choisit pas une profondeur pour chaque traceur.
\end{ideephys}

\section{Synthèse des traceurs}

\begin{table}[htbp]
\centering
\small
\caption{Familles de traceurs, processus enregistrés et prédictions de la théorie.}
\label{tab:traceurs}
\begin{tabularx}{\textwidth}{>{\raggedright\arraybackslash}p{2.6cm}
>{\raggedright\arraybackslash}X >{\raggedright\arraybackslash}X}
\toprule
\textbf{Traceur} & \textbf{Ce qu'il enregistre} & \textbf{Prédiction de la théorie} \\
\midrule
V, Cr, Mn & ségrégation du noyau à haute pression & appauvrissements lunaires semblables aux terrestres \\
W/U, W/Th & sidérophilie du tungstène & rapport lunaire proche du rapport terrestre \\
$\varepsilon^{182}$W & identité de source, chronologie & proche du manteau terrestre ; contrainte sur $t_B$ \\
$\Delta^{17}$O, $\varepsilon^{50}$Ti, $\varepsilon^{54}$Cr & part de $B$ & $f_B\ll1$ \\
FeO, $\phi_{\mathrm{met}}$ & fer oxydé et métal & excès lunaire produit dans le réservoir \\
W/V & profondeur et différenciation de la couche prélevée & lié à la latitude et à la profondeur arrachées \\
K/U, Zn & histoire thermique du réservoir & appauvrissement produit dans $\RD$ \\
H$_2$O, D/H & volatilité et source de l'eau & eau isotopiquement terrestre ; livrée par $B$ et par la traînée \\
\bottomrule
\end{tabularx}
\end{table}
% =====================================================================
\part{Domaine de viabilité}
\label{part:viabilite}
% =====================================================================

\noindent
Un chemin physique n'est crédible que s'il existe dans une région continue de paramètres, et non
pour une configuration isolée finement ajustée. Cette partie rassemble les paramètres du
Niveau~I, les portes que le chemin doit franchir simultanément et le programme numérique qui
produira les cartes de viabilité.

% ---------------------------------------------------------------------
\chapter{Paramètres, portes et domaine viable}
\label{ch:portes}
% ---------------------------------------------------------------------

\section{Les paramètres du Niveau I}

Le vecteur $\Theta_I$ (équation~\eqref{eq:Theta}) regroupe les paramètres de la proto-Terre, du
réservoir et de la rencontre. Le tableau~\ref{tab:grille} propose une grille initiale d'exploration ;
la bande dynamique de basses latitudes prograde y est désignée comme zone prioritaire.

\begin{table}[htbp]
\centering
\small
\caption{Grille initiale d'exploration des paramètres du Niveau I.}
\label{tab:grille}
\begin{tabular}{lll}
\toprule
\textbf{Paramètre} & \textbf{Domaine initial} & \textbf{Remarque} \\
\midrule
$\mu_B=M_B/\Mt$ & $0{,}05$--$1$ & masse de $B$ \\
$q_B$ & $\chi_{\mathrm{tide}}\,d_{R,B}$--$10\,\Rt$ & borne basse : porte~\eqref{eq:DtideB}, soit $\approx3{,}0\,\Rt$ \\
$v_\infty$ & $0$--$10$ km\,s$^{-1}$ & vitesse asymptotique relative \\
$t_{\mathrm{ej}}$ & $>t_B$ & élimination ou éjection ultérieure de $B$ \\
$i_B$ & $0$--$180^\circ$ & priorité : $|\lambda_B|\le\lambda_c$, prograde \\
$\omega_B$ & $0$--$180^\circ$ & priorité : voisinage de $90^\circ$ \\
$\tilde\omega$ & $0{,}8$--$1$ & rotation réduite de la proto-Terre \\
$\Req/\Rt$ & $1$--$2$ & dilatation ; configuration de référence à $1{,}7$ \\
Réservoir & $2$--$5\,\Rt$, $\Sigma\propto r^{-p}$, $p\in[0,3]$ & structure radiale \\
$M_{D,0}$ & $1$--$8\,M_L$ & masse initiale du réservoir ; priorité $2$--$3\,M_L$ \\
$\Omega_\oplus$ & $\tilde\omega\in[0{,}8,1]$ & contrôle $M_{\oplus\to D}$ (section~\ref{sec:bilan-injection}) \\
\bottomrule
\end{tabular}
\end{table}

\section{Les portes du domaine viable}

Le domaine total est défini par l'intersection
\begin{equation}
\boxed{\begin{aligned}
\mathcal{D}_{\mathrm{viable}}=\ &\mathcal{D}_{\mathrm{disque},0}\cap\mathcal{D}_{\mathrm{thermo}}
\cap\mathcal{D}_{\mathrm{grav}}\cap\mathcal{D}_{\mathrm{nc}}\cap\mathcal{D}_{\mathrm{nt}}
\cap\mathcal{D}_{B}\cap\mathcal{D}_{B,\mathrm{hel}}\cap\mathcal{D}_{\mathrm{tide},B}\\
&\cap\mathcal{D}_{\mathrm{fen}}\cap\mathcal{D}_{\mathrm{imp}}
\cap\mathcal{D}_{J}^{\mathrm{num}}\cap\mathcal{D}_{J}^{\mathrm{EM}}\cap\mathcal{D}_{\mathrm{evac}}
\cap\mathcal{D}_{\mathrm{accr}}\cap\mathcal{D}_{\mathrm{arch}}
\cap\mathcal{D}_{\mathrm{geo}}\cap\mathcal{D}_{\mathrm{iso}}\cap\mathcal{D}_{\mathrm{chrono}}.
\end{aligned}}
\label{eq:Dviable}
\end{equation}
Les portes ont les significations suivantes :
\begin{align}
\mathcal{D}_{\mathrm{disque},0}&=\left\{\Theta_I:\ \text{conditions~\eqref{eq:persistance} et~\eqref{eq:tpers} vérifiées à }t_B\right\},
\label{eq:Ddisque}\\
\mathcal{D}_{\mathrm{thermo}}&=\left\{\Theta_I:\ f_{\mathrm{sol}}(r,t)>0\ \text{pour }r>\aR\ \text{avant }
t^{\mathrm{start}}_{\mathrm{acc},L}\right\},
\label{eq:Dthermo}\\
\mathcal{D}_{\mathrm{grav}}&=\left\{\Theta_I:\ Q_T(r,t)\ \text{et}\ t_{\mathrm{sat}}\ \text{compatibles avec le régime de la section~\ref{sec:regimes}}\right\},
\label{eq:Dgrav}\\
\mathcal{D}_{\mathrm{nc}}&=\left\{\Theta_I:\ q_B>\Req+R_{B,\mathrm{eff}}\right\},
\label{eq:Dnc}\\
\mathcal{D}_{\mathrm{tide},B}&=\left\{\Theta_I:\ q_B>\chi_{\mathrm{tide}}\,d_{R,B},\ \chi_{\mathrm{tide}}>1\right\},
\label{eq:DtideB}\\
\mathcal{D}_{\mathrm{nt}}&=\left\{\Theta_I:\ R_{\mathrm{out,dense}}<\min(r^{+}_{\text{nœud}},r^{-}_{\text{nœud}})\right\},
\label{eq:Dnt}\\
\mathcal{D}_{B}&=\left\{\Theta_I:\ f_B=\frac{M_{B\to L}}{M_L}\le f_{B,\max}\right\},
\label{eq:DB}\\
\mathcal{D}_{B,\mathrm{hel}}&=\left\{\Theta_I^{+}:\ t_B<t_{\mathrm{ej}}
\text{ et }\mathcal{H}_{B,1}\to B\in\mathcal{P}_{\mathrm{elim}}\right\},
\label{eq:DBhel}\\
\mathcal{D}_{\mathrm{fen}}&=\left\{\Theta_I:\ \Delta\eta_L\ge\eta_{L,\min},\ \eta_{\mathrm{esc}}\le\eta_{\mathrm{esc},\max}\right\},
\label{eq:Dfen}\\
\mathcal{D}_{\mathrm{imp}}&=\left\{\Theta_I:\ \delta_{\mathrm{imp}}(\Theta_I)<\epsilon_{\mathrm{imp}}\right\},
\label{eq:Dimp}\\
\mathcal{D}_{J}^{\mathrm{num}}&=\left\{\Theta_I:\ \frac{|J_{\mathrm{avant}}-J_{\text{après}}|}{|J_{\mathrm{avant}}|}<\epsilon_J^{\mathrm{num}}\right\},
\label{eq:DJ}\\
\mathcal{D}_{J}^{\mathrm{EM}}&=\left\{\Theta_I:\ \frac{|J_{\oplus D,1}-J_{\mathrm{loss}}-L_{\mathrm{EM}}|}{L_{\mathrm{EM}}}<\epsilon_J\right\},
\label{eq:DJEM}\\
\mathcal{D}_{\mathrm{accr}}&=\left\{\Theta_I:\ M_{L,\mathrm{calc}}\in[M_{L,\min},M_{L,\max}]\right\},
\label{eq:Daccr}\\
\mathcal{D}_{\mathrm{geo}}&=\left\{\Theta_I:\ X^L_{\mathrm{FeO}},\ \phi_{\mathrm{met},L}\ \text{et sidérophiles dans les intervalles observés}\right\},
\label{eq:Dgeo}\\
\mathcal{D}_{\mathrm{iso}}&=\left\{\Theta_I:\ |R_X^L-R_X^\oplus|<\delta_X,\ X\in\{\mathrm{O},\mathrm{Ti},\mathrm{Cr},\mathrm{Si},\mathrm{W}\}\right\},
\label{eq:Diso}\\
\mathcal{D}_{\mathrm{chrono}}&=\left\{\Theta_I:\ \text{ordre~\eqref{eq:ordre} respecté et }
t^{\mathrm{accr}}_L\ \text{compatible avec l'âge de la Lune}\right\},
\label{eq:Dchrono}
\end{align}
et $\mathcal{D}_{\mathrm{arch}}$ est donné par l'équation~\eqref{eq:Darch}.

Une porte supplémentaire s'ajoute pour l'eau :
\begin{equation}
\mathcal{D}_{\mathrm{eau}}=\left\{\Theta_I:\ \text{H}_2\text{O lunaire et D/H compatibles avec la lignée terrestre}\right\}.
\label{eq:Deau}
\end{equation}


\paragraph{Vue d'ensemble.} Les portes sont nombreuses, et elles ont été introduites au fil des
chapitres, chacune à l'endroit où la physique la rendait nécessaire. Il est utile de les
rassembler. Le tableau~\ref{tab:portes} les récapitule, avec ce que chacune exige et l'état
d'avancement de sa vérification. La colonne de droite est la plus instructive : elle montre que la
théorie n'est pas également avancée partout, et où le travail reste à faire.

\begin{table}[htbp]
\centering
\small
\caption{Récapitulatif des portes du domaine viable. Une porte franchie est une contrainte
satisfaite ; une porte fermée exclut la branche correspondante. La colonne d'état indique si la
vérification a été conduite, esquissée, ou reste entièrement à faire.}
\label{tab:portes}
\begin{tabularx}{\textwidth}{>{\raggedright\arraybackslash}p{2.6cm}
>{\raggedright\arraybackslash}X >{\raggedright\arraybackslash}p{2.6cm}}
\toprule
\textbf{Porte} & \textbf{Ce qu'elle exige} & \textbf{État} \\
\midrule
$\mathcal{D}_{\mathrm{disque},0}$ & le réservoir existe et survit jusqu'au passage & à chiffrer \\
$\mathcal{D}_{\mathrm{thermo}}$ & l'état thermique du réservoir est cohérent & à chiffrer \\
$\mathcal{D}_{\mathrm{grav}}$ & le réservoir reste gravitationnellement admissible & esquissée \\
\midrule
$\mathcal{D}_{\mathrm{nc}}$ & $B$ ne touche pas la proto-Terre & franchie \\
$\mathcal{D}_{\mathrm{tide},B}$ & $B$ ne se disloque pas par marée & franchie \\
$\mathcal{D}_{\mathrm{nt}}$ & le passage reste non transitoire au sens requis & esquissée \\
$\mathcal{D}_{B}$ & la masse et la trajectoire de $B$ sont plausibles & esquissée \\
$\mathcal{D}_{B,\mathrm{hel}}$ & le devenir héliocentrique de $B$ est admissible & esquissée \\
\midrule
$\mathcal{D}_{\mathrm{fen}}$ & une fraction utile atteint la fenêtre lunaire & franchie \\
$\mathcal{D}_{\mathrm{imp}}$ & la forme close de l'impulsion reste valable & à chiffrer \\
$\mathcal{D}_{J}^{\mathrm{num}}$ & l'intégrateur conserve le moment cinétique & à chiffrer \\
$\mathcal{D}_{J}^{\mathrm{EM}}$ & l'état final rejoint $L_{\mathrm{EM}}$ & voie ouverte \\
$\mathcal{D}_{\mathrm{evac}}$ & $\phi_{\oplus\to B}\ge0{,}20$ part avec $B$ & à chiffrer \\
\midrule
$\mathcal{D}_{\mathrm{accr}}$ & la masse lunaire calculée est la bonne & franchie \\
$\mathcal{D}_{\mathrm{arch}}$ & un satellite dominant émerge & esquissée \\
$\mathcal{D}_{\mathrm{geo}}$ & fer, métal et sidérophiles aux teneurs observées & esquissée \\
$\mathcal{D}_{\mathrm{iso}}$ & O, Ti, Cr, Si et W indiscernables du manteau & esquissée \\
$\mathcal{D}_{\mathrm{eau}}$ & l'eau lunaire et le rapport D/H sont terrestres & esquissée \\
$\mathcal{D}_{\mathrm{chrono}}$ & l'ordre des événements et l'âge sont respectés & esquissée \\
\bottomrule
\end{tabularx}
\end{table}

Trois portes sont aujourd'hui franchies sur des bases quantitatives : la non-collision,
l'intégrité de $B$, et l'accrétion d'une masse lunaire. Cinq demandent un calcul qui n'a pas
encore été mené. Les autres reposent sur des arguments établis mais non chiffrés. C'est une carte
honnête de l'état d'avancement, et elle vaut mieux qu'une affirmation globale de cohérence.

\paragraph{La porte $\mathcal{D}_{\mathrm{imp}}$.}
Cette porte contrôle que l'approximation rectiligne utilisée pour la forme close de l'impulsion
reste acceptable dans le domaine viable. Elle est franchie lorsque l'écart entre la
forme~\eqref{eq:impulsion-close} et l'intégration directe sur la conique hyperbolique reste
inférieur à un seuil fixé, par exemple $10\,\%$ sur la fraction $\eta_L$. Ce seuil est un choix
de modélisation, à justifier en fonction de la précision recherchée sur le bilan d'injection. Si la
porte n'est pas franchie, la branche correspondante du domaine viable est exclue, et le calcul doit
être repris avec la forme close hyperbolique exacte (intégration numérique directe).

\paragraph{La porte $\mathcal{D}_{\mathrm{eau}}$.}
Cette porte vérifie la cohérence de l'apport d'eau avec la lignée terrestre. Elle exige que la
teneur en eau pré-éruptive des échantillons lunaires reste dans l'intervalle observé (quelques
centaines de ppm) et que le rapport D/H lunaire reste indiscernable du rapport terrestre
\cite{Saal2008,Saal2013}. Si, dans le domaine viable, l'apport d'eau par $B$ impliquait un rapport
D/H mesurablement non terrestre, la porte se fermerait sur la branche correspondante — ce qui
contraindrait la région de formation de $B$.

Trois précisions sur ces portes. La première distingue deux exigences qui étaient auparavant
confondues : $\mathcal{D}_J^{\mathrm{num}}$ vérifie que l'intégrateur conserve le moment cinétique
total, ce qui valide le calcul ; $\mathcal{D}_J^{\mathrm{EM}}$ vérifie que l'état laissé par le
passage rejoint $L_{\mathrm{EM}}$ au terme de l'évolution. Pour que la seconde soit falsifiable,
$J_{\mathrm{loss}}$ doit être une \emph{sortie} de la Phase~C, produite par le mécanisme
d'évacuation intégré, et non un paramètre ajusté a posteriori.

La deuxième porte sur l'intégrité de $B$. Les deux conditions doivent être évaluées avec les
paramètres de la configuration de référence. Pour $\mu_B=0{,}95$, le rayon propre de $B$ vaut
$R_{B,\mathrm{eff}}\simeq0{,}99\,\Rt$, de sorte que la condition de non-contact s'écrit
$q_B>\Req+R_{B,\mathrm{eff}}\simeq2{,}69\,\Rt$, tandis que la dislocation de marée impose
$d_{R,B}\simeq2{,}75\,\Rt$ d'après l'équation~\eqref{eq:RocheB}. Les deux seuils sont donc
désormais voisins, et la seconde reste la plus contraignante.

Le périastre nominal $q_B=3{,}0\,\Rt$ satisfait l'une et l'autre, avec une marge de $12\,\%$ sur le
contact et de $9\,\%$ sur la dislocation, c'est-à-dire $\chi_{\mathrm{tide}}\simeq1{,}09$. La marge
adoptée dans le domaine viable est $\chi_{\mathrm{tide}}=1{,}05$, que la configuration nominale
franchit confortablement ; une exigence plus sévère, $\chi_{\mathrm{tide}}=1{,}2$, imposerait
$q_B>3{,}30\,\Rt$ et placerait la configuration nominale hors du domaine. La borne basse du
tableau~\ref{tab:grille} doit donc être choisie en cohérence avec la valeur retenue pour
$\chi_{\mathrm{tide}}$.

La troisième concerne $\mathcal{D}_{\mathrm{accr}}$ : la masse à comparer est
$M_{L,\mathrm{calc}}$ au sens de l'équation~\eqref{eq:MLcalc-eff}, donc somme du terme de tri et du
terme d'injection, et non le seul produit $\eta_L M_{D,0}$.

\begin{falsif}
La théorie demande une région continue de paramètres satisfaisant toutes les portes ensemble.
Un cas isolé extrêmement ajusté ne constitue pas un domaine de plausibilité robuste.
\end{falsif}

% =====================================================================

\begin{figure}[htbp]
\centering
\begin{tikzpicture}[y=0.62cm]
  \foreach \w/\y/\lab/\c in {
     12/0/{$\Theta_I$ : ensemble des paramètres}/gray!10,
     10.5/-1.25/{$\mathcal{D}_{\mathrm{disque},0}$, $\mathcal{D}_{\mathrm{thermo}}$, $\mathcal{D}_{\mathrm{grav}}$ : réservoir persistant}/blue!8,
     9/-2.5/{$\mathcal{D}_{\mathrm{nc}}$, $\mathcal{D}_{\mathrm{nt}}$, $\mathcal{D}_B$ : passage non collisionnel}/violet!8,
     7.5/-3.75/{$\mathcal{D}_{\mathrm{fen}}$, $\mathcal{D}_J$ : tri orbital et bilans}/violet!14,
     6/-5/{$\mathcal{D}_{\mathrm{accr}}$, $\mathcal{D}_{\mathrm{arch}}$ : une Lune dominante}/teal!10,
     4.5/-6.25/{$\mathcal{D}_{\mathrm{geo}}$, $\mathcal{D}_{\mathrm{iso}}$, $\mathcal{D}_{\mathrm{eau}}$, $\mathcal{D}_{\mathrm{chrono}}$}/green!10,
     3/-7.5/{$\mathcal{D}_{\mathrm{viable}}$}/orange!25}
  { \node[draw,rounded corners=2pt,fill=\c,minimum width=\w cm,minimum height=0.62cm,font=\scriptsize]
      at (0,\y) {\lab}; }
\end{tikzpicture}
\caption{Le domaine viable comme intersection de portes successives. Seules les configurations
qui franchissent toutes les portes définissent le chemin physique.}
\label{fig:portes}
\end{figure}

% ---------------------------------------------------------------------
\chapter{Programme numérique de Niveau I}
\label{ch:programme}
% ---------------------------------------------------------------------

La séparation des échelles de temps (chapitre~\ref{ch:chronologie}) permet de découper le calcul en
trois phases enchaînées, au cours desquelles chaque élément de matière conserve ses étiquettes
dynamiques, thermiques et chimiques.

\section{Phase A : formation et évolution du disque avant \texorpdfstring{$B$}{B}}

La première simulation calcule
\begin{equation}
\mathcal{A}(t)=\left\{\dot M_{\mathrm{in}},\ \Sigma(r,t),\ T(r,t),\ P(r,t),\ L_D(t),\
R_{\mathrm{out}}(t),\ \mathbf{X}(r,t),\ \mathbf{f}_{\mathrm{phase}}(r,t),\ Q_T(r,t),\ \mathcal{G}(t)\right\}.
\label{eq:phaseA}
\end{equation}
Sa sortie est une famille d'états $\mathcal{D}_0$ accessibles au temps $t_B$.

\section{Phase B : passage rasant de \texorpdfstring{$B$}{B}}

La rencontre, intégrée avec le champ exact~\eqref{eq:adiff} et fermée dans le repère
barycentrique, fournit
\begin{equation}
\mathcal{B}=\left\{M_{D,1},\ L_{D,1},\ \Delta\mathbf{S}_\oplus,\ \Delta\mathbf{L}_{\mathrm{orb}},\
\Delta\mathbf{S}_B,\ M_{\oplus\to D},\ M_{B\to D},\ M_{\mathrm{esc}},\ i_{D,1},\ e_D(r),\ \chi_J,\ \mathcal{H}_{B,1}\right\}.
\label{eq:phaseB}
\end{equation}
La distribution des particules ou éléments de fluide est résumée par
\begin{equation}
\Pi_{\mathrm{dyn}}=\Pi_{\mathrm{dyn}}\left(\varepsilon_{\mathrm{orb}},j_z,q_p,i,r_0,\lambda_\Omega,z
\,\middle|\,\Theta_I\right).
\label{eq:Pidyn}
\end{equation}

\section{Phase C : relaxation et accrétion}

La phase finale calcule
\begin{equation}
\mathcal{C}=\left\{M_L,\ L_L,\ i_{L,0},\ \mathbf{X}_L,\ t_{\mathrm{acc},L},\ N_{\mathrm{sat}},\
f_{\mathrm{dom}},\ M_{\mathrm{reaccr}},\ M_{\mathrm{esc}}\right\}.
\label{eq:phaseC}
\end{equation}

\begin{figure}[htbp]
\centering
\begin{tikzpicture}[node distance=0.8cm,
  ph/.style={draw,rounded corners=3pt,align=center,font=\small,minimum height=1.3cm,text width=3.6cm,fill=#1}]
  \node[ph=blue!8] (A) {\textbf{Phase A}\\formation et persistance\\du disque};
  \node[ph=violet!10,right=of A] (B) {\textbf{Phase B}\\passage rasant\\non collisionnel de $B$};
  \node[ph=teal!8,right=of B] (C) {\textbf{Phase C}\\relaxation et\\accrétion lunaire};
  \draw[-{Stealth},thick] (A) -- (B);
  \draw[-{Stealth},thick] (B) -- (C);
  \node[font=\footnotesize,align=center,below=0.45cm of B,text width=13cm]
    {Les mêmes éléments de matière conservent leurs étiquettes dynamiques, thermiques et
     chimiques tout au long de la chaîne.};
\end{tikzpicture}
\caption{Architecture du programme numérique de la théorie.}
\label{fig:programme}
\end{figure}

% =====================================================================

\section{Protocole de la Phase B et cartes de viabilité}

\paragraph{Modèle.}
Les éléments du réservoir sont d'abord traités comme des particules test dans le champ de la
proto-Terre, $B$ suivant son hyperbole exacte. L'intégration utilise l'accélération
différentielle~\eqref{eq:adiff} et s'arrête lorsque le critère~\eqref{eq:atol} est satisfait. Chaque
élément est ensuite classé selon $(\varepsilon,\,j_z,\,e,\,q_p,\,r_{\mathrm{circ}})$ et le
critère~\eqref{eq:fenetre}.

\paragraph{Validation croisée de l'impulsion.}
La Phase~B calcule $\Delta\mathbf{v}$ par intégration directe de l'accélération
différentielle~\eqref{eq:adiff} sur la vraie conique hyperbolique de $B$. Le résultat est comparé à
la forme close rectiligne de l'annexe~\ref{app:impulsion}, et l'écart est quantifié en fonction de
$e_B$ et de $r/q_B$. Cet écart est une \emph{sortie} de la Phase~B, et non un paramètre ajusté. Il
alimente la porte $\mathcal{D}_{\mathrm{imp}}$.

\paragraph{Noyau de réponse.}
Pour découpler la réponse du profil de densité, on lance des anneaux d'éléments à une trentaine
de rayons initiaux $r_0$ entre $2$ et $5\,\Rt$, avec plusieurs centaines d'azimuts chacun. On obtient
la probabilité $P_k(r_0)$ qu'un élément parti de $r_0$ finisse dans la catégorie
$k\in\{L,\mathrm{in},\mathrm{reaccr},\mathrm{esc},\to B\}$. Pour tout profil $\Sigma(r_0)$,
\begin{equation}
\eta_k=\frac{1}{M_{D,0}}\int_{R_{\mathrm{in}}}^{R_{\mathrm{out}}}\Sigma(r_0)\,P_k(r_0)\,2\pi r_0\,\dd r_0.
\label{eq:noyau}
\end{equation}
Une seule intégration par jeu de paramètres de $B$ sert ainsi pour tous les profils du réservoir.

\paragraph{Variables réduites.}
La réponse dépend de $\mu_B$, $q_B/\Rt$, $v_\infty/v_{\mathrm{esc}}$, $i_B$ et $\omega_B$. Le moment
cinétique cédé par le réservoir est obtenu en pondérant les éléments ; la
fermeture~\eqref{eq:fermeture} en attribue l'opposé à l'orbite de $B$, ce qui reste valable tant que
$M_D\ll M_B$.

\paragraph{Cartes à produire.}
Les cartes de viabilité sont les sorties principales du programme. Elles comprennent les
fractions $\Delta\eta_L(q_B,i_B)$, $\eta_{\mathrm{esc}}(q_B,i_B)$ et $\eta_{\mathrm{reaccr}}(q_B,i_B)$ pour
plusieurs valeurs de $\mu_B$ et de $v_\infty$ ; la fraction échangée $\chi_J$ ; la distribution
de la matière injectée en latitude et en profondeur, $\dd M_{\oplus\to D}/\dd\lambda_\Omega\,\dd z$ ; et l'inclinaison $i_{D,1}$ du
réservoir réorganisé. Ces cartes sont les résultats à calculer ; la figure~\ref{fig:tri} en donne
l'anticipation analytique.
La même phase doit conserver l'orbite héliocentrique post-passage $\mathcal{H}_{B,1}$ afin de
permettre, dans un calcul ultérieur, le test de la porte $\mathcal{D}_{B,\mathrm{hel}}$.

\section{Bilan d'injection : évaluation de \texorpdfstring{$M_{\oplus\to D}$}{M_Terre-vers-D}}
\label{sec:bilan-injection}

Le canal d'injection introduit au chapitre~\ref{ch:injection} apporte au disque de la matière
superficielle terrestre, déjà appauvrie en fer métallique, en nickel, en cobalt et en éléments
sidérophiles. C'est le seul terme du bilan~\eqref{eq:MLcalc-eff} qui soit une \emph{source} : il ne
prélève pas sur $M_{D,0}$ et n'est donc pas borné par $\eta_L$. Il porte à la fois la masse et la
signature géochimique, ce qui en fait la quantité la plus rentable à chiffrer. Le protocole suivant
le calcule en quatre étapes.

\paragraph{Étape 1 — critère de départ.} Un élément de surface situé à la colatitude $\theta$ quitte
la proto-Terre lorsque la somme de l'accélération centrifuge associée à $\Omega_\oplus$ et de la
composante radiale du champ différentiel~\eqref{eq:adiff} dépasse la gravité locale. Le seuil~\eqref{eq:seuil} donne cette condition au point sous-$B$ ; elle doit être évaluée sur toute la
surface et à chaque instant, car la zone active se déplace pendant le passage.

\paragraph{Étape 2 — moment cinétique emporté.} Chaque élément part avec
\begin{equation}
j_0=\Omega_\oplus\,\Req^2\sin^2\theta+\Delta j_B,
\label{eq:j0-injection}
\end{equation}
où $\Delta j_B$ est l'apport tangentiel de $B$. Le critère de la fenêtre~\eqref{eq:fenetre}
s'applique ensuite sans modification : $r_{\mathrm{circ}}=j_0^2/G\Mt\ge\aR$. Cette étape est
décisive, car elle décide si la matière arrachée atteint la région d'accrétion ou retombe. Le
rapport à franchir vaut
\begin{equation}
\frac{j(\aR)}{j(\Req)}=\sqrt{\frac{\aR}{\Req}}\simeq1{,}55
\qquad(\Req=1{,}20\,\Rt,\ \aR=2{,}9\,\Rt),
\label{eq:saut-j}
\end{equation}
alors que l'impulsion de marée~\eqref{eq:impulsion} ne fournit qu'environ
$0{,}04\,v_K$ à $r=\Req$ pour $\mu_B=0{,}1$ et $q_B=3\,\Rt$. La rotation propre doit donc fournir
l'essentiel de $j_0$, et $B$ franchir le dernier seuil : c'est pourquoi $\tilde\omega$ est le
paramètre dominant de ce bilan, et pourquoi $M_{\oplus\to D}$ y est très sensible.

\paragraph{Étape 3 — intégration.} La surface est pavée en $(\theta,\varphi)$ et en profondeur, chaque
élément recevant ses étiquettes $(\lambda_\Omega,z)$. Les éléments franchissant le seuil sont
intégrés dans le champ combiné de la proto-Terre aplatie et de $B$, le long de la trajectoire
hyperbolique, jusqu'à la fin de la rencontre.

\paragraph{Étape 4 — sorties.} Le calcul produit $M_{\oplus\to D}$ et sa répartition entre la part
au-delà de Roche, $M_{\oplus\to L}$, et la part intérieure à Roche, qui alimente
$\eta_{\mathrm{in}}$ ; la distribution en $r_{\mathrm{circ}}$ ; le moment cinétique emporté, à
porter au bilan~\eqref{eq:fermeture} ; et la distribution en profondeur et en latitude, qui alimente
directement le rapport W/V de l'équation~\eqref{eq:WV}.

\paragraph{Où atterrit la matière injectée.} Une remarque de cadrage, qui change la lecture des
résultats attendus. Le rayon de circularisation d'un élément quittant la surface en corotation
vaut, à l'ordre dominant,
\begin{equation}
r_{\mathrm{circ}}\simeq\tilde\omega^2\,\Req,
\label{eq:rcirc-corot}
\end{equation}
puisque son moment cinétique spécifique vaut $\Omega_\oplus\Req^2$ à l'équateur. Même à la limite
de rotation, $\tilde\omega=1$, la matière injectée se circularise donc au voisinage de $\Req$, soit
$1{,}2$ à $1{,}3\,\Rt$, et non au-delà de $\aR$. L'impulsion de marée fournie par $B$ à $r=\Req$,
qui reste de l'ordre de quelques centièmes de $v_K$ d'après l'équation~\eqref{eq:impulsion-close},
ne comble pas l'écart.

On propose de ne pas lire ce résultat comme une perte. La matière injectée n'est ni éjectée ni
réaccrétée : elle rejoint le disque intérieur à Roche, c'est-à-dire la fraction
$\eta_{\mathrm{in}}$ de la partition~\eqref{eq:partition}, où elle est disponible pour le transport
de moment cinétique décrit à la section~\ref{sec:disque-accretion}. La chaîne complète devient
alors : $B$ prélève de la matière superficielle terrestre appauvrie en métal, cette matière se
circularise vers $1{,}3\,\Rt$, elle alimente le réservoir intérieur, et le transport visqueux ou
gravitationnel la porte ensuite au-delà de $\aR$ où elle s'accrète. La lignée matérielle est la
même d'un bout à l'autre de la chaîne, et les prédictions géochimiques du
chapitre~\ref{ch:traceurs} sont donc inchangées : c'est la même matière appauvrie qui arrive à la
Lune, quel que soit le chemin qu'elle emprunte pour franchir $\aR$.

La conséquence est méthodologique. Pour le canal d'injection, $\eta_L$ n'est pas la bonne mesure
d'efficacité, puisqu'il compte le placement \emph{direct} dans la fenêtre lunaire. L'efficacité
pertinente est celle de l'étalement, c'est-à-dire la fraction de $\eta_{\mathrm{in}}$ que le
transport porte au-delà de $\aR$ avant que le réservoir ne se dissipe. Le bilan
\eqref{eq:MLcalc-eff} doit donc être lu avec, pour le terme $M_{\oplus\to L}$, une efficacité
d'étalement et non une efficacité de tri.

\begin{controle}
Le balayage porte sur trois paramètres seulement, $\tilde\omega$, $q_B$ et $\mu_B$, les autres étant
fixés par la géométrie. La sensibilité attendue à $\tilde\omega$ est forte, et l'équation
\eqref{eq:rcirc-corot} indique que le résultat utile de ce calcul n'est pas la fraction atteignant
directement $\aR$, mais la masse totale injectée, sa distribution en $r_{\mathrm{circ}}$ et la
largeur de la fenêtre en rotation qui la permet, à mettre en regard du spin correspondant dans la
porte~\eqref{eq:DJEM}.
\end{controle}

\begin{falsif}
Si, sur toute la grille, la masse injectée reste négligeable devant $M_L$, ou si le temps
d'étalement nécessaire pour porter $\eta_{\mathrm{in}}$ au-delà de $\aR$ excède la durée de
persistance~\eqref{eq:tpers}, alors le canal d'injection ne peut pas compenser un réservoir léger,
et seule la branche massive du réservoir demeure ouverte, avec les contraintes de moment cinétique
associées.
\end{falsif}


\section{L'équation de la Lune}
\label{sec:equation-lune}

Les résultats des sections précédentes se rassemblent en une seule expression. La réduction repose
sur une simplification qui n'était pas visible au départ.

\paragraph{Les deux termes ont la même forme.} Une couche partant du rayon équatorial $\varpi$
emporte $j=\Omega_\oplus\varpi^2+\varpi\,\Delta v_B$. Le second terme, une fois développé avec la
forme close de l'impulsion, vaut
$\varpi\sqrt{G\Mt/\varpi}\;K\,(\varpi/q_B)^{3/2}$, où les puissances de $\varpi$ se recombinent
exactement en $\varpi^{2}$. Les deux contributions sont donc proportionnelles au même facteur, et
\begin{equation}
j(\varpi)=\sqrt{G\Mt\Rt}\;x^{2}\,\Lambda,
\qquad
\Lambda=\tilde\omega\left(\frac{\Rt}{\Req}\right)^{3/2}
+\frac{\sqrt2\,\mu_B}{\sqrt{1+\mu_B}}\left(\frac{\Rt}{q_B}\right)^{3/2},
\label{eq:Lambda}
\end{equation}
avec $x=\varpi/\Rt$. Le rayon de circularisation s'écrit alors, exactement,
\begin{equation}
\frac{r_{\mathrm{circ}}}{\Rt}=x^{4}\,\Lambda^{2}.
\label{eq:rcirc-close}
\end{equation}

\paragraph{Une seule équation.} La couche la plus externe, en $x=a\equiv\Req/\Rt$, est celle qui va
le plus loin. En posant $b\equiv q_B/\Rt$, la condition pour qu'elle franchisse la limite de Roche
se ramène à une quantité unique :
\begin{equation}
\boxed{\;
\mathcal{L}
=\tilde\omega\,\sqrt{a}
\;\pm\;\frac{\sqrt2\,\mu_B}{\sqrt{1+\mu_B}}\cdot\frac{a^{2}}{b^{3/2}}\;}
\label{eq:equation-lune}
\end{equation}
gouvernée par deux valeurs seuils :
\begin{equation}
\mathcal{L}\ \ge\ \sqrt{\aR/\Rt}=1{,}704
\quad\text{(existence)},
\qquad
\mathcal{L}\simeq1{,}80
\quad\text{(une masse lunaire)}.
\label{eq:seuils-lune}
\end{equation}

\paragraph{Le signe, et pourquoi il décide de tout.} L'impulsion de marée est tangentielle. Elle
s'ajoute au moment cinétique de rotation lorsque $B$ tourne dans le même sens que la proto-Terre,
et s'en retranche dans le cas contraire. Le signe $+$ correspond donc au passage prograde, le signe
$-$ au passage rétrograde, et l'écart entre les deux est considérable. Pour la configuration de
référence ci-dessous, le terme de rotation vaut $1{,}265$ et le terme de marée $0{,}535$, d'où
$\mathcal{L}=1{,}800$ en prograde et $\mathcal{L}=0{,}730$ en rétrograde. Le second est très
au-dessous du seuil d'existence : lors d'un passage rétrograde, \emph{aucune matière du canal
d'injection superficielle} ne franchit la limite de Roche. La fermeture est complète pour ce canal.
Elle ne s'étend pas au réservoir préexistant, qui relève d'un autre mécanisme et dont le calcul
pilote montre qu'une fraction résiduelle de $4\,\%$ atteint encore la fenêtre en rétrograde
(tableau~\ref{tab:pilote-etaL}).

\paragraph{L'origine physique de cet écart.} Elle tient à la cohérence temporelle. Pour la
configuration de référence, la vitesse angulaire de rotation de la proto-Terre vaut
$5{,}43\times10^{-4}$~rad\,s$^{-1}$ et celle de $B$ au périastre $4{,}83\times10^{-4}$, soit un
rapport de $1{,}12$. Les deux sont donc presque synchrones : un élément de la ceinture équatoriale
reste pratiquement sous $B$ pendant toute la rencontre, et les impulsions qu'il reçoit s'accumulent.
En rétrograde, les deux vitesses angulaires s'additionnent au lieu de se compenser, le déphasage
défile deux fois plus vite, et les impulsions alternées s'annulent presque. C'est le même
mécanisme de cohérence que révèle le rapport d'efficacité de $4{,}6$ entre les deux contrôles du
calcul pilote (tableau~\ref{tab:pilote-etaL}), mais on en tient ici le nombre.

\begin{ideephys}[Configuration de référence]
Les valeurs retenues dans la suite de cet ouvrage sont
$\tilde\omega=0{,}97$, $\Req=1{,}7\,\Rt$, $\mu_B=0{,}95$, $q_B=3{,}0\,\Rt$, passage prograde. Elles
donnent $\mathcal{L}=1{,}800$, soit exactement le seuil de masse lunaire, et
$M_{>\mathrm{R}}=1{,}05\,M_L$. Trois de leurs propriétés méritent d'être signalées. La dilatation
exigée du corps n'est que de $13\,\%$, ce qui est atteignable pour un manteau entièrement fondu,
là où un rayon de $2\,\Rt$ en aurait demandé $33\,\%$. Les marges géométriques, définies comme le rapport de $q_B$ au seuil concerné diminué de l'unité,
valent $12\,\%$ sur le contact et $9\,\%$ sur la dislocation de $B$, soit
$\chi_{\mathrm{tide}}\simeq1{,}09$. Et le rapport des
vitesses angulaires vaut $1{,}12$, c'est-à-dire que la configuration se place d'elle-même dans le
régime de cohérence maximale, sans que ce soit une condition imposée.
\end{ideephys}

Le premier seuil est exact et universel : c'est de la cinématique, et aucune combinaison de
paramètres ne le contourne. Le second dépend de la répartition de la masse dans les couches
externes, donc de la structure interne retenue.

\paragraph{Vérification.} Le long du contour produisant une masse lunaire, extrait numériquement
et reporté à la section précédente, $\mathcal{L}$ reste constante à $0{,}4\,\%$ :

\begin{center}
\begin{tabular}{lccccccc}
\toprule
$\mu_B$ & $0{,}30$ & $0{,}40$ & $0{,}50$ & $0{,}60$ & $0{,}70$ & $0{,}80$ & $0{,}85$ \\
$\Req/\Rt$ & $2{,}23$ & $2{,}09$ & $1{,}99$ & $1{,}89$ & $1{,}84$ & $1{,}79$ & $1{,}75$ \\
$\mathcal{L}$ & $1{,}805$ & $1{,}804$ & $1{,}808$ & $1{,}795$ & $1{,}811$ & $1{,}818$ & $1{,}804$ \\
\bottomrule
\end{tabular}
\end{center}

\begin{ideephys}[Ce que résume l'équation de la Lune]
Trois paramètres seulement entrent dans $\mathcal{L}$ : la rotation réduite de la proto-Terre, sa
dilatation, et le couple formé par la masse et la distance d'approche de $B$. Tout le reste de la
dynamique du passage se condense dans ces trois nombres. L'équation dit que ces trois quantités ne
comptent pas séparément mais par la valeur d'une somme unique, et que la Lune correspond à une
valeur particulière de cette somme. Le domaine viable n'est donc pas un volume dans l'espace des
paramètres : c'est une surface de niveau, et toute la question devient celle de la valeur du
niveau.
\end{ideephys}

\begin{controle}
Deux réserves d'emploi. L'équation~\eqref{eq:equation-lune} ne décrit que le canal d'injection ; si
le réservoir préexistant contribue par $\eta_L M_{D,0}$, le seuil de masse descend sous $1{,}81$.
Et la constante $1{,}81$ hérite de l'incertitude sur la structure interne, qui vaut un facteur
neuf sur la masse produite. La forme de la surface de niveau est donc mieux établie que sa
position, et déterminer cette constante est l'objectif prioritaire du programme numérique.
\end{controle}



\section{L'évacuation du moment cinétique par emport de matière}
\label{sec:evacuation}

Le passage laisse le système avec plus de moment cinétique qu'il n'en faut. Cette section propose
un mécanisme d'évacuation qui ne demande aucun ingrédient extérieur, car il est déjà contenu dans
le mécanisme d'alimentation.

\paragraph{De combien est l'excès.} Il faut d'abord l'évaluer correctement, et le point demande
un peu de soin. On écrit souvent le moment cinétique de rotation sous la forme
$S_\oplus=k\,\Mt\,\Omega_\oplus\Req^{2}$ avec $k\simeq0{,}25$, ce qui reviendrait à supposer que
toute la planète se dilate. Or le noyau métallique, qui porte environ un tiers de la masse, reste
compact au centre, où il ne contribue presque pas au moment d'inertie. En traitant séparément un
noyau de rayon inchangé et un manteau dilaté, le coefficient effectif tombe à
$k_{\mathrm{eff}}\simeq0{,}19$, et le spin correspondant passe de $2{,}9$ à $2{,}2\,L_{\mathrm{EM}}$
pour $\Req=2\,\Rt$ et $\tilde\omega=0{,}97$. L'excès à évacuer est donc un facteur deux, non un
facteur trois.

\paragraph{Ce que le survol ne fait pas.} Il serait tentant d'attribuer l'évacuation à
l'assistance gravitationnelle que subit $B$. Ce n'est pas le bon canal. Un survol est élastique
dans le référentiel de la proto-Terre : $B$ repart avec la même vitesse relative, seule sa
direction change. Ce qui s'échange est du moment cinétique \emph{héliocentrique}, dont le réservoir
vaut quelque $7{,}6\times10^{5}\,L_{\mathrm{EM}}$ et qui constitue un budget distinct de celui de
la rotation propre. Quant au couple de marée exercé sur la proto-Terre pendant la rencontre, il
existe mais reste faible : pour $k_2/Q$ compris entre $0{,}02$ et $0{,}10$, un passage unique
n'emporte que $0{,}02$ à $0{,}12\,L_{\mathrm{EM}}$. Ces deux voies sont donc à écarter, et il est
utile de le dire explicitement.

\paragraph{Ce que l'emport de matière fait.} La matière qui quitte définitivement le système
emporte son moment cinétique avec elle. Et comme elle part de la ceinture équatoriale, son moment
cinétique spécifique vaut $\Omega_\oplus\Req^{2}$, soit $1/k_{\mathrm{eff}}\simeq5{,}4$ fois la
valeur moyenne du corps. Chaque unité de masse perdue emporte donc cinq fois plus de rotation
qu'une unité de masse prise au hasard dans la planète. Le levier est considérable, et il conduit à
des masses modestes :

\begin{center}
\begin{tabular}{cc}
\toprule
\textbf{Spin résiduel visé} & \textbf{Masse à emporter} \\
\midrule
$1{,}5\,L_{\mathrm{EM}}$ & $4{,}8\,M_L$ \\
$1{,}2\,L_{\mathrm{EM}}$ & $6{,}9\,M_L$ \\
$1{,}0\,L_{\mathrm{EM}}$ & $8{,}3\,M_L$ \\
\bottomrule
\end{tabular}
\end{center}

\begin{ideephys}[Le mécanisme d'alimentation est aussi le mécanisme d'évacuation]
La masse mobilisée par le débordement atteint une trentaine de masses lunaires. Il suffit donc
qu'environ un cinquième de cette matière quitte définitivement le système pour que le bilan de
moment cinétique se ferme. Or c'est précisément ce que fait un débordement de lobe de Roche dans
une binaire : la matière franchit le point de Lagrange intérieur et s'écoule vers le compagnon. $B$
se trouve à quelques rayons terrestres et repart sur une trajectoire hyperbolique ; tout ce qu'il
capture part avec lui sans retour. Le même passage fournit ainsi la masse de la Lune, sa signature
chimique, et l'évacuation du moment cinétique excédentaire, sans qu'aucun mécanisme extérieur soit
invoqué.
\end{ideephys}


\begin{figure}[htbp]
\centering
\begin{tikzpicture}[scale=1.0,>=Stealth]
% limite de Roche
\draw[dashed,gray!70] (0,0) circle (3.0);
\node[gray!70,font=\scriptsize,fill=white,inner sep=1pt] at (-2.35,2.35) {limite de Roche};
% proto-Terre aplatie
\draw[fill=orange!25,draw=orange!70!black,thick] (0,0) ellipse (1.30 and 0.95);
% noyau
\draw[fill=gray!55,draw=gray!70!black] (0,0) circle (0.40);
% bande equatoriale
\draw[line width=2.6pt,red!65!black,opacity=0.75] (-1.30,0.06) -- (1.30,0.06);
\draw[line width=2.6pt,red!65!black,opacity=0.75] (-1.30,-0.06) -- (1.30,-0.06);
\node[red!60!black,font=\scriptsize,anchor=north,fill=white,inner sep=1pt] at (0,-1.12) {bande intertropicale};
% planete B et sa trajectoire
\draw[thick,blue!60!black,->] (4.2,-3.1) .. controls (2.5,-1.0) and (2.5,1.0) .. (4.2,3.1);
\draw[fill=blue!35,draw=blue!60!black] (2.15,0) circle (0.34);
\node[blue!50!black,font=\scriptsize,anchor=west] at (2.55,0.0) {$B$};
% flux vers B
\draw[->,line width=1.5pt,red!70!black] (1.35,0.10) -- (1.78,0.10);
\node[red!65!black,font=\scriptsize,anchor=south,align=center] at (1.6,0.30) {vers $B$,\\ sans retour};
% matiere qui retombe dans le disque interieur
\draw[->,line width=1.1pt,teal!65!black] (-1.35,0.15) arc (170:250:1.9);
\node[teal!55!black,font=\scriptsize,anchor=east,align=right] at (-1.6,0.55) {vers le disque\\ intérieur};
% matiere qui atteint la fenetre lunaire
\draw[->,line width=1.1pt,violet!70!black] (-1.25,0.55) .. controls (-2.4,1.9) .. (-0.4,3.1);
\node[violet!65!black,font=\scriptsize,anchor=south,align=center] at (-0.4,3.2) {vers la fenêtre lunaire};
\node[font=\scriptsize,anchor=north,align=center] at (0,-1.72) {noyau métallique compact,\\ manteau silicaté dilaté};
\end{tikzpicture}
\caption{Les trois destinations de la matière soulevée. La bande équatoriale, en rouge, est la seule
région d'où la matière peut partir. Une petite fraction atteint la fenêtre lunaire au-delà de la
limite de Roche et formera la Lune. La plus grande part retombe dans le disque intérieur, où le
transport visqueux la reprendra. Et une fraction quitte le système avec $B$, emportant avec elle le
moment cinétique excédentaire. Le même passage remplit ainsi les trois fonctions.}
\label{fig:evacuation}
\end{figure}

\paragraph{Une porte supplémentaire.} Ce canal se formule comme une condition sur la fraction
emportée :
\begin{equation}
\mathcal{D}_{\mathrm{evac}}=\left\{\Theta_I:\ \phi_{\oplus\to B}(\Theta_I)\ \ge\ \phi_{\oplus\to B}^{\min}\right\},
\qquad \phi_{\oplus\to B}^{\min}\simeq0{,}20 ,
\label{eq:Devac}
\end{equation}
où $\phi_{\oplus\to B}$ désigne la part de la masse \emph{arrachée à la proto-Terre} qui quitte la sphère de
Hill terrestre liée à $B$. Cette quantité ne doit pas être confondue avec $\eta_{D\to B}$ du
chapitre~\ref{ch:fenetre}, qui rapporte au réservoir préexistant $M_{D,0}$ la fraction de
\emph{disque} capturée par $B$. Les deux populations et les deux dénominateurs diffèrent, et le
calcul pilote donne d'ailleurs $\eta_{D\to B}\simeq0$ en prograde sans que cela contraigne
$\phi_{\oplus\to B}$, qui porte sur une tout autre matière. Cette définition restrictive est
délibérée. Il serait tentant d'y ajouter $\eta_{\mathrm{esc}}$, la matière simplement non liée à la
proto-Terre, mais celle-ci ne disparaît pas : elle reste sur des orbites héliocentriques voisines
de $1$~UA et une part reviendra s'accréter en quelques millions d'années. Seule la matière partant
avec $B$ est perdue sans retour, puisqu'elle franchit la sphère de Hill en sa compagnie, et le sort
héliocentrique ultérieur de $B$ n'y change rien. $\phi_{\oplus\to B}$ est ainsi une sortie du calcul de
Phase~B, qui suit déjà chaque élément et sait lequel finit lié à $B$.

\paragraph{Conséquence sur les constantes du problème.} Une proto-Terre plus massive déplace
plusieurs quantités de référence, et il vaut mieux les chiffrer que les laisser en suspens. Pour un
surcroît de masse de $10\,\%$, la limite de Roche et le rayon de Hill, tous deux en
$\Mt^{1/3}$, augmentent de $3{,}3\,\%$. Le seuil d'existence de l'équation de la Lune passe de
$1{,}704$ à $1{,}73$, et la constante de masse lunaire de $1{,}80$ à $1{,}83$, soit un déplacement
de $1{,}6\,\%$. C'est quatre fois moins que l'incertitude que laisse déjà le choix du profil
interne. L'effet est donc absorbé, mais il doit être porté au bilan et non négligé par principe.

\begin{controle}
Ce canal impose une contrainte supplémentaire, et c'est une bonne chose : il devient une prédiction.
Si $8\,M_L$ quittent le système, la proto-Terre était plus massive que la Terre actuelle d'environ
$10\,\%$. Cette valeur est compatible avec un scénario d'accrétion tardive, mais elle doit être
portée au bilan de masse et vérifiée pour sa cohérence chronologique. Elle est par ailleurs sans
incidence sur les contraintes isotopiques, puisque la matière emportée par $B$ ne participe pas à
la Lune.
\end{controle}


\section{Le bilan de moment cinétique est vectoriel}
\label{sec:vectoriel}

Le bilan a été conduit jusqu'ici sur des grandeurs scalaires, ce qui suffit tant que le passage se
fait dans le plan équatorial de la proto-Terre. Dès que $i_B\ne0$, cette simplification cesse d'être
légitime, et il vaut la peine de dire pourquoi, car c'est peut-être là que se trouve le résultat le
plus intéressant encore à obtenir.

\paragraph{Ce qu'un traitement scalaire laisse de côté.} Le couple exercé par $B$ sur une
proto-Terre aplatie n'est pas aligné avec l'axe de rotation lorsque le passage est incliné. Il
possède une composante perpendiculaire, qui fait basculer cet axe. Un bilan scalaire enregistre la
variation du module et ignore ce basculement. Or le basculement est précisément ce qui produit une
obliquité, et l'obliquité est une observable.

\paragraph{Deux cibles.} Le système Terre--Lune présente deux orientations qu'aucun bilan scalaire
ne peut engendrer. L'obliquité terrestre vaut $23{,}4^\circ$. Et l'orbite lunaire est inclinée de
$5{,}1^\circ$ sur l'écliptique, alors qu'une lune formée dans un disque équatorial devrait naître
dans le plan de l'équateur. Cette seconde valeur est un problème ouvert et reconnu
\cite{PahlevanMorbidelli2015}.

\begin{ideephys}[Une prédiction encore inexploitée]
Le chemin proposé dispose d'un paramètre libre dont il ne tire pour l'instant aucun parti :
l'inclinaison $i_B$ du passage. Un traitement vectoriel du couple exercé par $B$ relierait cette
inclinaison à l'orientation finale du système, et donc à deux quantités mesurées plutôt qu'à une
seule. Si une même valeur de $i_B$ rendait compte à la fois de l'obliquité terrestre et de
l'inclinaison de l'orbite lunaire, la théorie gagnerait une contrainte qu'aucune de ses versions
scalaires ne peut fournir. C'est, à notre sens, le développement le plus prometteur à entreprendre
après la fermeture du bilan de masse.
\end{ideephys}

\paragraph{Ce que cela demande.} Le formalisme de la Phase~B suit déjà chaque élément en trois
dimensions ; il suffit d'accumuler le moment cinétique sous forme vectorielle plutôt que scalaire,
et de projeter l'état final dans le repère barycentrique. La prédiction P7, qui porte sur
$i_{D,1}$, devient alors le premier terme d'un couple $(i_{D,1},\varepsilon_\oplus)$ à comparer aux
deux valeurs observées. Aucune physique nouvelle n'est requise, seulement un changement de nature
des quantités suivies.


\paragraph{Durée effective du passage.}\label{par:duree} Une distinction s'impose ici, car elle a été source de
confusion. Le temps caractéristique $\sqrt{q_B^{3}/G(\Mt+M_B)}$ vaut environ une heure pour la
configuration de référence, mais ce n'est pas la durée de la rencontre : c'est seulement l'échelle
de temps au périastre. La durée pendant laquelle le champ différentiel de $B$ reste comparable à sa
valeur maximale est bien plus longue.

\begin{center}
\begin{tabular}{lc}
\toprule
\textbf{Grandeur mesurée} & \textbf{Durée} \\
\midrule
temps caractéristique $\sqrt{q_B^{3}/G M_{\mathrm{tot}}}$ & $0{,}95$~h \\
temps passé en deçà de $2\,q_B$ & $3{,}3$~h \\
temps passé en deçà de $3\,q_B$ & $5{,}8$~h \\
temps passé en deçà de $5\,q_B$ & $11{,}2$~h \\
\bottomrule
\end{tabular}
\end{center}

La rencontre effective dure donc entre cinq et douze heures, soit cinq à dix fois le temps
caractéristique. Ce point est important pour le bilan d'injection : la période de rotation de la
proto-Terre étant de $3{,}2$~h, une rencontre d'une heure ne présenterait qu'un quart de la
ceinture équatoriale au champ de $B$, tandis qu'une rencontre de six heures lui en présente près de
deux tours complets. Toute la bande est balayée, et la zone la plus favorable l'est plusieurs fois.
Le paramètre $\eta_\oplus=\Omega_\oplus t_{p}$ vaut alors environ $1{,}5$ et non $1$ ;
on reste dans le régime où les deux échelles se confondent, ce qui justifie la porte
$\mathcal{D}_{\mathrm{imp}}$, mais la valeur doit être corrigée. On note enfin que la vitesse à
l'infini n'a presque aucune influence sur ces durées, la vitesse au périastre étant dominée par la
chute gravitationnelle ; le seul levier réel est $q_B$, en $q_B^{3/2}$.



\section{Un corps froid dans un milieu torride}
\label{sec:eau}

La section précédente a établi que le débordement fonctionne dans les deux sens. Il reste à
examiner ce que $B$ cède, et la réponse dépend entièrement de ce qu'il porte à sa surface.

\paragraph{Le contraste thermique.} $B$ arrive des régions extérieures du système, où les corps
rocheux se couvrent de glace. Il entre dans un environnement dont la pièce centrale est un océan
magmatique. Le flux reçu à $3\,\Rt$ d'une surface portée à $3700$~K vaut
$3{,}4\times10^{6}$~W\,m$^{-2}$, soit environ deux mille cinq cents fois la constante solaire, et
porte la face exposée de $B$ à l'équilibre radiatif vers $2800$~K. La transition est brutale : en
quelques minutes, une surface glacée passe à une température qui dépasse la fusion du silicate.

\paragraph{Le chauffage ne détache rien.} Il faut se garder d'une conclusion trop rapide. La
vitesse de libération de $B$ vaut environ $11$~km\,s$^{-1}$, tandis que la vitesse thermique de la
vapeur d'eau, même à $2800$~K, atteint seulement $2$~km\,s$^{-1}$, soit moins d'un cinquième. La
sublimation à elle seule ne livre d'ailleurs presque rien : en six heures, l'énergie reçue permet
de sublimer $6{,}5\times10^{18}$~kg, soit cinq millièmes d'océan terrestre. La chaleur latente de
sublimation est trop grande et la rencontre trop brève.

\begin{ideephys}[Le chauffage rend la matière livrable, la marée la livre]
Le rôle du flash thermique est ailleurs, et il est décisif. Une coquille de glace solide ne
s'écoule pas ; une atmosphère de vapeur, si. Or la hauteur d'échelle d'une enveloppe de vapeur
d'eau passe de $14$~km à $300$~K à $133$~km à $2800$~K : l'enveloppe est dilatée d'un facteur dix.
Le sommet de l'enveloppe se rapproche alors du lobe de Roche de $B$, et ce qui le franchit s'écoule.
C'est le champ de marée qui expulse, le chauffage se contente de rendre la matière expulsable.
\end{ideephys}

\paragraph{Les quantités en jeu.} Pour $\mu_B=0{,}95$ et $q_B=3{,}0\,\Rt$, le lobe de Roche de $B$
se situe à $1{,}12\,\Rt$ tandis que son corps rocheux mesure $0{,}99\,\Rt$. Une enveloppe de
volatils de $874$~km suffit donc à l'atteindre ; elle représenterait sept pour cent de la masse de
$B$, proportion modeste pour un corps formé au-delà de la ligne des glaces, et l'équivalent de
plusieurs centaines d'océans terrestres.

Du côté terrestre, la demande est dérisoire : un océan correspond à $0{,}025\,\%$ de la masse de
$B$, dix océans à $0{,}25\,\%$. Le mécanisme se situe donc très loin de tout réglage fin, et livre
largement au-delà du nécessaire bien avant de livrer trop peu. C'est un régime confortable, et il
faut le signaler comme tel.

\paragraph{L'eau lunaire appartient à la même lignée.} Il serait tentant de conclure que la Lune
naît sèche, le disque à la limite de Roche étant porté vers $2000$~K, bien au-dessus de la
condensation de l'eau. Les mesures interdisent cette conclusion. La teneur en eau pré-éruptive des
verres volcaniques lunaires est estimée à plusieurs centaines de parties par million
\cite{Saal2008}, valeur comparable à celle du manteau supérieur terrestre, et la composition
isotopique de l'hydrogène lunaire s'est révélée indiscernable de celle de l'eau terrestre
\cite{Saal2013}.

Cette observation renforce la théorie plutôt qu'elle ne la gêne, et pour la raison qui vaut déjà
pour l'oxygène, le titane, le chrome et le tungstène : \emph{c'est la même lignée de matière}. La
Lune contient une eau isotopiquement terrestre parce qu'elle est faite de manteau terrestre. Le
déficit relatif s'explique par la volatilité dans un disque chaud, la parenté isotopique par la
provenance. L'hydrogène rejoint ainsi la liste des traceurs qui pointent vers une origine
terrestre commune.

\begin{controle}
Deux points restent à établir, et ils sont contraignants. D'abord la chronologie : une eau arrivée
pendant les quelques heures de la rencontre aurait peine à se mélanger au manteau avant le départ
de la matière lunaire ; la livraison étalée décrite à la section~\ref{sec:trainee} lève cette
difficulté, mais l'ordre des événements doit être explicité. Ensuite la composition : si $B$
apporte une part de l'eau terrestre, son rapport deutérium sur hydrogène devient une contrainte
forte, puisque l'eau terrestre correspond aux chondrites carbonées et non aux comètes. La région de
formation de $B$ s'en trouverait contrainte, ce qui constitue un test du présent scénario.
\end{controle}

\section{La traînée laissée par \texorpdfstring{$B$}{B}}
\label{sec:trainee}

Une conséquence du mécanisme mérite d'être signalée, non parce qu'elle explique quelque chose, mais
parce qu'elle se produit nécessairement et qu'elle laisse quelque chose derrière elle.

\paragraph{Le débordement fonctionne dans les deux sens.} La section~\ref{sec:evacuation} a
considéré le débordement de la proto-Terre. Mais à l'approche minimale, les deux corps ont des lobes
de Roche presque identiques, et rien n'interdit à $B$ de déborder à son tour. Un corps venu des
régions froides du système, s'il porte une enveloppe de volatils de quelques centaines de
kilomètres, la voit dilatée par le rayonnement de l'océan magmatique : le flux reçu à $3\,\Rt$
d'une surface à $3700$~K vaut environ deux mille cinq cents fois la constante solaire, et porte la
face exposée de $B$ au voisinage de $2800$~K. La hauteur d'échelle de l'enveloppe est alors
multipliée par dix, et son sommet atteint le lobe.

Il faut souligner que le chauffage seul ne détache rien. La vitesse de libération de $B$ vaut
environ $11$~km\,s$^{-1}$, tandis que la vitesse thermique de la vapeur d'eau, même à $2800$~K,
n'atteint que $2$~km\,s$^{-1}$. Le chauffage ne libère pas la matière, il la rend libérable : c'est
le champ de marée qui l'expulse, et c'est pourquoi elle sort par les points de Lagrange plutôt que
dans toutes les directions.

\paragraph{Deux portes de sortie, deux destins.} Une enveloppe qui déborde ne s'échappe pas par un
seul endroit. Elle sort par les deux points de Lagrange de la configuration. Par le point intérieur,
la matière s'écoule vers la proto-Terre et son disque. Par le point extérieur, elle part vers
l'arrière et quitte le système pour de bon, laissant une traînée le long de la trajectoire de $B$.

Les deux flux sont concomitants : ils ont la même cause et l'un ne va pas sans l'autre. La traînée
n'est donc pas une hypothèse ajoutée au mécanisme, c'est une conséquence obligatoire de son
fonctionnement.

\begin{ideephys}[Une livraison étalée plutôt qu'instantanée]
La traînée est laissée au voisinage de l'orbite terrestre, sur des trajectoires héliocentriques
proches de celle de la Terre. Elle continue donc de croiser la Terre longtemps après le passage, et
se réaccrète progressivement sur des durées de l'ordre de dix à cent millions d'années. C'est une
propriété utile, car elle dissocie deux échelles de temps que le passage confondait : la rencontre
dure quelques heures, mais l'apport de matière qu'elle initie s'étale sur des dizaines de millions
d'années. Une objection chronologique naturelle, celle du temps nécessaire pour incorporer au
manteau une matière arrivée en six heures, se trouve levée : la plus grande part n'arrive
précisément pas en six heures.
\end{ideephys}

\paragraph{Un ordre de grandeur, et rien de plus.} Il existe une contrainte indépendante sur
l'apport de matière postérieur à la ségrégation du noyau terrestre, établie à partir de l'abondance
des éléments fortement sidérophiles dans le manteau. Elle correspond à quelques dixièmes de pour
cent de la masse terrestre, soit de l'ordre de $3\times10^{22}$~kg. Une enveloppe de volatils de
$874$~km sur $B$ en représenterait $4\times10^{23}$~kg. La fraction de la traînée qu'il faudrait
voir retomber sur la Terre pour atteindre cet ordre de grandeur est donc de quelques pour cent.

On se borne ici à ce constat de disponibilité. Il est présenté comme un apport du mécanisme, non
comme une explication de la contrainte.

\begin{controle}
Une réserve doit accompagner ce constat. Le rapport des apports tardifs entre la Terre et la Lune
est observé de l'ordre de $10^{3}$, alors que le rapport des sections efficaces de collision entre
les deux corps n'est que de quelques dizaines. Une traînée réaccrétée ne rend donc pas compte à
elle seule de cet écart, qui constitue un problème ouvert indépendant du présent travail. La
question reste ouverte, et il serait imprudent de la présenter autrement.
\end{controle}

\section{La constante de l'équation de la Lune}
\label{sec:constante-lune}

L'équation de la section précédente comporte deux nombres. Le premier, $1{,}704$, est exact et ne
dépend de rien : il exprime seulement la distance à franchir pour atteindre la limite de Roche. Le
second, celui qui correspond à une masse lunaire, dépend de la façon dont la matière est répartie
dans les couches externes de la proto-Terre. C'est le seul endroit de toute la théorie où la
structure interne du corps entre en jeu, et il vaut la peine de s'y arrêter.

\paragraph{Pourquoi la répartition compte.} La matière qui franchit la limite de Roche est celle
qui part du plus loin de l'axe de rotation, c'est-à-dire la plus superficielle. Sa quantité dépend
donc de l'épaisseur de la couche la plus externe, et cette épaisseur dépend à son tour de la
compressibilité du matériau. Un corps peu compressé garde beaucoup de matière près de sa surface ;
un corps fortement condensé concentre sa masse au centre et n'en laisse que peu en surface. À
passage identique, le premier fournit donc davantage que le second, et il faut un passage plus
énergique pour compenser.

On décrit conventionnellement ce degré de condensation par un indice : plus il est élevé, plus la
masse est concentrée au centre. La constante de l'équation de la Lune en dépend de la manière
suivante.

\begin{center}
\begin{tabular}{ccl}
\toprule
\textbf{Indice} & \textbf{Constante} & \textbf{État du corps correspondant} \\
\midrule
$0{,}5$ & $1{,}77$ & corps peu condensé, proche de l'incompressible \\
$1{,}0$ & $1{,}80$ & silicate comprimé, cas de référence \\
$1{,}5$ & $1{,}88$ & manteau silicaté chaud et convectif \\
$2{,}0$ & $2{,}07$ & manteau très chaud, océan magmatique étendu \\
\bottomrule
\end{tabular}
\end{center}

\begin{ideephys}[Une incertitude qui s'effondre]
Le tableau ci-dessus est, à notre sens, la meilleure justification du détour par l'équation de la
Lune. Sur la masse produite, passer d'un bout à l'autre de cet intervalle d'indices change le
résultat d'un facteur neuf : c'est l'incertitude qui pesait le plus lourd dans tout le programme
numérique. Sur la constante, le même changement ne déplace le résultat que de six pour cent autour
de $1{,}9$. La reformulation n'a pas supprimé l'incertitude, elle l'a confinée dans une quantité où
elle devient négligeable. On ne cherche donc plus une constante inconnue : on la connaît à quelques
pour cent près, et ce qui reste à faire est de choisir l'indice, non de recommencer le calcul.
\end{ideephys}

\paragraph{La forme du domaine ne bouge pas.} Un second résultat mérite d'être signalé, car il
n'allait pas de soi. Quel que soit l'indice retenu, la constante reste constante le long du
contour, à deux ou trois dixièmes de pour cent près, y compris lorsque l'on fait varier la distance
d'approche et pas seulement la masse du perturbateur. Autrement dit, changer de structure interne
déplace le niveau de la surface, mais ne la déforme pas. La propriété la plus utile du résultat,
celle qui ramène trois paramètres à un seul nombre, survit donc intacte au changement de profil.

\paragraph{Une simplification qui allège la démonstration.} Le calcul de la
section~\ref{sec:bilan-injection} passe par deux quantités intermédiaires : la masse mobilisable et
la fraction de celle-ci qui franchit la limite de Roche. Or la seconde est un rapport dont le
dénominateur est précisément la première. Lorsqu'on multiplie l'une par l'autre pour obtenir la
masse utile, la masse mobilisable se simplifie, et il ne reste que la masse située au-delà du rayon
seuil. Le lobe de Roche ne sert donc plus qu'à vérifier qu'il descend bien plus bas que ce rayon,
ce qui est le cas dans tout le domaine étudié, mais il n'intervient plus dans le résultat. Cette
remarque raccourcit la chaîne de raisonnement sans en changer les conclusions.

\begin{controle}
Le point contraignant est à la borne haute du tableau. Pour une proto-Terre dont les couches
externes seraient très chaudes et très dilatées, donc pour un indice voisin de $2$, le domaine se
ferme à la distance d'approche utilisée jusqu'ici : atteindre une masse lunaire exigerait une
dilatation telle que la proto-Terre et $B$ entreraient en contact, et la porte $\mathcal{D}_{\mathrm{nc}}$
se referme avant que la masse soit atteinte. Le domaine n'est pas perdu pour autant, mais il se
déplace : il rouvre pour un perturbateur de masse presque terrestre passant un peu plus loin, avec
une constante de $2{,}07$ et la même stabilité qu'ailleurs. La borne supérieure de l'intervalle
d'indices plausible est donc aussi une frontière du domaine viable, et déterminer de quel côté se
situe la proto-Terre réelle est la question la plus contraignante que pose ce chapitre.
\end{controle}

\section{Premier calcul pilote de tri orbital}
\label{sec:pilote-etaL}

Un calcul de Phase~B a été réalisé pour vérifier que la fenêtre lunaire \eqref{eq:fenetre} peut
être atteinte dans un domaine non ponctuel de paramètres. Il constitue un diagnostic de tendance
plutôt qu'une preuve hydrodynamique : le réservoir est représenté par $504$ particules test de même
masse, initialement circulaires et coplanaires, distribuées entre $\Req$ et $\aR$ avec
$\Sigma\propto r^{-1}$. Les paramètres sont ceux de la configuration de référence
(section~\ref{sec:equation-lune}) : $\Req=1{,}70\,\Rt$, $M_B=0{,}95\,\Mt$, $q_B=3{,}0\,\Rt$ et
$v_\infty=2$~km\,s$^{-1}$, avec $\omega_B=90^\circ$. L'accélération employée est le champ
différentiel exact de l'équation~\eqref{eq:adiff}, intégré sur la trajectoire hyperbolique complète.

\begin{table}[htbp]
\centering
\small
\caption{Partition du réservoir après le passage, en fonction de l'inclinaison, pour la
configuration de référence. Calcul par particules test, sans autogravité du disque, sans collisions
mutuelles, sans pression et sans hydrodynamique. La dernière ligne donne le contrôle rétrograde.}
\label{tab:pilote-etaL}
\begin{tabular}{l|ccccc}
\toprule
& $\eta_L$ & $\eta_{\mathrm{in}}$ & $\eta_{\mathrm{reaccr}}$ & $\eta_{\mathrm{esc}}$ &
$\eta_{D\to B}$ \\
\midrule
$i_B=20^\circ$ & $0{,}151$ & $0{,}675$ & $0{,}071$ & $0{,}103$ & $0{,}000$ \\
$i_B=25^\circ$ & $0{,}185$ & $0{,}655$ & $0{,}067$ & $0{,}093$ & $0{,}000$ \\
$i_B=30^\circ$ & $0{,}187$ & $0{,}653$ & $0{,}079$ & $0{,}081$ & $0{,}000$ \\
\midrule
rétrograde, $25^\circ$ & $0{,}040$ & $0{,}036$ & $0{,}246$ & $0{,}425$ & $0{,}254$ \\
\bottomrule
\end{tabular}
\end{table}

Trois enseignements se dégagent de ce tableau.

D'abord, le réservoir survit au passage. Malgré une masse du perturbateur voisine de celle de la
proto-Terre, la fraction échappée demeure inférieure à $11\,\%$ et les deux tiers du réservoir
restent liés dans le domaine intérieur, disponibles pour le transport décrit à la
section~\ref{sec:disque-accretion}. Une estimation fondée sur la seule amplitude de l'impulsion
\eqref{eq:impulsion-close} aurait conduit à prédire une dispersion complète ; l'intégration exacte
montre qu'il n'en est rien, l'impulsion n'étant que partiellement tangentielle et son signe
alternant avec l'azimut.

Ensuite, l'efficacité croît de $20$ à $25^\circ$ puis plafonne, ce qui recoupe indépendamment la
fenêtre d'inclinaison imposée par le survol du disque.

Enfin, le contrôle rétrograde est sans appel. Le rapport des efficacités atteint $4{,}6$, et
surtout la structure de la partition change de nature : en rétrograde, $42{,}5\,\%$ du réservoir
s'échappe et $25{,}4\,\%$ part avec $B$, contre $9\,\%$ et une fraction négligeable en prograde. La
dispersion du réservoir que l'on pouvait craindre se produit bien, mais uniquement dans le sens
rétrograde. La cohérence temporelle établie à la section~\ref{sec:equation-lune} protège donc aussi
le réservoir, et pas seulement le canal d'injection.

Pour le cas retenu, $q_B=3\,\Rt$ et $i_B=25^\circ$, on a donc
\begin{equation}
\eta_L=0{,}185,\qquad \eta_{\mathrm{in}}=0{,}655,\qquad
\eta_{\mathrm{reaccr}}=0{,}067,\qquad \eta_{\mathrm{esc}}=0{,}093 .
\label{eq:pilote-cas}
\end{equation}
À masse initiale fixée, la masse directement placée dans la fenêtre est
\begin{equation}
M_{L,\mathrm{fenetre}}\simeq0{,}185\,M_{D,0},
\label{eq:MLfenetre}
\end{equation}
ce qui signifie qu'une masse lunaire immédiatement placée dans cette fenêtre demanderait
$M_{D,<R}\simeq5{,}4\,M_L$ dans ce calcul test. La notation $M_{D,<R}$ rappelle que les $504$
particules sont toutes initialement intérieures à la limite de Roche : la masse inférée concerne
cette population testée, et non l'ensemble du réservoir, qui peut s'étendre au-delà. Aucune
particule ne satisfaisant le critère de fenêtre avant le passage, on a par construction
\begin{equation}
\eta_L^{(0)}=0,\qquad \Delta\eta_L=\eta_L,
\label{eq:deltaetaL}
\end{equation}
ce qui raccorde directement le tableau au critère de succès~\eqref{eq:fenetre-succes}.

Ces nombres sont des diagnostics de tri, non un bilan de formation lunaire. La quantité $\eta_L$
mesure une efficacité de tri orbital ; il faut encore que la matière triée survive à la relaxation
puis s'accrète. Le bilan de masse pertinent s'écrit donc
\begin{equation}
\boxed{\;
M_{L,\mathrm{calc}}=
\epsilon_{\mathrm{acc}}\,\epsilon_{\mathrm{surv}}\,\eta_L\,M_{D,0}
+M_{\oplus\to L}+M_{B\to L}
\;\pm\;\delta M_{L,\mathrm{imp}},
\;}
\label{eq:MLcalc-eff}
\end{equation}
où $\delta M_{L,\mathrm{imp}}$ est l'incertitude systématique introduite par l'approximation
rectiligne, estimée par la comparaison de la forme close~\eqref{eq:impulsion-close} avec
l'intégration directe sur la conique hyperbolique (porte~\eqref{eq:Dimp}).

Deux enseignements en découlent. D'abord, comme
$\epsilon_{\mathrm{acc}}\epsilon_{\mathrm{surv}}\le1$, tout raffinement de ces efficacités augmente
la masse de réservoir requise à Lune fixée, et donc aussi $L_{D,0}$, qui lui est proportionnel : le
gain se paie sur la porte $\mathcal{D}_J^{\mathrm{EM}}$~\eqref{eq:DJEM}. Ensuite, et c'est le point
important, le terme $M_{\oplus\to L}$ ne dépend pas du tout de $M_{D,0}$. C'est par lui que la
théorie peut satisfaire $\mathcal{D}_{\mathrm{accr}}$ avec un réservoir modeste, et c'est donc lui
qu'il faut chiffrer en priorité (section~\ref{sec:bilan-injection}).

\paragraph{Moment cinétique associé.}
Un réservoir de $5{,}4\,M_L$ ainsi structuré porte, en rotation képlérienne,
\begin{equation}
L_{D,0}=2\pi\!\int_{R_{\mathrm{in}}}^{R_{\mathrm{out}}}\!\Sigma(r)\sqrt{G\Mt r}\;r\,\dd r
\simeq3{,}0\times10^{34}\ \text{kg\,m}^2\,\text{s}^{-1}\simeq0{,}86\,L_{\mathrm{EM}},
\label{eq:LD-reservoir}
\end{equation}
auquel s'ajoute le spin de la proto-Terre, de l'ordre de $1{,}3\,L_{\mathrm{EM}}$ pour une période
de $3$~h. Deux conséquences utiles. D'une part, $L_{D,0}$ étant proportionnel à $M_{D,0}$, la
branche légère du réservoir allège d'autant le bilan : c'est une raison supplémentaire de la
préférer. D'autre part, le passage n'a pas à fournir le moment cinétique du système final, qui est
déjà disponible ; sa fonction est bien de le redistribuer. Il reste que l'état laissé par le
passage devra rejoindre $L_{\mathrm{EM}}$ au terme de l'évolution maréale, ce qui suppose un
mécanisme d'évacuation. Trois voies sont envisageables : une capture en résonance d'évection
solaire ou une transition de plan de Laplace \cite{MurrayDermott1999,TremaineTouma2009}, l'une et
l'autre relevant de la mécanique céleste classique du système Terre--satellite et intervenant en
Phase~C ; ou bien l'emport direct de matière par $B$ au moment du passage, qui relève de la Phase~B
et fait l'objet de la section~\ref{sec:evacuation}. C'est
l'objet de la porte $\mathcal{D}_J^{\mathrm{EM}}$~\eqref{eq:DJEM} et de la prédiction P8.

Le contrôle rétrograde du tableau~\ref{tab:pilote-etaL} isole le rôle du sens de passage. Pour $q_B=3\,\Rt$, à latitude de
périastre identique et avec les autres paramètres inchangés, on obtient
\begin{equation}
\eta_{L,\mathrm{retro}}\simeq0{,}040,
\label{eq:retro-pilote}
\end{equation}
soit une efficacité environ $4{,}6$ fois plus faible que dans le cas prograde. La partition change
en outre de nature, comme le montre le tableau~\ref{tab:pilote-etaL} : la fraction échappée passe
de $9{,}3$ à $42{,}5\,\%$ et la fraction emportée par $B$ de zéro à $25{,}4\,\%$. Ce résultat est cohérent
avec l'argument de durée de cohérence de l'équation~\eqref{eq:prograde}.

L'origine radiale de la matière promue confirme que le tri agit principalement sur la partie externe
du réservoir, voisine de Roche : pour le cas $q_B=3\,\Rt$, $i_B=30^\circ$, les contributions à la
cohorte $\eta_L$ proviennent à $35{,}0\,\%$ de l'anneau $2{,}53$--$2{,}71\,\Rt$ et à
$53{,}1\,\%$ de l'anneau $2{,}71$--$2{,}90\,\Rt$.

\begin{controle}
Ce calcul pilote ne ferme pas encore la théorie : il néglige la masse propre du réservoir, les chocs,
la viscosité, la pression, l'autogravité, les collisions mutuelles et l'aplatissement complet de la
proto-Terre. Deux précisions aident à le situer. D'abord son domaine de validité en masse : pour
$M_{D,0}=5{,}4\,M_L$, soit $0{,}066\,\Mt$, le réservoir reste sous un dixième de $M_B$ et son
autogravité demeure petite devant le forçage mesuré ; la branche légère, $M_{D,0}\simeq2$--$3\,M_L$,
vers laquelle pointent le contrôle~\eqref{eq:controle-Ida} et le
critère~\eqref{eq:sigmar-reservoir}, serait plus confortable encore. Le
tableau~\ref{tab:pilote-etaL} se lit donc comme une carte de tendance en $(q_B,i_B)$, à masse
fixée. Ensuite son domaine de validité en mécanisme : les $504$ particules représentent le seul
réservoir préexistant, sans matière injectée depuis la surface terrestre ; la carte mesure
$\eta_L$, c'est-à-dire le premier terme de l'équation~\eqref{eq:MLcalc-eff}, et laisse
$M_{\oplus\to D}$ entièrement à évaluer. Il établit toutefois un résultat concret de Niveau~I : un passage prograde de $B$ peut
placer une fraction mesurable du réservoir préexistant dans la fenêtre lunaire sans transformer
l'essentiel du disque en matière échappée. La prochaine étape est de remplacer cette carte de
particules test par une carte massive et dissipative.
\end{controle}

\begin{falsif}
Si, sur l'ensemble de la grille, $\Delta\eta_L$ reste négligeable, ou n'augmente qu'au prix d'une
fraction échappée incompatible avec la masse du réservoir, le mécanisme de tri orbital est
insuffisant et la branche correspondante est exclue.
\end{falsif}

\begin{rappel}
Le domaine de viabilité fixe ce que le calcul doit démontrer. La dernière partie énonce ce que la
nature devrait montrer si la théorie est correcte, et comment le vérifier.
\end{rappel}

% =====================================================================
\part{Prédictions}
\label{part:predictions}
% =====================================================================

\noindent
Chaque prédiction est énoncée sous la même forme : ce que l'on doit observer \emph{si la théorie
est correcte}, puis \emph{comment le vérifier}, par la mesure d'échantillons lunaires ou par le
calcul. Les prédictions portent sur l'architecture, la composition, l'eau, la chronologie et la
dynamique.

% ---------------------------------------------------------------------
\chapter{Prédictions et vérifications}
\label{ch:predictions}
\label{sec:signatures}
% ---------------------------------------------------------------------

\begin{prediction}{P1. Un satellite dominant}
\sila Le réservoir trié et relaxé produit une masse satellitaire concentrée presque entièrement
dans un corps unique :
\begin{equation}
\mathcal{D}_{\mathrm{arch}}=\left\{\Theta_I:\ M_L\simeq M_{L,\mathrm{obs}},\
N_{\mathrm{sat,grand}}\simeq1,\ f_{\mathrm{dom}}\to1\right\}.
\label{eq:Darch}
\end{equation}
\verif Simulations à $N$ corps de la Phase~C, avec collisions et autogravité, à partir des
réservoirs réorganisés de la Phase~B ; suivi de $N_{\mathrm{sat}}$ et de $f_{\mathrm{dom}}$.
\end{prediction}

\begin{prediction}{P2. Une part négligeable de $B$}
\sila La fraction de matière lunaire issue de $B$,
\begin{equation}
f_B=\frac{M_{B\to L}}{M_L},
\label{eq:fB}
\end{equation}
est faible, et les compositions isotopiques de l'oxygène, du titane et du chrome de la Lune sont
indiscernables de celles de la Terre silicatée, dans les limites~\eqref{eq:bornefB}.
\verif Mesures isotopiques de haute précision sur des échantillons lunaires variés ; calcul de
$M_{B\to D}$ en Phase~B pour les passages avec $q_B$ proche de $d_{R,B}$.
\end{prediction}

\begin{prediction}{P3. Un tungstène terrestre}
\sila Après correction de l'accrétion tardive, $\varepsilon^{182}$W lunaire est identique à celui du
manteau terrestre, et le rapport Hf/W de la matière lunaire ne produit pas d'excès de $^{182}$W
incompatible avec l'instant $t_B$.
\verif Mesures de $\varepsilon^{182}$W et de Hf/W sur des roches lunaires à faible exposition
cosmique ; comparaison avec la chronologie de la Phase~A.
\end{prediction}

\begin{prediction}{P4. L'héritage de la ségrégation du noyau terrestre}
\sila Les appauvrissements lunaires en V, Cr et Mn reproduisent ceux du manteau terrestre, et le
rapport W/U lunaire est proche du rapport terrestre.
\verif Estimations de la composition du manteau lunaire à partir des basaltes et des roches
primitives ; expériences de partage métal--silicate à haute pression.
\end{prediction}

\begin{prediction}{P5. Une corrélation dynamique--chimique}
\sila Une même cohorte de matière porte simultanément une histoire orbitale et une composition :
\begin{equation}
(r_0,\lambda_\Omega,z)\longleftrightarrow(\varepsilon_{\mathrm{orb}},j_z,q_p,\mathbf{X}),
\label{eq:correlation}
\end{equation}
et le rapport W/V de la matière injectée~\eqref{eq:WV} reflète la profondeur et la latitude
arrachées.
\verif Suivi des étiquettes $(\lambda_\Omega,z)$ en Phase~B ; comparaison du W/V calculé avec les
rapports mesurés, normalisés à des éléments de même comportement.
\end{prediction}

\begin{prediction}{P6. Un réservoir de composition terrestre}
\sila Le réservoir circumterrestre préexistant est de nature silicatée terrestre, alimenté
principalement par la perte de masse équatoriale et les éjectas du manteau, et l'excès de FeO
lunaire est produit par les transformations internes au réservoir.
\verif Bilans de Phase~A sur les canaux~\eqref{eq:canaux} ; calcul de
$\boldsymbol{\Pi}_D$ et comparaison avec les teneurs lunaires en FeO et en sidérophiles.
\end{prediction}

\begin{prediction}{P7. Une inclinaison initiale non nulle}
\sila Un passage légèrement incliné donne au réservoir réorganisé un moment cinétique non
colinéaire avec le spin terrestre :
\begin{equation}
i_{D,1}=\arccos\!\left(\frac{\mathbf{L}_{D,1}\cdot\mathbf{S}_{\oplus,1}}
{|\mathbf{L}_{D,1}|\,|\mathbf{S}_{\oplus,1}|}\right).
\label{eq:inclinaison}
\end{equation}
Cette grandeur constitue une condition initiale de l'évolution maréale du système Terre--Lune,
dont l'inclinaison orbitale lunaire est l'une des énigmes classiques \cite{PahlevanMorbidelli2015}.
\verif Calcul de $i_{D,1}$ en Phase~B, puis intégration de l'évolution maréale jusqu'à l'orbite
actuelle.
\end{prediction}

\begin{prediction}{P8. Un moment cinétique compatible après évacuation}
\sila Le moment cinétique du système Terre--réservoir laissé par le passage, $J_{\oplus D,1}$,
dépasse $L_{\mathrm{EM}}$ d'autant plus que le réservoir est massif et la proto-Terre rapide. La
théorie prédit donc qu'une évacuation post-accrétionnelle est nécessaire, et que le couple
$(J_{\oplus D,1},i_{D,1})$ produit par le passage appartient au domaine de capture d'un mécanisme
connu.
\verif Fermeture du bilan~\eqref{eq:fermeture} en Phase~B, puis intégration maréale de
$J_{\oplus D,1}$ jusqu'à $L_{\mathrm{EM}}$, $J_{\mathrm{loss}}$ étant une sortie du calcul
par résonance d'évection ou transition de plan de Laplace
\cite{MurrayDermott1999,TremaineTouma2009}. Si aucun couple $(J_{\oplus D,1},i_{D,1})$ du domaine
viable n'est capturable, la porte~\eqref{eq:DJEM} se ferme sur la branche correspondante.
\end{prediction}

\begin{prediction}{P9. Un tri orbital efficace dans la bande dynamique de basses latitudes}
\sila Il existe une région continue de la bande dynamique de basses latitudes prograde où $\Delta\eta_L$ est
significatif et $\eta_{\mathrm{esc}}$ modéré, conformément à~\eqref{eq:fenetre-succes}, tout en
franchissant les portes $\mathcal{D}_J^{\mathrm{EM}}$, $\mathcal{D}_{\mathrm{grav}}$,
$\mathcal{D}_{\mathrm{disque},0}$, $\mathcal{D}_{\mathrm{accr}}$,
$\mathcal{D}_{\mathrm{tide},B}$ et $\mathcal{D}_{\mathrm{imp}}$.
\verif Cartes de viabilité de la Phase~B (chapitre~\ref{ch:programme}), croisées avec le bilan
d'injection de la section~\ref{sec:bilan-injection} et le bilan de moment cinétique.
\end{prediction}

\begin{prediction}{P10. Un destin héliocentrique éliminé pour \texorpdfstring{$B$}{B}}
\sila Le corps $B$ appartient à la population des corps planétaires primitifs ensuite éliminés :
\begin{equation}
B\in\mathcal{P}_{\mathrm{elim}},\qquad t_B<t_{\mathrm{ej}},
\label{eq:predB}
\end{equation}
et l'orbite héliocentrique post-passage $\mathcal{H}_{B,1}$ est compatible avec une diffusion
ultérieure vers une collision, une accrétion par une planète ou une éjection du Système solaire.
\verif Intégrations $N$-corps héliocentriques démarrant à $\mathcal{H}_{B,1}$, avec scénarios
d'instabilité précoce et tardive, et vérification que les architectures finales restent compatibles
avec le Système solaire actuel.
\end{prediction}

\begin{prediction}{P11. Une eau lunaire isotopiquement terrestre}
\sila La Lune contient de l'eau en quantité comparable à celle du manteau supérieur terrestre
(quelques centaines de ppm), et son rapport D/H est indiscernable de celui de l'eau terrestre.
L'eau provient de la lignée matérielle terrestre, complétée par un apport modéré de $B$ lors du
passage.
\verif Mesures de la teneur en eau pré-éruptive des verres volcaniques et des inclusions fondues
lunaires \cite{Saal2008} ; mesures du rapport D/H lunaire et comparaison avec les chondrites
carbonées \cite{Saal2013} ; calcul du bilan d'eau en Phase~B et en Phase~C, avec suivi des
étiquettes isotopiques de l'hydrogène.
\end{prediction}

\chapter*{Conclusion générale}
\addcontentsline{toc}{part}{Conclusion générale}
\markboth{Conclusion générale}{Conclusion générale}
% =====================================================================

La question posée au début de cet ouvrage était celle d'un chemin : comment le système
Terre--Lune observé aujourd'hui a-t-il pu émerger du système proto-terrestre par une évolution
physiquement cohérente ?

La théorie propose une réponse précise. Pendant sa croissance, la proto-Terre est accompagnée
d'un réservoir secondaire circumterrestre, analogue par son principe aux systèmes d'accrétion des
planètes géantes et au disque de PDS~70c, mais dominé par les solides et alimenté principalement
par sa propre couche silicatée. Ce réservoir contient déjà de la masse et du moment cinétique
lorsque le corps planétaire massif $B$ passe.

$B$ n'apporte pas l'essentiel de la matière lunaire. C'est un agent
gravitationnel transitoire de redistribution des flux de masse et de moment cinétique entre les
réservoirs. Les passages rasants sont au moins aussi fréquents que les collisions
(équation~\eqref{eq:frequence}) ; la trajectoire de $B$ porte un moment cinétique de l'ordre de celui
du système Terre--Lune (équation~\eqref{eq:LBrel}) ; dans la bande dynamique de basses latitudes d'une proto-Terre
oblate en rotation rapide, son passage prograde injecte une matière superficielle déjà appauvrie
en métal. Son résultat central est un tri orbital : une cohorte liée est placée, par le couple de $B$,
dans la fenêtre lunaire au-delà de la limite de Roche (équations~\eqref{eq:fenetre}
et~\eqref{eq:etaL}), tandis que d'autres fractions réaccrètent ou s'échappent.

Le calcul de Phase~B conduit avec la configuration de référence donne à ce tri une première
valeur. Pour $i_B=25^\circ$, $18{,}5\,\%$ du réservoir satisfait directement le critère de fenêtre
lunaire et $65{,}5\,\%$ demeure lié dans le disque intérieur, soit $84\,\%$ de masse conservée dans
le domaine circumterrestre, contre $9{,}3\,\%$ échappés. Le réservoir survit donc au passage, bien
que la masse du perturbateur soit voisine de celle de la proto-Terre. Le contrôle rétrograde, conduit
à inclinaison identique, sépare nettement les deux régimes : la fraction de fenêtre tombe à
$4{,}0\,\%$, la fraction échappée monte à $42{,}5\,\%$ et un quart du réservoir part avec $B$. La
cohérence temporelle du passage prograde protège ainsi le réservoir autant qu'elle alimente
l'injection.

Le réservoir trié se relaxe et accrète la Lune, en concentrant la masse satellitaire dans un corps
dominant. La géochimie appartient à la même histoire : le passage non collisionnel favorise une
faible contribution de $B$, ce qui fait de la similitude isotopique une contrainte attendue du
domaine viable ; le vanadium et le tungstène, lus conjointement, tracent les flux de matière
silicatée terrestre et la profondeur de la couche prélevée.

Le même passage fournit un mécanisme physique d'apport d'eau. Venu des régions froides, $B$ entre dans un environnement
dont la pièce centrale est un océan magmatique, y reçoit un flux de l'ordre de deux mille cinq cents
constantes solaires, et voit son enveloppe de volatils dilatée au point de déborder son propre lobe
de Roche. Le chauffage ne détache rien par lui-même, la vitesse de libération étant cinq fois
supérieure à la vitesse thermique de la vapeur ; il rend la matière expulsable et c'est le champ de
marée qui l'expulse. Deux flux en résultent, l'un vers le système terrestre, l'autre vers une traînée
dont la réaccrétion s'étale sur des dizaines de millions d'années. L'eau lunaire, dont la signature
isotopique est indiscernable de l'eau terrestre, appartient quant à elle à la lignée terrestre, au
même titre que l'oxygène, le titane, le chrome et le tungstène.

La logique du modèle se referme avec le destin de $B$. L'agent qui réorganise le réservoir
circumterrestre n'a pas à demeurer dans le Système solaire final : il peut appartenir à la population
des corps planétaires primitifs ensuite éliminés par l'instabilité orbitale globale, pourvu que son
passage vérifie $t_B<t_{\mathrm{ej}}$. Cette hypothèse relie la formation Terre--Lune à la migration
et à l'instabilité du Système solaire primitif, sans imposer que $B$ soit identifié à une planète
manquante unique.

La question scientifique finale s'écrit
\begin{equation}
\boxed{\;\mathcal{S}_0=\{\Theta_I\}\longrightarrow\mathcal{D}_0
\xrightarrow{\;\text{tri orbital par }B\;}\mathcal{D}_1\longrightarrow
\left(M_L,\,L_L,\,i_{L,0},\,\mathbf{X}_L\right)\in\mathcal{D}_{\mathrm{viable}}
\longrightarrow\mathcal{S}_f.\;}
\label{eq:question}
\end{equation}
La théorie ne prétend pas démontrer une histoire unique. Elle définit un chemin physique
quantitatif et falsifiable, et elle indique, pour chacune de ses étapes, comment le confirmer, le
restreindre ou l'exclure. La loi de transport par passage rasant est une \emph{forme analytique
sous hypothèses}, dont la validation sur la conique hyperbolique est une sortie de la Phase~B et
non un paramètre ajusté. C'est cette validation qui transforme le seuil statique en modèle fermé.

% =====================================================================
\appendix
% =====================================================================

\chapter{Démonstrations}
\label{app:demo}

\section{Rayon de circularisation de la matière capturée}
\label{app:rc}

Une matière capturée à travers une zone d'accrétion de rayon $R_{\mathrm{acc}}$ dans un
écoulement cisaillé de fréquence $\Omega$ porte un moment cinétique spécifique d'ordre
$j\simeq\Omega R_{\mathrm{acc}}^2/4$ \cite{QuillenTrilling1998}. Elle se circularise au rayon où son
moment cinétique égale le moment képlérien, $j^2=G\Mt r_c$ :
\begin{equation}
r_c=\frac{\Omega^2R_{\mathrm{acc}}^4}{16\,G\Mt}.
\end{equation}
Pour $R_{\mathrm{acc}}=R_H$, la définition du rayon de Hill, $R_H^3=G\Mt/(3\Omega^2)$, donne
$\Omega^2=G\Mt/(3R_H^3)$ et $r_c=R_H/48$. Pour une accrétion limitée à la sphère de Bondi,
$r_c$ est multiplié par $(R_{\mathrm{B}}/R_H)^4$, ce qui conduit à l'équation~\eqref{eq:rc}.

\section{Fréquence des passages rasants}

Pour une trajectoire hyperbolique de paramètre d'impact $b$, la conservation du moment
cinétique et de l'énergie donne $b\,v_\infty=q\,v_p$ et $v_p^2=v_\infty^2+2G\Mtot/q$, d'où
\begin{equation}
b^2=q^2\left(1+\frac{2G\Mtot}{q\,v_\infty^2}\right),\qquad \sigma(q)=\pi b^2,
\end{equation}
qui est l'équation~\eqref{eq:sigma}. Avec $\Theta=2G\Mtot/(R_cv_\infty^2)$,
$\sigma(R_c)=\pi R_c^2(1+\Theta)$ et $\sigma(2R_c)=\pi R_c^2(4+2\Theta)$ ; le rapport
$[\sigma(2R_c)-\sigma(R_c)]/\sigma(R_c)=(3+\Theta)/(1+\Theta)$ décroît de 3 à 1 lorsque $\Theta$ croît
de 0 à l'infini.

\section{Nœuds de la trajectoire et non-traversée}

L'équation de la conique s'écrit $r=p/(1+e\cos\nu)$ avec $p=q(1+e)$. La trajectoire traverse le
plan de référence lorsque l'argument de latitude $\omega+\nu$ vaut $0$ ou $\pi$, soit $\nu=-\omega$ ou
$\nu=\pi-\omega$ ; les rayons correspondants sont $p/(1+e\cos\omega)$ et $p/(1-e\cos\omega)$, dont le
plus petit vaut $p/(1+e|\cos\omega|)$. C'est l'équation~\eqref{eq:nontraversee}.

\section{Rayon de circularisation et périgée}
\label{app:rcirc}

Pour une orbite liée, $j^2=G\Mt\,a(1-e^2)$ et $q_p=a(1-e)$, donc
$r_{\mathrm{circ}}=j^2/(G\Mt)=a(1-e)(1+e)=q_p(1+e)$. Comme $0\le e<1$,
$q_p\le r_{\mathrm{circ}}<2q_p$, ce qui établit l'équation~\eqref{eq:rcirc-qp}.

\section{Frontières du tri orbital}
\label{app:tri}

Un élément sur orbite circulaire de rayon $r_0$ a la vitesse $v_K=\sqrt{G\Mt/r_0}$. Après une
impulsion tangentielle, $v=v_K(1+\delta)$ et $j=r_0v_K(1+\delta)$, d'où
$r_{\mathrm{circ}}=r_0(1+\delta)^2$. Son énergie spécifique vaut
\begin{equation}
\varepsilon=\frac{G\Mt}{r_0}\left[\frac{(1+\delta)^2}{2}-1\right],
\end{equation}
qui s'annule pour $(1+\delta)^2=2$ : $\delta_{\mathrm{esc}}=\sqrt2-1$. Pour $\delta<0$, $r_0$ devient
l'apogée ; le demi-grand axe vaut $a=r_0/[2-(1+\delta)^2]$ et le périgée
$q_p=2a-r_0=r_0(1+\delta)^2/[2-(1+\delta)^2]$. La condition $q_p=\Req$ donne
$(1+\delta)^2=2\Req/(r_0+\Req)$, soit $\delta_{\mathrm{re}}$. La condition $r_{\mathrm{circ}}=\aR$
donne $\delta_L=\sqrt{\aR/r_0}-1$. On obtient l'équation~\eqref{eq:frontieres}.

\section{Impulsion de marée : forme analytique sous hypothèses}
\label{app:impulsion}

\paragraph{Statut de ce résultat.}
Cette section établit une \emph{forme analytique de référence} pour l'impulsion de marée reçue par
un élément du réservoir pendant le passage de $B$. Elle n'est pas exacte au sens strict : elle
repose sur trois hypothèses explicites, énoncées ci-dessous. Son domaine de validité est précisé,
et la Phase~B du programme numérique est chargée de la valider par intégration directe sur la vraie
conique hyperbolique.

\paragraph{Hypothèses.}
\begin{enumerate}
\item \textbf{Trajectoire rectiligne.} La trajectoire de $B$ est assimilée à une droite de vitesse
$v_p$ et de paramètre d'impact $q_B$ pendant toute la durée du passage.
\item \textbf{Forçage impulsif.} La durée de la rencontre est courte devant la période orbitale
locale : $\eta(r)=\Omega_K(r)\,t_{p}\ll1$.
\item \textbf{Particule fixe.} L'élément de matière ne se déplace pas pendant le passage ; son
mouvement orbital propre est négligé devant l'impulsion reçue.
\end{enumerate}

\paragraph{Conventions.}
Repère géocentrique cartésien $(O;x,y,z)$ : l'axe $x$ est porté par la vitesse de $B$ au périastre,
l'axe $y$ est dirigé du centre de la proto-Terre vers $B$ au périastre, l'axe $z$ complète le trièdre
direct. La trajectoire nominale s'écrit $\mathbf{r}_B(t)=(v_p t,\,q_B,\,0)$.

\paragraph{Dérivation.}
Le champ différentiel exact est donné par l'équation~\eqref{eq:adiff}. En intégrant séparément le
terme direct et le terme indirect sur la trajectoire rectiligne, on obtient la forme close
suivante, valable pour tout $r<q_B$ :
\begin{equation}
\boxed{\;
\Delta\mathbf{v}(\mathbf{r})
=
-\frac{GM_B}{v_p\,q_B}
\left(
0,\;
\frac{2(\eta-1)}{(\eta-1)^2+\zeta^2}+2,\;
\frac{2\zeta}{(\eta-1)^2+\zeta^2}
\right)
\;}
\label{eq:impulsion-close}
\end{equation}
avec les coordonnées réduites $\eta=y/q_B$ et $\zeta=z/q_B$.

L'approximation au premier ordre en $r/q_B$ s'en déduit :
\begin{equation}
\Delta\mathbf{v}^{\text{approx}}(\mathbf{r})
=
-\frac{2GM_B}{v_p\,q_B^2}\,(0,\;y,\;-z)
+
O\!\left(\frac{r^2}{q_B^3}\right).
\label{eq:impulsion-premier-ordre}
\end{equation}
C'est cette expression qui figure historiquement dans la littérature sous le nom d'approximation
impulsionnelle.

\paragraph{Comparaison et domaine de validité.}
Le tableau~\ref{tab:impulsion-validite} compare les deux formes pour une particule située sur l'axe
du paramètre d'impact ($\zeta=0$) puis sur l'axe perpendiculaire ($\eta=0$). L'écart relatif est
calculé par rapport à la forme close~\eqref{eq:impulsion-close}.

\begin{table}[htbp]
\centering
\small
\caption{Écart entre la forme close~\eqref{eq:impulsion-close} et l'approximation au premier
ordre~\eqref{eq:impulsion-premier-ordre}.}
\label{tab:impulsion-validite}
\begin{tabular}{ccc}
\toprule
$r/q_B$ & Écart sur l'axe du paramètre d'impact & Écart sur l'axe perpendiculaire \\
\midrule
$0{,}1$ & $+11\,\%$ & $-1\,\%$ \\
$0{,}3$ & $+43\,\%$ & $-8\,\%$ \\
$0{,}5$ & $+100\,\%$ & $-20\,\%$ \\
$0{,}7$ & $+233\,\%$ & $-33\,\%$ \\
$0{,}9$ & $+900\,\%$ & $-45\,\%$ \\
\bottomrule
\end{tabular}
\end{table}

L'approximation au premier ordre perd toute précision dès $r/q_B\gtrsim0{,}3$. Pour le cas de
référence du manuscrit ($q_B=3\,\Rt$, réservoir jusqu'à $\sim2{,}9\,\Rt$), le rapport $r/q_B$ varie
de $0{,}53$ à $0{,}97$ : la forme close~\eqref{eq:impulsion-close} doit donc être utilisée comme
référence, et l'approximation au premier ordre réservée aux contrôles d'ordre de grandeur.

\paragraph{Borne d'ordre de grandeur.}
Dans la limite $r\ll q_B$, la forme close redonne l'estimation utilisée dans le corps du texte :
\begin{equation}
\frac{\Delta v}{v_K}\sim\sqrt2\,\frac{\mu_B}{\sqrt{1+\mu_B}}
\left(\frac{r}{q_B}\right)^{3/2}.
\label{eq:impulsion-borne}
\end{equation}
Cette borne est utile pour fixer les ordres de grandeur ; elle ne remplace pas la
forme~\eqref{eq:impulsion-close}.

\paragraph{Limite résiduelle.}
La forme close~\eqref{eq:impulsion-close} suppose encore la trajectoire rectiligne. Pour le cas de
référence ($e_B\simeq1{,}1$), la courbure n'est pas négligeable et une intégration directe sur la
vraie conique hyperbolique est nécessaire. Cette intégration est réalisée en Phase~B, et l'écart
avec la forme close est une \emph{sortie} du calcul, quantifiée en fonction de $e_B$ et de
$r/q_B$. C'est cette validation croisée qui autorise à présenter la loi de transport comme
\emph{analytique sous hypothèses}, et non comme exacte au sens strict.

\section{Énergie de circularisation}

Une parcelle d'énergie $\varepsilon_i=-G\Mt/(2a)$ se circularise à moment cinétique conservé au
rayon $a_{\mathrm{circ}}=a(1-e^2)$, d'énergie $\varepsilon_f=-G\Mt/(2a_{\mathrm{circ}})$. L'énergie
libérée vaut
$\varepsilon_i-\varepsilon_f=\frac{G\Mt}{2a}\left[\frac{1}{1-e^2}-1\right]=\frac{G\Mt}{2a}\frac{e^2}{1-e^2}$,
ce qui est l'équation~\eqref{eq:circ}.

\section{Seuil d'injection}

Au point sous $B$, l'accélération de marée radiale vaut, lorsque $\lambda_B=0$,
\begin{equation}
GM_B\left[(q_B-\Req)^{-2}-q_B^{-2}\right]\simeq2GM_B\Req/q_B^3.
\end{equation}
À basse latitude, sa projection sur la direction radiale est multipliée par
$F(\lambda_B)=(3\cos^2\lambda_B-1)/2$. L'égaler à la gravité effective
$(G\Mt/\Req^2)(1-\tilde\omega^2)$ donne
\begin{equation}
1-\tilde\omega^2\lesssim2F(\lambda_B)\,\frac{M_B}{\Mt}\left(\frac{\Req}{q_B}\right)^3,
\end{equation}
qui est l'équation~\eqref{eq:seuil}.

\section{Borne sur la part chondritique du réservoir}
\label{app:borneW}

Avec $f_B=0$ et $\varepsilon^{182}\mathrm{W}_\oplus=0$, l'équation~\eqref{eq:melange} s'écrit
$\varepsilon^{182}\mathrm{W}_L=f_0C_{\mathrm{W},0}\varepsilon^{182}\mathrm{W}_0/
[f_0C_{\mathrm{W},0}+(1-f_0)C_{\mathrm{W},\oplus}]$. Pour $f_0C_{\mathrm{W},0}\ll C_{\mathrm{W},\oplus}$,
le dénominateur se réduit à $C_{\mathrm{W},\oplus}$, et la condition
$|\varepsilon^{182}\mathrm{W}_L|<\delta\varepsilon$ donne l'équation~\eqref{eq:bornef0}.

% ---------------------------------------------------------------------
\chapter{Notations}
\label{app:notations}
% ---------------------------------------------------------------------

\small
\begin{longtable}{>{\raggedright\arraybackslash}p{3.2cm}>{\raggedright\arraybackslash}p{11.6cm}}
\toprule
\textbf{Symbole} & \textbf{Signification} \\
\midrule
\endhead
$\mathcal{S}_0$, $\mathcal{S}_f$ & état initial du système proto-terrestre, état observé \\
$\RT$, $\RD$, $\RB$, $\Rinf$ & réservoirs : proto-Terre, réservoir circumterrestre, corps $B$, milieu extérieur \\
$\Phi_{i\to j}$ & flux de masse du réservoir $i$ vers le réservoir $j$ \\
$\boldsymbol{\mathcal{T}}_{j\to i}$ & couple exercé par le réservoir $j$ sur le réservoir $i$ \\
$\mathbf{X}_i$, $\mathbf{X}_{i\to j}$ & composition d'un réservoir, d'un flux \\
$\Mt$, $\Req$, $\Omega_\oplus$, $S_\oplus$ & masse, rayon équatorial, vitesse angulaire et spin de la proto-Terre \\
$\tilde\omega$ & rotation réduite $\Omega_\oplus/\sqrt{G\Mt/\Req^3}$ \\
$J_2$ & coefficient d'aplatissement gravitationnel \\
$M_D$, $L_D$, $\Sigma(r)$ & masse, moment cinétique et densité surfacique du réservoir \\
$\mathcal{D}_0$, $\mathcal{D}_1$ & état du réservoir avant et après la rencontre \\
$\mathcal{G}(t)$ & population d'agrégats et de lunules \\
$R_{\mathrm{in}}$, $R_{\mathrm{out}}$ & rayons interne et externe du réservoir \\
$R_H$, $R_{\mathrm{B}}$ & rayons de Hill et de Bondi \\
$\aR$, $a_{\mathrm{sync}}$ & limite de Roche, orbite synchrone \\
$M_B$, $\mu_B$, $R_B$, $\rho_B$ & masse, masse relative, rayon et densité de $B$ \\
$q_B$, $v_\infty$, $v_p$, $e_B$ & périastre, vitesse asymptotique, vitesse au périastre, excentricité de la trajectoire de $B$ \\
$i_B$, $\omega_B$ & inclinaison sur l'équateur rotationnel, argument du périastre \\
$\lambda_B$ & latitude du point sous $B$ au périastre \\
$\lambda_c$ & demi-largeur de la bande dynamique de basses latitudes \\
$t_{p}$, $\eta(r)$ & durée de la rencontre, paramètre de régime $\Omega_Kt_{p}$ \\
$d_{R,B}$ & distance de Roche de $B$ autour de la proto-Terre \\
$L_{B,\mathrm{rel}}$, $\chi_J$ & moment orbital relatif de $B$, fraction échangée \\
$\varepsilon$, $j$, $j_z$, $e$, $q_p$ & énergie, moment cinétique, projection, excentricité et périgée d'un élément \\
$r_{\mathrm{circ}}$ & rayon de circularisation $j^2/(G\Mt)$ \\
$\eta_L$, $\Delta\eta_L$ & fraction dans la fenêtre lunaire, contribution propre de $B$ \\
$\eta_{\mathrm{in}}$, $\eta_{\mathrm{reaccr}}$, $\eta_{\mathrm{esc}}$, $\eta_{D\to B}$ & fractions : disque intérieur, réaccrétion, échappement, capture par $B$ \\
$\delta$ & impulsion tangentielle relative $\Delta v_t/v_K$ \\
$f_0$, $f_\oplus$, $f_B$ & fractions de provenance de la matière lunaire \\
$f_{\mathrm{dom}}$, $N_{\mathrm{sat}}$ & facteur de dominance, nombre de satellites \\
$\lambda_\Omega$, $z$ & latitude dynamique et profondeur d'origine de la matière injectée \\
$\mathcal{P}_{\mathrm{elim}}$ & population des corps planétaires primitifs ensuite éliminés par collision, accrétion, diffusion ou éjection \\
$t_{\mathrm{ej}}$ & instant d'éjection ou d'élimination dynamique de $B$ \\
$\mathcal{H}_{B,1}$ & orbite héliocentrique de $B$ après son passage près de la proto-Terre \\
$\Theta_I$ & vecteur des paramètres du Niveau I \\
$\mathcal{D}_{\mathrm{viable}}$ & domaine de viabilité \\
$\mathcal{D}_{B,\mathrm{hel}}$ & porte héliocentrique vérifiant le destin éliminé de $B$ \\
$\mathcal{D}_{\mathrm{imp}}$ & porte de validation de l'impulsion sur la conique hyperbolique \\
$\mathcal{D}_{\mathrm{eau}}$ & porte de cohérence de l'eau lunaire et du rapport D/H \\
\bottomrule
\end{longtable}
\normalsize

% ---------------------------------------------------------------------
\chapter{Constantes et valeurs numériques}
\label{app:constantes}
% ---------------------------------------------------------------------

\begin{table}[htbp]
\centering
\small
\caption{Constantes et valeurs de référence utilisées dans l'ouvrage.}
\label{tab:constantes}
\begin{tabular}{lll}
\toprule
\textbf{Grandeur} & \textbf{Symbole} & \textbf{Valeur} \\
\midrule
Paramètre gravitationnel terrestre & $G\Mt$ & $3{,}986\times10^{14}$ m$^3$\,s$^{-2}$ \\
Masse de la Terre & $\Mt$ & $5{,}972\times10^{24}$ kg \\
Rayon terrestre moyen & $\Rt$ & $6{,}371\times10^{6}$ m \\
Masse de la Lune & $M_L$ & $7{,}342\times10^{22}$ kg \\
Masse du Soleil & $M_\odot$ & $1{,}989\times10^{30}$ kg \\
Unité astronomique & $a_\oplus$ & $1{,}496\times10^{11}$ m \\
Densité lunaire moyenne & $\rho_s$ & $3{,}34$ g\,cm$^{-3}$ \\
Densité terrestre moyenne & $\bar\rho_\oplus$ & $5{,}51$ g\,cm$^{-3}$ \\
Moment cinétique Terre--Lune & $L_{\mathrm{EM}}$ & $\approx3{,}5\times10^{34}$ kg\,m$^2$\,s$^{-1}$ \\
Limite de Roche ($\rho_s=3{,}34$) & $\aR$ & $\approx2{,}9\,\Rt$ \\
Demi-vie du $^{182}$Hf & $t_{1/2}$ & $8{,}9$ Ma \\
\bottomrule
\end{tabular}
\end{table}

\clearpage
\phantomsection
\addcontentsline{toc}{part}{Bibliographie}
\markboth{Bibliographie}{Bibliographie}
\begin{thebibliography}{99}
\small

\bibitem{MurrayDermott1999} C.~D. Murray et S.~F. Dermott, \emph{Solar System Dynamics},
Cambridge University Press, 1999.

\bibitem{TremaineTouma2009} S. Tremaine, J. Touma et F. Namouni, « Satellite dynamics on the
Laplace surface », \emph{The Astronomical Journal}, 137 (2009), 3706--3717.
doi:10.1088/0004-6256/137/3/3706.

\bibitem{Saal2008} A.~E. Saal, E.~H. Hauri, M. Lo Cascio, J.~A. Van Orman, M.~C. Rutherford et
R.~F. Cooper, « Volatile content of lunar volcanic glasses and the presence of water in the Moon's
interior », \emph{Nature}, 454 (2008), 192--195. doi:10.1038/nature07047.

\bibitem{Saal2013} A.~E. Saal, E.~H. Hauri, J.~A. Van Orman et M.~J. Rutherford, « Hydrogen
isotopes in lunar volcanic glasses and melt inclusions reveal a carbonaceous chondrite heritage »,
\emph{Science}, 340 (2013), 1317--1320. doi:10.1126/science.1235142.

\bibitem{BalbusHawley1991} S.~A. Balbus et J.~F. Hawley, « A powerful local shear instability in
weakly magnetized disks. I. Linear analysis », \emph{The Astrophysical Journal}, 376 (1991),
214--222. doi:10.1086/170270.

\bibitem{ShakuraSunyaev1973} N.~I. Shakura et R.~A. Sunyaev, « Black holes in binary systems.
Observational appearance », \emph{Astronomy and Astrophysics}, 24 (1973), 337--355.

\bibitem{FrankKingRaine2002} J. Frank, A. King et D. Raine, \emph{Accretion Power in Astrophysics},
3\ieme{} édition, Cambridge University Press, 2002.

\bibitem{Pringle1981} J.~E. Pringle, « Accretion discs in astrophysics », \emph{Annual Review of
Astronomy and Astrophysics}, 19 (1981), 137--160. doi:10.1146/annurev.aa.19.090181.001033.

\bibitem{JohansenLacerda2010} A. Johansen et P. Lacerda, « Prograde rotation of protoplanets by
accretion of pebbles in a gaseous environment », \emph{Monthly Notices of the Royal Astronomical
Society}, 404 (2010), 475--485. doi:10.1111/j.1365-2966.2010.16309.x.

\bibitem{Visser2020} R.~G. Visser, C.~W. Ormel, C. Dominik et S. Ida, « Spinning up planetary
bodies by pebble accretion », \emph{Icarus}, 335 (2020), 113380. doi:10.1016/j.icarus.2019.07.014.

\bibitem{Johansen2021} A. Johansen, T. Ronnet, M. Bizzarro, M. Schiller, M. Lambrechts,
Å. Nordlund et H. Lammer, « A pebble accretion model for the formation of the terrestrial planets
in the Solar System », \emph{Science Advances}, 7 (2021), eabc0444. doi:10.1126/sciadv.abc0444.

\bibitem{CanupWard2006} R.~M. Canup et W.~R. Ward, « A common mass scaling for satellite systems
of gaseous planets », \emph{Nature}, 441 (2006), 834--839. doi:10.1038/nature04860.

\bibitem{Toomre1964} A. Toomre, « On the gravitational stability of a disk of stars »,
\emph{The Astrophysical Journal}, 139 (1964), 1217--1238. doi:10.1086/147861.

\bibitem{Ida1997} S. Ida, R.~M. Canup et G.~R. Stewart, « Lunar accretion from an impact-generated
disk », \emph{Nature}, 389 (1997), 353--357. doi:10.1038/38669.

\bibitem{SalmonCanup2012} J. Salmon et R.~M. Canup, « Lunar accretion from a Roche-interior fluid
disk », \emph{The Astrophysical Journal}, 760 (2012), 83. doi:10.1088/0004-637X/760/1/83.

\bibitem{CanupWard2002} R.~M. Canup et W.~R. Ward, « Formation of the Galilean satellites:
conditions of accretion », \emph{The Astronomical Journal}, 124 (2002), 3404--3423.
doi:10.1086/344684.

\bibitem{Nimmo2026} F. Nimmo, R. Canup, Y. Fujii, A. Izidoro, J. Ji, W. McKinnon, O. Mousis,
C. Ormel, M. Ogihara, D. Stevenson et M. Blanc, « Origin and evolution of the Galilean satellites
within the Jovian system », \emph{Space Science Reviews}, 222 (2026), article~46.
doi:10.1007/s11214-026-01295-6.

\bibitem{Charnoz2010} S. Charnoz, J. Salmon et A. Crida, « The recent formation of Saturn's
moonlets from viscous spreading of the main rings », \emph{Nature}, 465 (2010), 752--754.
doi:10.1038/nature09096.

\bibitem{Blanc2025} M. Blanc, A. Crida, Y. Shibaike, S. Charnoz, M. El~Moutamid, P. Estrada,
O. Mousis, J. Salmon, A. Schneeberger et P. Vernazza, « Understanding the formation of Saturn's
regular moons in the context of giant planet moons formation scenarios », \emph{Space Science
Reviews}, 221 (2025), article~35. doi:10.1007/s11214-025-01156-8.

\bibitem{Helled2020} R. Helled, N. Nettelmann et T. Guillot, « Uranus and Neptune: origin,
evolution and internal structure », \emph{Space Science Reviews}, 216 (2020).
doi:10.1007/s11214-020-00660-3.

\bibitem{Benisty2021} M. Benisty, J. Bae, S. Facchini, M. Keppler, R. Teague, A. Isella,
N.~T. Kurtovic, L.~M. Pérez, A. Sierra, S.~M. Andrews et al., « A circumplanetary disk around
PDS70c », \emph{The Astrophysical Journal Letters}, 916 (2021), L2. doi:10.3847/2041-8213/ac0f83.

\bibitem{JPL} Jet Propulsion Laboratory, Solar System Dynamics, « Planetary satellite physical
parameters and mean elements », JPL/Caltech, données consultées en septembre 2026.

\bibitem{Popovas2018} A. Popovas, Å. Nordlund, J.~P. Ramsey et C.~W. Ormel, « Pebble dynamics and
accretion on to rocky planets -- I. Adiabatic and convective models », \emph{Monthly Notices of the
Royal Astronomical Society}, 479 (2018), 5136--5156. doi:10.1093/mnras/sty1752.

\bibitem{Ruskol1960} E.~L. Ruskol, « The origin of the Moon. I. Formation of a swarm of bodies
around the Earth », \emph{Soviet Astronomy}, 4 (1960), 657.

\bibitem{Weidenschilling1986} S.~J. Weidenschilling, R. Greenberg, C.~R. Chapman, F. Herbert,
D.~R. Davis, M.~J. Drake, J. Jones et W.~K. Hartmann, « Origin of the Moon from a circumterrestrial
disk », in W.~K. Hartmann, R.~J. Phillips et G.~J. Taylor (éds.), \emph{Origin of the Moon}, Lunar
and Planetary Institute, Houston (1986), 519--550.

\bibitem{QuillenTrilling1998} A.~C. Quillen et D.~E. Trilling, « Do proto-Jovian planets drive
outflows? », \emph{The Astrophysical Journal}, 508 (1998), 707--713.

\bibitem{MartinLubow2011} R.~G. Martin et S.~H. Lubow, « Tidal truncation of circumplanetary
discs », \emph{Monthly Notices of the Royal Astronomical Society}, 413 (2011), 1447--1461.

\bibitem{Ormel2015} C.~W. Ormel, J.-M. Shi et R. Kuiper, « Hydrodynamics of embedded planets' first
atmospheres -- II. A rapid recycling of atmospheric gas », \emph{Monthly Notices of the Royal
Astronomical Society}, 447 (2015), 3512--3525.

\bibitem{Fung2015} J. Fung, P. Artymowicz et Y. Wu, « The 3D flow field around an embedded
planet », \emph{The Astrophysical Journal}, 811 (2015), 101.

\bibitem{Daisaka2001} H. Daisaka, H. Tanaka et S. Ida, « Viscosity in a dense planetary ring with
self-gravitating particles », \emph{Icarus}, 154 (2001), 296--312.

\bibitem{Safronov1960} V.~S. Safronov, « On the gravitational instability in flattened systems with
axial symmetry and non-uniform rotation », \emph{Annales d'Astrophysique}, 23 (1960), 979.

\bibitem{GoldreichLyndenBell1965} P. Goldreich et D. Lynden-Bell, « Gravitational stability of
uniformly rotating disks », \emph{Monthly Notices of the Royal Astronomical Society}, 130 (1965),
97--124.

\bibitem{Kokubo2000} E. Kokubo, S. Ida et J. Makino, « Evolution of a circumterrestrial disk and
formation of a single moon », \emph{Icarus}, 148 (2000), 419--436.

\bibitem{PahlevanMorbidelli2015} K. Pahlevan et A. Morbidelli, « Collisionless encounters and the
origin of the lunar inclination », \emph{Nature}, 527 (2015), 492--494. doi:10.1038/nature16137.

\bibitem{Breslau2014} A. Breslau, M. Steinhausen, K. Vincke et S. Pfalzner, « Sizes of
protoplanetary discs after star-disc encounters », \emph{Astronomy \& Astrophysics}, 565 (2014),
A130.

\bibitem{Kruijer2015} T.~S. Kruijer, T. Kleine, M. Fischer-Gödde et P. Sprung, « Lunar tungsten
isotopic evidence for the late veneer », \emph{Nature}, 520 (2015), 534--537.
doi:10.1038/nature14360.

\bibitem{Touboul2015} M. Touboul, I.~S. Puchtel et R.~J. Walker, « Tungsten isotopic evidence for
disproportional late accretion to the Earth and Moon », \emph{Nature}, 520 (2015), 530--533.
doi:10.1038/nature14355.

\bibitem{Young2016} E.~D. Young, I.~E. Kohl, P.~H. Warren, D.~C. Rubie, S.~A. Jacobson et
A. Morbidelli, « Oxygen isotopic evidence for vigorous mixing during the Moon-forming giant
impact », \emph{Science}, 351 (2016), 493--496. doi:10.1126/science.aad0525.

\bibitem{Zhang2012} J. Zhang, N. Dauphas, A.~M. Davis, I. Leya et A. Fedkin, « The proto-Earth as
a significant source of lunar material », \emph{Nature Geoscience}, 5 (2012), 251--255.
doi:10.1038/ngeo1429.


\bibitem{Ringwood1990} A.~E. Ringwood, T. Kato, W. Hibberson et N. Ware, « High-pressure
geochemistry of Cr, V and Mn and implications for the origin of the Moon », \emph{Nature}, 347
(1990), 174--176. doi:10.1038/347174a0.

\bibitem{Drake1989} M.~J. Drake, H.~E. Newsom et C.~J. Capobianco, « V, Cr, and Mn in the Earth,
Moon, EPB, and SPB and the origin of the Moon: experimental studies », \emph{Geochimica et
Cosmochimica Acta}, 53 (1989), 2101--2111.

\bibitem{GessmannRubie2000} C.~K. Gessmann et D.~C. Rubie, « The origin of the depletions of V,
Cr and Mn in the mantles of the Earth and Moon », \emph{Earth and Planetary Science Letters}, 184
(2000), 95--107.

\bibitem{Newsom1986} H.~E. Newsom, « Constraints on the origin of the Moon from the abundance of
molybdenum and other siderophile elements », in W.~K. Hartmann, R.~J. Phillips et G.~J. Taylor
(éds.), \emph{Origin of the Moon}, Lunar and Planetary Institute, Houston (1986), 203--229.

\bibitem{Canup2023} R.~M. Canup, K. Righter, N. Dauphas et al., « Origin of the Moon »,
\emph{Reviews in Mineralogy and Geochemistry}, 89 (2023), 53--102. arXiv:2103.02045.

\bibitem{Spitzer1958} L. Spitzer, « Disruption of galactic clusters », \emph{The Astrophysical
Journal}, 127 (1958), 17--27.

\bibitem{BinneyTremaine2008} J. Binney et S. Tremaine, \emph{Galactic Dynamics}, 2\textsuperscript{e}
édition, Princeton University Press, Princeton (2008).

\bibitem{ToomreToomre1972} A. Toomre et J. Toomre, « Galactic bridges and tails », \emph{The
Astrophysical Journal}, 178 (1972), 623--666.


\bibitem{MosqueiraEstrada2003a} I. Mosqueira et P.~R. Estrada, « Formation of the regular
satellites of giant planets in an extended gaseous nebula I: Subnebula model and accretion of
satellites », \emph{Icarus}, 163 (2003), 198--231. doi:10.1016/S0019-1035(03)00076-9.

\bibitem{CridaCharnoz2012} A. Crida et S. Charnoz, « Formation of regular satellites from ancient
massive rings in the Solar System », \emph{Science}, 338 (2012), 1196--1199.
doi:10.1126/science.1226477.

\bibitem{AgnorHamilton2006} C.~B. Agnor et D.~P. Hamilton, « Neptune's capture of its moon Triton
in a binary--planet gravitational encounter », \emph{Nature}, 441 (2006), 192--194.
doi:10.1038/nature04792.

\bibitem{BanfieldMurray1992} D. Banfield et N. Murray, « A dynamical history of the inner
Neptunian satellites », \emph{Icarus}, 99 (1992), 390--401.
doi:10.1016/0019-1035(92)90155-Z.

\bibitem{Showalter2019} M.~R. Showalter, I. de Pater, J.~J. Lissauer et R.~S. French, « The
seventh inner moon of Neptune », \emph{Nature}, 566 (2019), 350--353.
doi:10.1038/s41586-019-0909-9.

\bibitem{Szulagyi2018} J. Szulágyi, M. Cilibrasi et L. Mayer, « In situ formation of icy moons of
Uranus and Neptune », \emph{The Astrophysical Journal Letters}, 868 (2018), L13.
doi:10.3847/2041-8213/aaeed6.

\bibitem{Davis2026} M.~R. Davis, M. Belyakov, I. Wong, Z. Milby et M.~E. Brown, « Neptune's Inner
Moons and Rings Are Exposed Icy Body Interiors », arXiv:2607.27481 (2026).


\bibitem{Tsiganis2005} K.~Tsiganis, R.~Gomes, A.~Morbidelli et H.~F. Levison, « Origin of the
orbital architecture of the giant planets of the Solar System », \emph{Nature}, 435 (2005),
459--461. doi:10.1038/nature03539.

\bibitem{Gomes2005} R.~Gomes, H.~F. Levison, K.~Tsiganis et A.~Morbidelli, « Origin of the
cataclysmic Late Heavy Bombardment period of the terrestrial planets », \emph{Nature}, 435
(2005), 466--469. doi:10.1038/nature03676.

\bibitem{Nesvorny2011} D.~Nesvorn\'y, « Young Solar System's fifth giant planet? »,
\emph{The Astrophysical Journal Letters}, 742 (2011), L22. doi:10.1088/2041-8205/742/2/L22.

\bibitem{NesvornyMorbidelli2012} D.~Nesvorn\'y et A.~Morbidelli, « Statistical study of the early
Solar System's instability with four, five, and six giant planets », \emph{The Astronomical
Journal}, 144 (2012), 117. doi:10.1088/0004-6256/144/4/117.

\bibitem{Liu2022} B.~Liu, S.~N. Raymond et S.~A. Jacobson, « Early Solar System instability
triggered by dispersal of the gaseous disk », \emph{Nature}, 604 (2022), 643--646.
doi:10.1038/s41586-022-04535-1.

\bibitem{Clement2018} M.~S. Clement, N.~A. Kaib, S.~N. Raymond et K.~J. Walsh, « Mars' growth
stunted by an early giant planet instability », \emph{Icarus}, 311 (2018), 340--356.
doi:10.1016/j.icarus.2018.04.008.

\end{thebibliography}

\end{document}